% Chapter written by unit TR-1 (plans/contraste/manuscripts-plan.md in CAOS_MANAGE).
% Every number below is a macro (\cs...) generated by tables/make_credit-scoring.py from data/derived; the tables and figures
% are generated by tables/make_credit-scoring.py and figures/make_credit-scoring.py.
\section{Methods: credit scoring and application models}
\label{sec:credit-scoring}
\makeatletter
\providecommand*{\input@path}{}
\g@addto@macro\input@path{{tables/}}
\makeatother
% Generated by tables/make_credit-scoring.py from data/derived; do not edit.

\newcommand{\csCthreeACnAuc}{0.5000}
\newcommand{\csCthreeACnAucHi}{0.5000}
\newcommand{\csCthreeACnAucLo}{0.5000}
\newcommand{\csCthreeACnBrier}{0.2195}
\newcommand{\csCthreeACnCal}{+0.1792}
\newcommand{\csCthreeACrisis}{no}
\newcommand{\csCthreeAEbAuc}{0.7340}
\newcommand{\csCthreeAEbAucHi}{0.7366}
\newcommand{\csCthreeAEbAucLo}{0.7314}
\newcommand{\csCthreeAEbBrier}{0.1754}
\newcommand{\csCthreeAEbCal}{+0.1102}
\newcommand{\csCthreeAEvents}{23{,}320}
\newcommand{\csCthreeALast}{2000}
\newcommand{\csCthreeALgAuc}{0.7967}
\newcommand{\csCthreeALgAucHi}{0.7990}
\newcommand{\csCthreeALgAucLo}{0.7944}
\newcommand{\csCthreeALgBrier}{0.1786}
\newcommand{\csCthreeALgCal}{+0.1324}
\newcommand{\csCthreeALlAuc}{0.8264}
\newcommand{\csCthreeALlAucHi}{0.8284}
\newcommand{\csCthreeALlAucLo}{0.8244}
\newcommand{\csCthreeALlBrier}{0.1885}
\newcommand{\csCthreeALlCal}{+0.1646}
\newcommand{\csCthreeALoans}{330{,}682}
\newcommand{\csCthreeAPlAuc}{0.7696}
\newcommand{\csCthreeAPlAucHi}{0.7722}
\newcommand{\csCthreeAPlAucLo}{0.7671}
\newcommand{\csCthreeAPlBrier}{0.1631}
\newcommand{\csCthreeAPlCal}{+0.0967}
\newcommand{\csCthreeAScAuc}{0.8252}
\newcommand{\csCthreeAScAucHi}{0.8272}
\newcommand{\csCthreeAScAucLo}{0.8231}
\newcommand{\csCthreeAScBrier}{0.1866}
\newcommand{\csCthreeAScCal}{+0.1619}
\newcommand{\csCthreeBCnAuc}{0.5000}
\newcommand{\csCthreeBCnAucHi}{0.5000}
\newcommand{\csCthreeBCnAucLo}{0.5000}
\newcommand{\csCthreeBCnBrier}{0.1179}
\newcommand{\csCthreeBCnCal}{+0.0625}
\newcommand{\csCthreeBCrisis}{no}
\newcommand{\csCthreeBEbAuc}{0.7769}
\newcommand{\csCthreeBEbAucHi}{0.7802}
\newcommand{\csCthreeBEbAucLo}{0.7736}
\newcommand{\csCthreeBEbBrier}{0.1062}
\newcommand{\csCthreeBEbCal}{+0.0497}
\newcommand{\csCthreeBEvents}{28{,}393}
\newcommand{\csCthreeBLast}{2002}
\newcommand{\csCthreeBLgAuc}{0.8323}
\newcommand{\csCthreeBLgAucHi}{0.8349}
\newcommand{\csCthreeBLgAucLo}{0.8297}
\newcommand{\csCthreeBLgBrier}{0.1034}
\newcommand{\csCthreeBLgCal}{+0.0572}
\newcommand{\csCthreeBLlAuc}{0.8137}
\newcommand{\csCthreeBLlAucHi}{0.8165}
\newcommand{\csCthreeBLlAucLo}{0.8109}
\newcommand{\csCthreeBLlBrier}{0.1070}
\newcommand{\csCthreeBLlCal}{+0.0671}
\newcommand{\csCthreeBLoans}{413{,}083}
\newcommand{\csCthreeBPlAuc}{0.7748}
\newcommand{\csCthreeBPlAucHi}{0.7786}
\newcommand{\csCthreeBPlAucLo}{0.7710}
\newcommand{\csCthreeBPlBrier}{0.1003}
\newcommand{\csCthreeBPlCal}{+0.0517}
\newcommand{\csCthreeBScAuc}{0.8149}
\newcommand{\csCthreeBScAucHi}{0.8178}
\newcommand{\csCthreeBScAucLo}{0.8121}
\newcommand{\csCthreeBScBrier}{0.1064}
\newcommand{\csCthreeBScCal}{+0.0665}
\newcommand{\csCthreeCCnAuc}{0.5000}
\newcommand{\csCthreeCCnAucHi}{0.5000}
\newcommand{\csCthreeCCnAucLo}{0.5000}
\newcommand{\csCthreeCCnBrier}{0.0349}
\newcommand{\csCthreeCCnCal}{$-$0.0467}
\newcommand{\csCthreeCCrisis}{yes}
\newcommand{\csCthreeCEbAuc}{0.8277}
\newcommand{\csCthreeCEbAucHi}{0.8331}
\newcommand{\csCthreeCEbAucLo}{0.8224}
\newcommand{\csCthreeCEbBrier}{0.0316}
\newcommand{\csCthreeCEbCal}{$-$0.0233}
\newcommand{\csCthreeCEvents}{50{,}614}
\newcommand{\csCthreeCLast}{2005}
\newcommand{\csCthreeCLgAuc}{0.8732}
\newcommand{\csCthreeCLgAucHi}{0.8769}
\newcommand{\csCthreeCLgAucLo}{0.8695}
\newcommand{\csCthreeCLgBrier}{0.0308}
\newcommand{\csCthreeCLgCal}{$-$0.0174}
\newcommand{\csCthreeCLlAuc}{0.8355}
\newcommand{\csCthreeCLlAucHi}{0.8406}
\newcommand{\csCthreeCLlAucLo}{0.8304}
\newcommand{\csCthreeCLlBrier}{0.0316}
\newcommand{\csCthreeCLlCal}{$-$0.0256}
\newcommand{\csCthreeCLoans}{627{,}594}
\newcommand{\csCthreeCPlAuc}{0.8139}
\newcommand{\csCthreeCPlAucHi}{0.8204}
\newcommand{\csCthreeCPlAucLo}{0.8075}
\newcommand{\csCthreeCPlBrier}{0.0308}
\newcommand{\csCthreeCPlCal}{$-$0.0261}
\newcommand{\csCthreeCScAuc}{0.8170}
\newcommand{\csCthreeCScAucHi}{0.8226}
\newcommand{\csCthreeCScAucLo}{0.8113}
\newcommand{\csCthreeCScBrier}{0.0321}
\newcommand{\csCthreeCScCal}{$-$0.0301}
\newcommand{\csCthreeCalRows}{20{,}000}
\newcommand{\csCthreeDCnAuc}{0.5000}
\newcommand{\csCthreeDCnAucHi}{0.5000}
\newcommand{\csCthreeDCnAucLo}{0.5000}
\newcommand{\csCthreeDCnBrier}{0.0453}
\newcommand{\csCthreeDCnCal}{$-$0.0825}
\newcommand{\csCthreeDCrisis}{yes}
\newcommand{\csCthreeDEbAuc}{0.7951}
\newcommand{\csCthreeDEbAucHi}{0.8004}
\newcommand{\csCthreeDEbAucLo}{0.7897}
\newcommand{\csCthreeDEbBrier}{0.0411}
\newcommand{\csCthreeDEbCal}{$-$0.0442}
\newcommand{\csCthreeDEvents}{110{,}349}
\newcommand{\csCthreeDLast}{2009}
\newcommand{\csCthreeDLgAuc}{0.8367}
\newcommand{\csCthreeDLgAucHi}{0.8412}
\newcommand{\csCthreeDLgAucLo}{0.8322}
\newcommand{\csCthreeDLgBrier}{0.0395}
\newcommand{\csCthreeDLgCal}{$-$0.0379}
\newcommand{\csCthreeDLlAuc}{0.8055}
\newcommand{\csCthreeDLlAucHi}{0.8106}
\newcommand{\csCthreeDLlAucLo}{0.8005}
\newcommand{\csCthreeDLlBrier}{0.0396}
\newcommand{\csCthreeDLlCal}{$-$0.0411}
\newcommand{\csCthreeDLoans}{900{,}254}
\newcommand{\csCthreeDPlAuc}{0.8489}
\newcommand{\csCthreeDPlAucHi}{0.8529}
\newcommand{\csCthreeDPlAucLo}{0.8448}
\newcommand{\csCthreeDPlBrier}{0.0396}
\newcommand{\csCthreeDPlCal}{$-$0.0411}
\newcommand{\csCthreeDScAuc}{0.7935}
\newcommand{\csCthreeDScAucHi}{0.7989}
\newcommand{\csCthreeDScAucLo}{0.7881}
\newcommand{\csCthreeDScBrier}{0.0403}
\newcommand{\csCthreeDScCal}{$-$0.0439}
\newcommand{\csCthreeECnAuc}{0.5000}
\newcommand{\csCthreeECnAucHi}{0.5000}
\newcommand{\csCthreeECnAucLo}{0.5000}
\newcommand{\csCthreeECnBrier}{0.0425}
\newcommand{\csCthreeECnCal}{$-$0.0739}
\newcommand{\csCthreeECrisis}{yes}
\newcommand{\csCthreeEEbAuc}{0.7460}
\newcommand{\csCthreeEEbAucHi}{0.7517}
\newcommand{\csCthreeEEbAucLo}{0.7404}
\newcommand{\csCthreeEEbBrier}{0.0395}
\newcommand{\csCthreeEEbCal}{$-$0.0252}
\newcommand{\csCthreeEEvents}{115{,}237}
\newcommand{\csCthreeELast}{2012}
\newcommand{\csCthreeELgAuc}{0.8305}
\newcommand{\csCthreeELgAucHi}{0.8344}
\newcommand{\csCthreeELgAucLo}{0.8265}
\newcommand{\csCthreeELgBrier}{0.0382}
\newcommand{\csCthreeELgCal}{$-$0.0159}
\newcommand{\csCthreeELlAuc}{0.8341}
\newcommand{\csCthreeELlAucHi}{0.8379}
\newcommand{\csCthreeELlAucLo}{0.8304}
\newcommand{\csCthreeELlBrier}{0.0383}
\newcommand{\csCthreeELlCal}{$-$0.0194}
\newcommand{\csCthreeELoans}{1{,}024{,}597}
\newcommand{\csCthreeEPlAuc}{0.7393}
\newcommand{\csCthreeEPlAucHi}{0.7462}
\newcommand{\csCthreeEPlAucLo}{0.7324}
\newcommand{\csCthreeEPlBrier}{0.0398}
\newcommand{\csCthreeEPlCal}{$-$0.0425}
\newcommand{\csCthreeEScAuc}{0.8264}
\newcommand{\csCthreeEScAucHi}{0.8303}
\newcommand{\csCthreeEScAucLo}{0.8225}
\newcommand{\csCthreeEScBrier}{0.0390}
\newcommand{\csCthreeEScCal}{$-$0.0221}
\newcommand{\csCthreeFitRows}{40{,}000}
\newcommand{\csCthreeLaterMax}{5.0}
\newcommand{\csCthreeLaterMin}{1.8}
\newcommand{\csCthreeVAAuc}{0.8332}
\newcommand{\csCthreeVAGap}{+13.51}
\newcommand{\csCthreeVAN}{86{,}393}
\newcommand{\csCthreeVAObs}{22.31}
\newcommand{\csCthreeVAPred}{8.80}
\newcommand{\csCthreeVBAuc}{0.8204}
\newcommand{\csCthreeVBGap}{+18.82}
\newcommand{\csCthreeVBN}{88{,}134}
\newcommand{\csCthreeVBObs}{27.58}
\newcommand{\csCthreeVBPred}{8.76}
\newcommand{\csCthreeVCAuc}{0.8132}
\newcommand{\csCthreeVCGap}{+13.91}
\newcommand{\csCthreeVCN}{61{,}539}
\newcommand{\csCthreeVCObs}{21.27}
\newcommand{\csCthreeVCPred}{7.35}
\newcommand{\csCthreeVDAuc}{0.7845}
\newcommand{\csCthreeVDGap}{+2.23}
\newcommand{\csCthreeVDN}{36{,}594}
\newcommand{\csCthreeVDObs}{8.38}
\newcommand{\csCthreeVDPred}{6.15}
\newcommand{\csCthreeVEAuc}{0.8102}
\newcommand{\csCthreeVEGap}{$-$0.51}
\newcommand{\csCthreeVEN}{39{,}890}
\newcommand{\csCthreeVEObs}{4.90}
\newcommand{\csCthreeVEPred}{5.41}
\newcommand{\csCthreeVFAuc}{0.7983}
\newcommand{\csCthreeVFGap}{$-$2.84}
\newcommand{\csCthreeVFN}{45{,}592}
\newcommand{\csCthreeVFObs}{3.65}
\newcommand{\csCthreeVFPred}{6.49}
\newcommand{\csCthreeVGAuc}{0.8261}
\newcommand{\csCthreeVGGap}{$-$3.27}
\newcommand{\csCthreeVGN}{38{,}861}
\newcommand{\csCthreeVGObs}{3.27}
\newcommand{\csCthreeVGPred}{6.55}
\newcommand{\csCthreeVHAuc}{0.8236}
\newcommand{\csCthreeVHGap}{$-$3.13}
\newcommand{\csCthreeVHN}{40{,}400}
\newcommand{\csCthreeVHObs}{3.07}
\newcommand{\csCthreeVHPred}{6.20}
\newcommand{\csCthreeVIAuc}{0.8235}
\newcommand{\csCthreeVIGap}{$-$2.84}
\newcommand{\csCthreeVIN}{45{,}942}
\newcommand{\csCthreeVIObs}{3.52}
\newcommand{\csCthreeVIPred}{6.36}
\newcommand{\csCthreeVJAuc}{0.8085}
\newcommand{\csCthreeVJGap}{$-$4.97}
\newcommand{\csCthreeVJN}{55{,}387}
\newcommand{\csCthreeVJObs}{3.85}
\newcommand{\csCthreeVJPred}{8.82}
\newcommand{\csCthreeVKAuc}{0.7932}
\newcommand{\csCthreeVKGap}{$-$4.49}
\newcommand{\csCthreeVKN}{56{,}759}
\newcommand{\csCthreeVKObs}{3.87}
\newcommand{\csCthreeVKPred}{8.36}
\newcommand{\csCthreeVLAuc}{0.7841}
\newcommand{\csCthreeVLGap}{$-$3.73}
\newcommand{\csCthreeVLN}{56{,}057}
\newcommand{\csCthreeVLObs}{4.31}
\newcommand{\csCthreeVLPred}{8.04}
\newcommand{\csCthreeVMAuc}{0.8104}
\newcommand{\csCthreeVMGap}{$-$2.01}
\newcommand{\csCthreeVMN}{54{,}171}
\newcommand{\csCthreeVMObs}{4.66}
\newcommand{\csCthreeVMPred}{6.68}
\newcommand{\csCthreeVNAuc}{0.8321}
\newcommand{\csCthreeVNGap}{$-$1.85}
\newcommand{\csCthreeVNN}{45{,}664}
\newcommand{\csCthreeVNObs}{3.91}
\newcommand{\csCthreeVNPred}{5.76}
\newcommand{\csCthreeVOAuc}{0.8418}
\newcommand{\csCthreeVOGap}{$-$2.96}
\newcommand{\csCthreeVON}{36{,}456}
\newcommand{\csCthreeVOObs}{2.59}
\newcommand{\csCthreeVOPred}{5.55}
\newcommand{\csCtwoADisagree}{39.7}
\newcommand{\csCtwoAEbAuc}{0.9538}
\newcommand{\csCtwoAEbAucHi}{0.9720}
\newcommand{\csCtwoAEbAucLo}{0.9355}
\newcommand{\csCtwoAEbBrier}{0.0289}
\newcommand{\csCtwoAEbBrierDiff}{$-$0.0168}
\newcommand{\csCtwoAEbDiff}{+0.0431}
\newcommand{\csCtwoAEbDiffHi}{+0.0645}
\newcommand{\csCtwoAEbDiffLo}{+0.0218}
\newcommand{\csCtwoAEbDlP}{\ensuremath{7.4\times10^{-5}}}
\newcommand{\csCtwoAEbGap}{0.0385}
\newcommand{\csCtwoAFlags}{1}
\newcommand{\csCtwoALgAuc}{0.9520}
\newcommand{\csCtwoALgAucHi}{0.9713}
\newcommand{\csCtwoALgAucLo}{0.9327}
\newcommand{\csCtwoALgBrier}{0.0306}
\newcommand{\csCtwoALgBrierDiff}{$-$0.0151}
\newcommand{\csCtwoALgDiff}{+0.0414}
\newcommand{\csCtwoALgDiffHi}{+0.0584}
\newcommand{\csCtwoALgDiffLo}{+0.0244}
\newcommand{\csCtwoALgDlP}{\ensuremath{1.8\times10^{-6}}}
\newcommand{\csCtwoALgGap}{0.0479}
\newcommand{\csCtwoALlAuc}{0.9115}
\newcommand{\csCtwoALlAucHi}{0.9359}
\newcommand{\csCtwoALlAucLo}{0.8871}
\newcommand{\csCtwoALlBrier}{0.0459}
\newcommand{\csCtwoALlBrierDiff}{+0.0002}
\newcommand{\csCtwoALlDiff}{+0.0009}
\newcommand{\csCtwoALlDiffHi}{+0.0147}
\newcommand{\csCtwoALlDiffLo}{$-$0.0128}
\newcommand{\csCtwoALlDlP}{0.897}
\newcommand{\csCtwoALlGap}{0.0226}
\newcommand{\csCtwoAPlAuc}{0.9343}
\newcommand{\csCtwoAPlAucHi}{0.9530}
\newcommand{\csCtwoAPlAucLo}{0.9156}
\newcommand{\csCtwoAPlBrier}{0.0412}
\newcommand{\csCtwoAPlBrierDiff}{$-$0.0044}
\newcommand{\csCtwoAPlDiff}{+0.0237}
\newcommand{\csCtwoAPlDiffHi}{+0.0390}
\newcommand{\csCtwoAPlDiffLo}{+0.0083}
\newcommand{\csCtwoAPlDlP}{0.003}
\newcommand{\csCtwoAPlGap}{0.0394}
\newcommand{\csCtwoAPzAuc}{0.8128}
\newcommand{\csCtwoAPzAucHi}{0.8562}
\newcommand{\csCtwoAPzAucLo}{0.7695}
\newcommand{\csCtwoAPzBrier}{0.0527}
\newcommand{\csCtwoAPzGap}{-0.0269}
\newcommand{\csCtwoAScAuc}{0.9106}
\newcommand{\csCtwoAScAucHi}{0.9343}
\newcommand{\csCtwoAScAucLo}{0.8869}
\newcommand{\csCtwoAScBrier}{0.0457}
\newcommand{\csCtwoAScGap}{0.0052}
\newcommand{\csCtwoAScHalf}{0.0237}
\newcommand{\csCtwoAXgAuc}{0.9412}
\newcommand{\csCtwoAXgAucHi}{0.9633}
\newcommand{\csCtwoAXgAucLo}{0.9191}
\newcommand{\csCtwoAXgBrier}{0.0316}
\newcommand{\csCtwoAXgBrierDiff}{$-$0.0140}
\newcommand{\csCtwoAXgDiff}{+0.0306}
\newcommand{\csCtwoAXgDiffHi}{+0.0490}
\newcommand{\csCtwoAXgDiffLo}{+0.0122}
\newcommand{\csCtwoAXgDlP}{0.001}
\newcommand{\csCtwoAXgGap}{0.0551}
\newcommand{\csCtwoBDisagree}{49.5}
\newcommand{\csCtwoBEbAuc}{0.9217}
\newcommand{\csCtwoBEbAucHi}{0.9457}
\newcommand{\csCtwoBEbAucLo}{0.8977}
\newcommand{\csCtwoBEbBrier}{0.0270}
\newcommand{\csCtwoBEbBrierDiff}{$-$0.0171}
\newcommand{\csCtwoBEbDiff}{+0.1117}
\newcommand{\csCtwoBEbDiffHi}{+0.1401}
\newcommand{\csCtwoBEbDiffLo}{+0.0833}
\newcommand{\csCtwoBEbDlP}{\ensuremath{1.2\times10^{-14}}}
\newcommand{\csCtwoBEbGap}{0.0531}
\newcommand{\csCtwoBFlags}{4}
\newcommand{\csCtwoBLgAuc}{0.9228}
\newcommand{\csCtwoBLgAucHi}{0.9464}
\newcommand{\csCtwoBLgAucLo}{0.8992}
\newcommand{\csCtwoBLgBrier}{0.0291}
\newcommand{\csCtwoBLgBrierDiff}{$-$0.0151}
\newcommand{\csCtwoBLgDiff}{+0.1128}
\newcommand{\csCtwoBLgDiffHi}{+0.1401}
\newcommand{\csCtwoBLgDiffLo}{+0.0854}
\newcommand{\csCtwoBLgDlP}{\ensuremath{6.3\times10^{-16}}}
\newcommand{\csCtwoBLgGap}{0.0769}
\newcommand{\csCtwoBLlAuc}{0.8348}
\newcommand{\csCtwoBLlAucHi}{0.8687}
\newcommand{\csCtwoBLlAucLo}{0.8010}
\newcommand{\csCtwoBLlBrier}{0.0434}
\newcommand{\csCtwoBLlBrierDiff}{$-$0.0008}
\newcommand{\csCtwoBLlDiff}{+0.0248}
\newcommand{\csCtwoBLlDiffHi}{+0.0376}
\newcommand{\csCtwoBLlDiffLo}{+0.0120}
\newcommand{\csCtwoBLlDlP}{\ensuremath{1.5\times10^{-4}}}
\newcommand{\csCtwoBLlGap}{0.0293}
\newcommand{\csCtwoBPlAuc}{0.9036}
\newcommand{\csCtwoBPlAucHi}{0.9273}
\newcommand{\csCtwoBPlAucLo}{0.8799}
\newcommand{\csCtwoBPlBrier}{0.0335}
\newcommand{\csCtwoBPlBrierDiff}{$-$0.0107}
\newcommand{\csCtwoBPlDiff}{+0.0936}
\newcommand{\csCtwoBPlDiffHi}{+0.1217}
\newcommand{\csCtwoBPlDiffLo}{+0.0655}
\newcommand{\csCtwoBPlDlP}{\ensuremath{6.5\times10^{-11}}}
\newcommand{\csCtwoBPlGap}{0.0705}
\newcommand{\csCtwoBPzAuc}{0.6698}
\newcommand{\csCtwoBPzAucHi}{0.7126}
\newcommand{\csCtwoBPzAucLo}{0.6270}
\newcommand{\csCtwoBPzBrier}{0.0489}
\newcommand{\csCtwoBPzGap}{0.0492}
\newcommand{\csCtwoBScAuc}{0.8100}
\newcommand{\csCtwoBScAucHi}{0.8478}
\newcommand{\csCtwoBScAucLo}{0.7723}
\newcommand{\csCtwoBScBrier}{0.0441}
\newcommand{\csCtwoBScGap}{0.0333}
\newcommand{\csCtwoBScHalf}{0.0377}
\newcommand{\csCtwoBXgAuc}{0.9124}
\newcommand{\csCtwoBXgAucHi}{0.9368}
\newcommand{\csCtwoBXgAucLo}{0.8880}
\newcommand{\csCtwoBXgBrier}{0.0313}
\newcommand{\csCtwoBXgBrierDiff}{$-$0.0128}
\newcommand{\csCtwoBXgDiff}{+0.1023}
\newcommand{\csCtwoBXgDiffHi}{+0.1288}
\newcommand{\csCtwoBXgDiffLo}{+0.0758}
\newcommand{\csCtwoBXgDlP}{\ensuremath{3.8\times10^{-14}}}
\newcommand{\csCtwoBXgGap}{0.0735}
\newcommand{\csCtwoCDisagree}{43.4}
\newcommand{\csCtwoCEbAuc}{0.8983}
\newcommand{\csCtwoCEbAucHi}{0.9242}
\newcommand{\csCtwoCEbAucLo}{0.8724}
\newcommand{\csCtwoCEbBrier}{0.0301}
\newcommand{\csCtwoCEbBrierDiff}{$-$0.0084}
\newcommand{\csCtwoCEbDiff}{+0.0595}
\newcommand{\csCtwoCEbDiffHi}{+0.0900}
\newcommand{\csCtwoCEbDiffLo}{+0.0290}
\newcommand{\csCtwoCEbDlP}{\ensuremath{1.3\times10^{-4}}}
\newcommand{\csCtwoCEbGap}{0.0812}
\newcommand{\csCtwoCFlags}{4}
\newcommand{\csCtwoCLgAuc}{0.9253}
\newcommand{\csCtwoCLgAucHi}{0.9472}
\newcommand{\csCtwoCLgAucLo}{0.9034}
\newcommand{\csCtwoCLgBrier}{0.0274}
\newcommand{\csCtwoCLgBrierDiff}{$-$0.0111}
\newcommand{\csCtwoCLgDiff}{+0.0865}
\newcommand{\csCtwoCLgDiffHi}{+0.1145}
\newcommand{\csCtwoCLgDiffLo}{+0.0585}
\newcommand{\csCtwoCLgDlP}{\ensuremath{1.4\times10^{-9}}}
\newcommand{\csCtwoCLgGap}{0.0738}
\newcommand{\csCtwoCLlAuc}{0.8670}
\newcommand{\csCtwoCLlAucHi}{0.8984}
\newcommand{\csCtwoCLlAucLo}{0.8356}
\newcommand{\csCtwoCLlBrier}{0.0371}
\newcommand{\csCtwoCLlBrierDiff}{$-$0.0014}
\newcommand{\csCtwoCLlDiff}{+0.0282}
\newcommand{\csCtwoCLlDiffHi}{+0.0456}
\newcommand{\csCtwoCLlDiffLo}{+0.0107}
\newcommand{\csCtwoCLlDlP}{0.002}
\newcommand{\csCtwoCLlGap}{0.0147}
\newcommand{\csCtwoCPlAuc}{0.8895}
\newcommand{\csCtwoCPlAucHi}{0.9170}
\newcommand{\csCtwoCPlAucLo}{0.8621}
\newcommand{\csCtwoCPlBrier}{0.0328}
\newcommand{\csCtwoCPlBrierDiff}{$-$0.0057}
\newcommand{\csCtwoCPlDiff}{+0.0507}
\newcommand{\csCtwoCPlDiffHi}{+0.0783}
\newcommand{\csCtwoCPlDiffLo}{+0.0232}
\newcommand{\csCtwoCPlDlP}{\ensuremath{3.1\times10^{-4}}}
\newcommand{\csCtwoCPlGap}{0.0923}
\newcommand{\csCtwoCPzAuc}{0.6120}
\newcommand{\csCtwoCPzAucHi}{0.6484}
\newcommand{\csCtwoCPzAucLo}{0.5755}
\newcommand{\csCtwoCPzBrier}{0.0438}
\newcommand{\csCtwoCPzGap}{0.0496}
\newcommand{\csCtwoCScAuc}{0.8388}
\newcommand{\csCtwoCScAucHi}{0.8756}
\newcommand{\csCtwoCScAucLo}{0.8021}
\newcommand{\csCtwoCScBrier}{0.0385}
\newcommand{\csCtwoCScGap}{0.0179}
\newcommand{\csCtwoCScHalf}{0.0368}
\newcommand{\csCtwoCXgAuc}{0.9009}
\newcommand{\csCtwoCXgAucHi}{0.9292}
\newcommand{\csCtwoCXgAucLo}{0.8727}
\newcommand{\csCtwoCXgBrier}{0.0300}
\newcommand{\csCtwoCXgBrierDiff}{$-$0.0085}
\newcommand{\csCtwoCXgDiff}{+0.0621}
\newcommand{\csCtwoCXgDiffHi}{+0.0882}
\newcommand{\csCtwoCXgDiffLo}{+0.0360}
\newcommand{\csCtwoCXgDlP}{\ensuremath{3.0\times10^{-6}}}
\newcommand{\csCtwoCXgGap}{0.0950}
\newcommand{\csCtwoDDisagree}{41.1}
\newcommand{\csCtwoDEbAuc}{0.9011}
\newcommand{\csCtwoDEbAucHi}{0.9291}
\newcommand{\csCtwoDEbAucLo}{0.8731}
\newcommand{\csCtwoDEbBrier}{0.0244}
\newcommand{\csCtwoDEbBrierDiff}{$-$0.0077}
\newcommand{\csCtwoDEbDiff}{+0.1091}
\newcommand{\csCtwoDEbDiffHi}{+0.1386}
\newcommand{\csCtwoDEbDiffLo}{+0.0797}
\newcommand{\csCtwoDEbDlP}{\ensuremath{3.7\times10^{-13}}}
\newcommand{\csCtwoDEbGap}{0.0748}
\newcommand{\csCtwoDFlags}{6}
\newcommand{\csCtwoDLgAuc}{0.9117}
\newcommand{\csCtwoDLgAucHi}{0.9405}
\newcommand{\csCtwoDLgAucLo}{0.8829}
\newcommand{\csCtwoDLgBrier}{0.0235}
\newcommand{\csCtwoDLgBrierDiff}{$-$0.0086}
\newcommand{\csCtwoDLgDiff}{+0.1197}
\newcommand{\csCtwoDLgDiffHi}{+0.1502}
\newcommand{\csCtwoDLgDiffLo}{+0.0891}
\newcommand{\csCtwoDLgDlP}{\ensuremath{1.6\times10^{-14}}}
\newcommand{\csCtwoDLgGap}{0.0878}
\newcommand{\csCtwoDLlAuc}{0.8123}
\newcommand{\csCtwoDLlAucHi}{0.8567}
\newcommand{\csCtwoDLlAucLo}{0.7679}
\newcommand{\csCtwoDLlBrier}{0.0304}
\newcommand{\csCtwoDLlBrierDiff}{$-$0.0017}
\newcommand{\csCtwoDLlDiff}{+0.0202}
\newcommand{\csCtwoDLlDiffHi}{+0.0403}
\newcommand{\csCtwoDLlDiffLo}{+0.0002}
\newcommand{\csCtwoDLlDlP}{0.048}
\newcommand{\csCtwoDLlGap}{0.0674}
\newcommand{\csCtwoDPlAuc}{0.8973}
\newcommand{\csCtwoDPlAucHi}{0.9283}
\newcommand{\csCtwoDPlAucLo}{0.8664}
\newcommand{\csCtwoDPlBrier}{0.0244}
\newcommand{\csCtwoDPlBrierDiff}{$-$0.0077}
\newcommand{\csCtwoDPlDiff}{+0.1053}
\newcommand{\csCtwoDPlDiffHi}{+0.1355}
\newcommand{\csCtwoDPlDiffLo}{+0.0751}
\newcommand{\csCtwoDPlDlP}{\ensuremath{8.3\times10^{-12}}}
\newcommand{\csCtwoDPlGap}{0.0849}
\newcommand{\csCtwoDPzAuc}{0.6737}
\newcommand{\csCtwoDPzAucHi}{0.7179}
\newcommand{\csCtwoDPzAucLo}{0.6295}
\newcommand{\csCtwoDPzBrier}{0.0350}
\newcommand{\csCtwoDPzGap}{-0.0156}
\newcommand{\csCtwoDScAuc}{0.7920}
\newcommand{\csCtwoDScAucHi}{0.8355}
\newcommand{\csCtwoDScAucLo}{0.7486}
\newcommand{\csCtwoDScBrier}{0.0321}
\newcommand{\csCtwoDScGap}{0.0655}
\newcommand{\csCtwoDScHalf}{0.0434}
\newcommand{\csCtwoDXgAuc}{0.8887}
\newcommand{\csCtwoDXgAucHi}{0.9220}
\newcommand{\csCtwoDXgAucLo}{0.8555}
\newcommand{\csCtwoDXgBrier}{0.0237}
\newcommand{\csCtwoDXgBrierDiff}{$-$0.0084}
\newcommand{\csCtwoDXgDiff}{+0.0967}
\newcommand{\csCtwoDXgDiffHi}{+0.1292}
\newcommand{\csCtwoDXgDiffLo}{+0.0642}
\newcommand{\csCtwoDXgDlP}{\ensuremath{5.6\times10^{-9}}}
\newcommand{\csCtwoDXgGap}{0.1016}
\newcommand{\csCtwoEDisagree}{50.0}
\newcommand{\csCtwoEEbAuc}{0.9320}
\newcommand{\csCtwoEEbAucHi}{0.9615}
\newcommand{\csCtwoEEbAucLo}{0.9024}
\newcommand{\csCtwoEEbBrier}{0.0184}
\newcommand{\csCtwoEEbBrierDiff}{$-$0.0185}
\newcommand{\csCtwoEEbDiff}{+0.2041}
\newcommand{\csCtwoEEbDiffHi}{+0.2589}
\newcommand{\csCtwoEEbDiffLo}{+0.1493}
\newcommand{\csCtwoEEbDlP}{\ensuremath{2.8\times10^{-13}}}
\newcommand{\csCtwoEEbGap}{0.0620}
\newcommand{\csCtwoEFlags}{6}
\newcommand{\csCtwoELgAuc}{0.9428}
\newcommand{\csCtwoELgAucHi}{0.9677}
\newcommand{\csCtwoELgAucLo}{0.9179}
\newcommand{\csCtwoELgBrier}{0.0178}
\newcommand{\csCtwoELgBrierDiff}{$-$0.0191}
\newcommand{\csCtwoELgDiff}{+0.2149}
\newcommand{\csCtwoELgDiffHi}{+0.2680}
\newcommand{\csCtwoELgDiffLo}{+0.1618}
\newcommand{\csCtwoELgDlP}{\ensuremath{2.2\times10^{-15}}}
\newcommand{\csCtwoELgGap}{0.0571}
\newcommand{\csCtwoELlAuc}{0.7583}
\newcommand{\csCtwoELlAucHi}{0.8081}
\newcommand{\csCtwoELlAucLo}{0.7084}
\newcommand{\csCtwoELlBrier}{0.0359}
\newcommand{\csCtwoELlBrierDiff}{$-$0.0011}
\newcommand{\csCtwoELlDiff}{+0.0304}
\newcommand{\csCtwoELlDiffHi}{+0.0482}
\newcommand{\csCtwoELlDiffLo}{+0.0125}
\newcommand{\csCtwoELlDlP}{\ensuremath{8.5\times10^{-4}}}
\newcommand{\csCtwoELlGap}{0.0773}
\newcommand{\csCtwoEPlAuc}{0.9029}
\newcommand{\csCtwoEPlAucHi}{0.9414}
\newcommand{\csCtwoEPlAucLo}{0.8644}
\newcommand{\csCtwoEPlBrier}{0.0237}
\newcommand{\csCtwoEPlBrierDiff}{$-$0.0132}
\newcommand{\csCtwoEPlDiff}{+0.1750}
\newcommand{\csCtwoEPlDiffHi}{+0.2349}
\newcommand{\csCtwoEPlDiffLo}{+0.1151}
\newcommand{\csCtwoEPlDlP}{\ensuremath{1.0\times10^{-8}}}
\newcommand{\csCtwoEPlGap}{0.0789}
\newcommand{\csCtwoEPzAuc}{0.7433}
\newcommand{\csCtwoEPzAucHi}{0.7891}
\newcommand{\csCtwoEPzAucLo}{0.6974}
\newcommand{\csCtwoEPzBrier}{0.0359}
\newcommand{\csCtwoEPzGap}{-0.0050}
\newcommand{\csCtwoEScAuc}{0.7279}
\newcommand{\csCtwoEScAucHi}{0.7827}
\newcommand{\csCtwoEScAucLo}{0.6731}
\newcommand{\csCtwoEScBrier}{0.0369}
\newcommand{\csCtwoEScGap}{0.0922}
\newcommand{\csCtwoEScHalf}{0.0548}
\newcommand{\csCtwoEXgAuc}{0.9269}
\newcommand{\csCtwoEXgAucHi}{0.9608}
\newcommand{\csCtwoEXgAucLo}{0.8930}
\newcommand{\csCtwoEXgBrier}{0.0198}
\newcommand{\csCtwoEXgBrierDiff}{$-$0.0171}
\newcommand{\csCtwoEXgDiff}{+0.1990}
\newcommand{\csCtwoEXgDiffHi}{+0.2536}
\newcommand{\csCtwoEXgDiffLo}{+0.1445}
\newcommand{\csCtwoEXgDlP}{\ensuremath{8.8\times10^{-13}}}
\newcommand{\csCtwoEXgGap}{0.0626}
\newcommand{\csCtwoEventsA}{410}
\newcommand{\csCtwoEventsB}{515}
\newcommand{\csCtwoEventsC}{495}
\newcommand{\csCtwoEventsD}{400}
\newcommand{\csCtwoEventsE}{271}
\newcommand{\csCtwoEventsP}{2{,}091}
\newcommand{\csCtwoFlagTotal}{21}
\newcommand{\csCtwoHoldA}{1{,}773}
\newcommand{\csCtwoHoldB}{2{,}938}
\newcommand{\csCtwoHoldC}{3{,}151}
\newcommand{\csCtwoHoldD}{3{,}052}
\newcommand{\csCtwoHoldE}{2{,}109}
\newcommand{\csCtwoHoldEvA}{123}
\newcommand{\csCtwoHoldEvB}{155}
\newcommand{\csCtwoHoldEvC}{149}
\newcommand{\csCtwoHoldEvD}{120}
\newcommand{\csCtwoHoldEvE}{81}
\newcommand{\csCtwoHoldEvP}{628}
\newcommand{\csCtwoHoldP}{13{,}023}
\newcommand{\csCtwoPDisagree}{47.4}
\newcommand{\csCtwoPEbAuc}{0.9450}
\newcommand{\csCtwoPEbAucHi}{0.9538}
\newcommand{\csCtwoPEbAucLo}{0.9362}
\newcommand{\csCtwoPEbBrier}{0.0232}
\newcommand{\csCtwoPEbBrierDiff}{$-$0.0141}
\newcommand{\csCtwoPEbDiff}{+0.0802}
\newcommand{\csCtwoPEbDiffHi}{+0.0918}
\newcommand{\csCtwoPEbDiffLo}{+0.0686}
\newcommand{\csCtwoPEbDlP}{\ensuremath{4.8\times10^{-42}}}
\newcommand{\csCtwoPEbGap}{0.0243}
\newcommand{\csCtwoPFlags}{0}
\newcommand{\csCtwoPLgAuc}{0.9703}
\newcommand{\csCtwoPLgAucHi}{0.9768}
\newcommand{\csCtwoPLgAucLo}{0.9638}
\newcommand{\csCtwoPLgBrier}{0.0176}
\newcommand{\csCtwoPLgBrierDiff}{$-$0.0197}
\newcommand{\csCtwoPLgDiff}{+0.1055}
\newcommand{\csCtwoPLgDiffHi}{+0.1180}
\newcommand{\csCtwoPLgDiffLo}{+0.0931}
\newcommand{\csCtwoPLgDlP}{\ensuremath{3.4\times10^{-62}}}
\newcommand{\csCtwoPLgGap}{0.0295}
\newcommand{\csCtwoPLlAuc}{0.8824}
\newcommand{\csCtwoPLlAucHi}{0.8963}
\newcommand{\csCtwoPLlAucLo}{0.8684}
\newcommand{\csCtwoPLlBrier}{0.0349}
\newcommand{\csCtwoPLlBrierDiff}{$-$0.0023}
\newcommand{\csCtwoPLlDiff}{+0.0176}
\newcommand{\csCtwoPLlDiffHi}{+0.0246}
\newcommand{\csCtwoPLlDiffLo}{+0.0106}
\newcommand{\csCtwoPLlDlP}{\ensuremath{8.8\times10^{-7}}}
\newcommand{\csCtwoPLlGap}{0.0039}
\newcommand{\csCtwoPPlAuc}{0.9168}
\newcommand{\csCtwoPPlAucHi}{0.9286}
\newcommand{\csCtwoPPlAucLo}{0.9049}
\newcommand{\csCtwoPPlBrier}{0.0262}
\newcommand{\csCtwoPPlBrierDiff}{$-$0.0111}
\newcommand{\csCtwoPPlDiff}{+0.0520}
\newcommand{\csCtwoPPlDiffHi}{+0.0640}
\newcommand{\csCtwoPPlDiffLo}{+0.0399}
\newcommand{\csCtwoPPlDlP}{\ensuremath{2.4\times10^{-17}}}
\newcommand{\csCtwoPPlGap}{0.0427}
\newcommand{\csCtwoPPzAuc}{0.6519}
\newcommand{\csCtwoPPzAucHi}{0.6708}
\newcommand{\csCtwoPPzAucLo}{0.6331}
\newcommand{\csCtwoPPzBrier}{0.0427}
\newcommand{\csCtwoPPzGap}{-0.0084}
\newcommand{\csCtwoPScAuc}{0.8648}
\newcommand{\csCtwoPScAucHi}{0.8796}
\newcommand{\csCtwoPScAucLo}{0.8500}
\newcommand{\csCtwoPScBrier}{0.0372}
\newcommand{\csCtwoPScGap}{-0.0000}
\newcommand{\csCtwoPScHalf}{0.0148}
\newcommand{\csCtwoPXgAuc}{0.9642}
\newcommand{\csCtwoPXgAucHi}{0.9715}
\newcommand{\csCtwoPXgAucLo}{0.9568}
\newcommand{\csCtwoPXgBrier}{0.0186}
\newcommand{\csCtwoPXgBrierDiff}{$-$0.0186}
\newcommand{\csCtwoPXgDiff}{+0.0993}
\newcommand{\csCtwoPXgDiffHi}{+0.1115}
\newcommand{\csCtwoPXgDiffLo}{+0.0872}
\newcommand{\csCtwoPXgDlP}{\ensuremath{2.1\times10^{-57}}}
\newcommand{\csCtwoPXgGap}{0.0352}
\newcommand{\csCtwoRateA}{6.94}
\newcommand{\csCtwoRateB}{5.26}
\newcommand{\csCtwoRateC}{4.71}
\newcommand{\csCtwoRateD}{3.93}
\newcommand{\csCtwoRateE}{3.86}
\newcommand{\csCtwoRateP}{4.82}
\newcommand{\csCtwoRowsA}{5{,}910}
\newcommand{\csCtwoRowsB}{9{,}792}
\newcommand{\csCtwoRowsC}{10{,}503}
\newcommand{\csCtwoRowsD}{10{,}173}
\newcommand{\csCtwoRowsE}{7{,}027}
\newcommand{\csCtwoRowsP}{43{,}405}
\newcommand{\csCtwoRungTotal}{35}
\newcommand{\csDmCnAuc}{0.5000}
\newcommand{\csDmCnAucHi}{0.5000}
\newcommand{\csDmCnAucLo}{0.5000}
\newcommand{\csDmCnBrier}{0.1924}
\newcommand{\csDmCnEce}{0.0361}
\newcommand{\csDmCnGini}{0.0000}
\newcommand{\csDmCnSkill}{-0.007}
\newcommand{\csDmCnSpP}{\ensuremath{7.2\times10^{-17}}}
\newcommand{\csDmEbAuc}{0.7953}
\newcommand{\csDmEbAucHi}{0.8062}
\newcommand{\csDmEbAucLo}{0.7844}
\newcommand{\csDmEbBrier}{0.1433}
\newcommand{\csDmEbDlDiff}{+0.0110}
\newcommand{\csDmEbDlP}{\ensuremath{4.2\times10^{-8}}}
\newcommand{\csDmEbEce}{0.0149}
\newcommand{\csDmEbGini}{0.5906}
\newcommand{\csDmEbHlP}{0.004}
\newcommand{\csDmEbKs}{0.4730}
\newcommand{\csDmEbPsi}{0.0325}
\newcommand{\csDmEbPsiBench}{0.0029}
\newcommand{\csDmEbSkill}{0.250}
\newcommand{\csDmEbSpP}{0.370}
\newcommand{\csDmEvents}{2{,}316}
\newcommand{\csDmLgAuc}{0.7994}
\newcommand{\csDmLgAucHi}{0.8100}
\newcommand{\csDmLgAucLo}{0.7887}
\newcommand{\csDmLgBrier}{0.1429}
\newcommand{\csDmLgDlDiff}{+0.0151}
\newcommand{\csDmLgDlP}{\ensuremath{4.6\times10^{-13}}}
\newcommand{\csDmLgEce}{0.0149}
\newcommand{\csDmLgGini}{0.5988}
\newcommand{\csDmLgHlP}{0.067}
\newcommand{\csDmLgKs}{0.4658}
\newcommand{\csDmLgPsi}{0.0334}
\newcommand{\csDmLgPsiBench}{0.0026}
\newcommand{\csDmLgSkill}{0.252}
\newcommand{\csDmLgSpP}{0.061}
\newcommand{\csDmLlAuc}{0.7842}
\newcommand{\csDmLlAucHi}{0.7956}
\newcommand{\csDmLlAucLo}{0.7729}
\newcommand{\csDmLlBrier}{0.1467}
\newcommand{\csDmLlDlDiff}{+0.0000}
\newcommand{\csDmLlDlP}{0.990}
\newcommand{\csDmLlEce}{0.0142}
\newcommand{\csDmLlGini}{0.5685}
\newcommand{\csDmLlHlP}{0.002}
\newcommand{\csDmLlKs}{0.4558}
\newcommand{\csDmLlPsi}{0.0324}
\newcommand{\csDmLlPsiBench}{0.0029}
\newcommand{\csDmLlSkill}{0.233}
\newcommand{\csDmLlSpP}{0.013}
\newcommand{\csDmN}{9{,}000}
\newcommand{\csDmPlAuc}{0.7966}
\newcommand{\csDmPlAucHi}{0.8075}
\newcommand{\csDmPlAucLo}{0.7857}
\newcommand{\csDmPlBrier}{0.1436}
\newcommand{\csDmPlDlDiff}{+0.0123}
\newcommand{\csDmPlDlP}{\ensuremath{3.2\times10^{-10}}}
\newcommand{\csDmPlEce}{0.0134}
\newcommand{\csDmPlGini}{0.5932}
\newcommand{\csDmPlHlP}{0.075}
\newcommand{\csDmPlKs}{0.4704}
\newcommand{\csDmPlPsi}{0.0325}
\newcommand{\csDmPlPsiBench}{0.0029}
\newcommand{\csDmPlSkill}{0.248}
\newcommand{\csDmPlSpP}{0.055}
\newcommand{\csDmPzAuc}{0.7340}
\newcommand{\csDmPzAucHi}{0.7459}
\newcommand{\csDmPzAucLo}{0.7220}
\newcommand{\csDmPzBrier}{0.1493}
\newcommand{\csDmPzDlDiff}{$-$0.0503}
\newcommand{\csDmPzDlP}{\ensuremath{2.5\times10^{-38}}}
\newcommand{\csDmPzEce}{0.0063}
\newcommand{\csDmPzGini}{0.4679}
\newcommand{\csDmPzHlP}{0.894}
\newcommand{\csDmPzKs}{0.4273}
\newcommand{\csDmPzPsi}{0.0451}
\newcommand{\csDmPzPsiBench}{0.0013}
\newcommand{\csDmPzSkill}{0.219}
\newcommand{\csDmPzSpP}{0.336}
\newcommand{\csDmRate}{25.73}
\newcommand{\csDmScAuc}{0.7842}
\newcommand{\csDmScAucHi}{0.7956}
\newcommand{\csDmScAucLo}{0.7729}
\newcommand{\csDmScBrier}{0.1467}
\newcommand{\csDmScEce}{0.0143}
\newcommand{\csDmScGini}{0.5685}
\newcommand{\csDmScHalf}{0.0114}
\newcommand{\csDmScHlP}{0.002}
\newcommand{\csDmScKs}{0.4559}
\newcommand{\csDmScPsi}{0.0325}
\newcommand{\csDmScPsiBench}{0.0029}
\newcommand{\csDmScSkill}{0.233}
\newcommand{\csDmScSpP}{0.012}
\newcommand{\csDmXgAuc}{0.7973}
\newcommand{\csDmXgAucHi}{0.8081}
\newcommand{\csDmXgAucLo}{0.7865}
\newcommand{\csDmXgBrier}{0.1429}
\newcommand{\csDmXgDlDiff}{+0.0131}
\newcommand{\csDmXgDlP}{\ensuremath{2.4\times10^{-11}}}
\newcommand{\csDmXgEce}{0.0130}
\newcommand{\csDmXgGini}{0.5946}
\newcommand{\csDmXgHlP}{0.138}
\newcommand{\csDmXgKs}{0.4669}
\newcommand{\csDmXgPsi}{0.0330}
\newcommand{\csDmXgPsiBench}{0.0024}
\newcommand{\csDmXgSkill}{0.252}
\newcommand{\csDmXgSpP}{0.187}
\newcommand{\csDsCnAuc}{0.5000}
\newcommand{\csDsCnAucHi}{0.5000}
\newcommand{\csDsCnAucLo}{0.5000}
\newcommand{\csDsCnBrier}{0.2530}
\newcommand{\csDsCnEce}{0.1448}
\newcommand{\csDsCnGini}{0.0000}
\newcommand{\csDsCnSkill}{-0.090}
\newcommand{\csDsCnSpP}{\ensuremath{1.5\times10^{-240}}}
\newcommand{\csDsEbAuc}{0.8230}
\newcommand{\csDsEbAucHi}{0.8319}
\newcommand{\csDsEbAucLo}{0.8140}
\newcommand{\csDsEbBrier}{0.1589}
\newcommand{\csDsEbDlDiff}{+0.0052}
\newcommand{\csDsEbDlP}{\ensuremath{3.1\times10^{-5}}}
\newcommand{\csDsEbEce}{0.0181}
\newcommand{\csDsEbGini}{0.6460}
\newcommand{\csDsEbHlP}{\ensuremath{1.4\times10^{-4}}}
\newcommand{\csDsEbKs}{0.5469}
\newcommand{\csDsEbPsi}{0.3380}
\newcommand{\csDsEbPsiBench}{0.0029}
\newcommand{\csDsEbSkill}{0.315}
\newcommand{\csDsEbSpP}{0.186}
\newcommand{\csDsEvents}{3{,}294}
\newcommand{\csDsLgAuc}{0.8277}
\newcommand{\csDsLgAucHi}{0.8365}
\newcommand{\csDsLgAucLo}{0.8189}
\newcommand{\csDsLgBrier}{0.1567}
\newcommand{\csDsLgDlDiff}{+0.0099}
\newcommand{\csDsLgDlP}{\ensuremath{1.4\times10^{-12}}}
\newcommand{\csDsLgEce}{0.0137}
\newcommand{\csDsLgGini}{0.6554}
\newcommand{\csDsLgHlP}{0.041}
\newcommand{\csDsLgKs}{0.5389}
\newcommand{\csDsLgPsi}{0.3395}
\newcommand{\csDsLgPsiBench}{0.0026}
\newcommand{\csDsLgSkill}{0.325}
\newcommand{\csDsLgSpP}{0.012}
\newcommand{\csDsLlAuc}{0.8178}
\newcommand{\csDsLlAucHi}{0.8270}
\newcommand{\csDsLlAucLo}{0.8086}
\newcommand{\csDsLlBrier}{0.1620}
\newcommand{\csDsLlDlDiff}{+0.0000}
\newcommand{\csDsLlDlP}{0.660}
\newcommand{\csDsLlEce}{0.0203}
\newcommand{\csDsLlGini}{0.6356}
\newcommand{\csDsLlHlP}{\ensuremath{2.3\times10^{-10}}}
\newcommand{\csDsLlKs}{0.5316}
\newcommand{\csDsLlPsi}{0.3415}
\newcommand{\csDsLlPsiBench}{0.0029}
\newcommand{\csDsLlSkill}{0.302}
\newcommand{\csDsLlSpP}{\ensuremath{1.6\times10^{-4}}}
\newcommand{\csDsN}{9{,}000}
\newcommand{\csDsPlAuc}{0.8244}
\newcommand{\csDsPlAucHi}{0.8333}
\newcommand{\csDsPlAucLo}{0.8155}
\newcommand{\csDsPlBrier}{0.1596}
\newcommand{\csDsPlDlDiff}{+0.0066}
\newcommand{\csDsPlDlP}{\ensuremath{2.8\times10^{-6}}}
\newcommand{\csDsPlEce}{0.0193}
\newcommand{\csDsPlGini}{0.6488}
\newcommand{\csDsPlHlP}{0.002}
\newcommand{\csDsPlKs}{0.5410}
\newcommand{\csDsPlPsi}{0.3277}
\newcommand{\csDsPlPsiBench}{0.0029}
\newcommand{\csDsPlSkill}{0.312}
\newcommand{\csDsPlSpP}{0.022}
\newcommand{\csDsPzAuc}{0.7888}
\newcommand{\csDsPzAucHi}{0.7982}
\newcommand{\csDsPzAucLo}{0.7794}
\newcommand{\csDsPzBrier}{0.1650}
\newcommand{\csDsPzDlDiff}{$-$0.0290}
\newcommand{\csDsPzDlP}{\ensuremath{4.0\times10^{-34}}}
\newcommand{\csDsPzEce}{0.0186}
\newcommand{\csDsPzGini}{0.5776}
\newcommand{\csDsPzHlP}{0.020}
\newcommand{\csDsPzKs}{0.5046}
\newcommand{\csDsPzPsi}{0.4108}
\newcommand{\csDsPzPsiBench}{0.0013}
\newcommand{\csDsPzSkill}{0.289}
\newcommand{\csDsPzSpP}{0.634}
\newcommand{\csDsRate}{36.60}
\newcommand{\csDsScAuc}{0.8178}
\newcommand{\csDsScAucHi}{0.8270}
\newcommand{\csDsScAucLo}{0.8086}
\newcommand{\csDsScBrier}{0.1621}
\newcommand{\csDsScEce}{0.0206}
\newcommand{\csDsScGini}{0.6356}
\newcommand{\csDsScHalf}{0.0092}
\newcommand{\csDsScHlP}{\ensuremath{8.0\times10^{-11}}}
\newcommand{\csDsScKs}{0.5318}
\newcommand{\csDsScPsi}{0.3414}
\newcommand{\csDsScPsiBench}{0.0029}
\newcommand{\csDsScSkill}{0.302}
\newcommand{\csDsScSpP}{\ensuremath{1.3\times10^{-4}}}
\newcommand{\csDsXgAuc}{0.8275}
\newcommand{\csDsXgAucHi}{0.8363}
\newcommand{\csDsXgAucLo}{0.8187}
\newcommand{\csDsXgBrier}{0.1568}
\newcommand{\csDsXgDlDiff}{+0.0097}
\newcommand{\csDsXgDlP}{\ensuremath{1.5\times10^{-13}}}
\newcommand{\csDsXgEce}{0.0172}
\newcommand{\csDsXgGini}{0.6550}
\newcommand{\csDsXgHlP}{0.003}
\newcommand{\csDsXgKs}{0.5408}
\newcommand{\csDsXgPsi}{0.3349}
\newcommand{\csDsXgPsiBench}{0.0024}
\newcommand{\csDsXgSkill}{0.324}
\newcommand{\csDsXgSpP}{0.049}
\newcommand{\csEbInteractions}{5}
\newcommand{\csEbMaxBins}{64}
\newcommand{\csEbMaxIntBins}{16}
\newcommand{\csEbOuterBags}{8}
\newcommand{\csEbmConstrained}{13}
\newcommand{\csEbmGap}{0}
\newcommand{\csEbmMonotone}{13}
\newcommand{\csEbmPairs}{5}
\newcommand{\csFitGeCal}{140}
\newcommand{\csFitGeCells}{64}
\newcommand{\csFitGeDropIv}{7}
\newcommand{\csFitGeDropSign}{0}
\newcommand{\csFitGeDropped}{7}
\newcommand{\csFitGeEbKnots}{20}
\newcommand{\csFitGeEbTerms}{22}
\newcommand{\csFitGeHold}{300}
\newcommand{\csFitGeInputs}{17}
\newcommand{\csFitGeIvMax}{0.612}
\newcommand{\csFitGeKept}{10}
\newcommand{\csFitGeLgConstrained}{2}
\newcommand{\csFitGeLgGrid}{30}
\newcommand{\csFitGeLgKnots}{20}
\newcommand{\csFitGeLgRounds}{235}
\newcommand{\csFitGeLgRows}{200}
\newcommand{\csFitGeLgViol}{0}
\newcommand{\csFitGeLlC}{1}
\newcommand{\csFitGeLlSel}{10}
\newcommand{\csFitGeLlSigns}{agree}
\newcommand{\csFitGeNeg}{10}
\newcommand{\csFitGePlBoot}{20}
\newcommand{\csFitGePlC}{0.3}
\newcommand{\csFitGePlCand}{275}
\newcommand{\csFitGePlJm}{0.280}
\newcommand{\csFitGePlJmin}{0.196}
\newcommand{\csFitGePlSel}{32}
\newcommand{\csFitGeReasonBoot}{10}
\newcommand{\csFitGeReasonMean}{0.731}
\newcommand{\csFitGeReasonMin}{0.624}
\newcommand{\csFitGeReasonRows}{200}
\newcommand{\csFitGeTrain}{560}
\newcommand{\csFitSmCal}{4{,}200}
\newcommand{\csFitSmCells}{68}
\newcommand{\csFitSmDropIv}{6}
\newcommand{\csFitSmDropSign}{4}
\newcommand{\csFitSmDropped}{10}
\newcommand{\csFitSmEbKnots}{56}
\newcommand{\csFitSmEbTerms}{25}
\newcommand{\csFitSmHold}{9{,}000}
\newcommand{\csFitSmInputs}{20}
\newcommand{\csFitSmIvMax}{0.923}
\newcommand{\csFitSmKept}{10}
\newcommand{\csFitSmLgConstrained}{13}
\newcommand{\csFitSmLgGrid}{30}
\newcommand{\csFitSmLgKnots}{56}
\newcommand{\csFitSmLgRounds}{67}
\newcommand{\csFitSmLgRows}{200}
\newcommand{\csFitSmLgViol}{0}
\newcommand{\csFitSmLlC}{0.3}
\newcommand{\csFitSmLlSel}{12}
\newcommand{\csFitSmLlSigns}{disagree}
\newcommand{\csFitSmNeg}{10}
\newcommand{\csFitSmPlBoot}{20}
\newcommand{\csFitSmPlC}{0.3}
\newcommand{\csFitSmPlCand}{467}
\newcommand{\csFitSmPlJm}{0.349}
\newcommand{\csFitSmPlJmin}{0.281}
\newcommand{\csFitSmPlSel}{89}
\newcommand{\csFitSmReasonBoot}{10}
\newcommand{\csFitSmReasonMean}{0.719}
\newcommand{\csFitSmReasonMin}{0.662}
\newcommand{\csFitSmReasonRows}{200}
\newcommand{\csFitSmTrain}{16{,}800}
\newcommand{\csFitTwCal}{4{,}200}
\newcommand{\csFitTwCells}{95}
\newcommand{\csFitTwDropIv}{7}
\newcommand{\csFitTwDropSign}{0}
\newcommand{\csFitTwDropped}{7}
\newcommand{\csFitTwEbKnots}{58}
\newcommand{\csFitTwEbTerms}{25}
\newcommand{\csFitTwHold}{9{,}000}
\newcommand{\csFitTwInputs}{20}
\newcommand{\csFitTwIvMax}{0.886}
\newcommand{\csFitTwKept}{13}
\newcommand{\csFitTwLgConstrained}{13}
\newcommand{\csFitTwLgGrid}{30}
\newcommand{\csFitTwLgKnots}{56}
\newcommand{\csFitTwLgRounds}{94}
\newcommand{\csFitTwLgRows}{200}
\newcommand{\csFitTwLgViol}{0}
\newcommand{\csFitTwLlC}{1}
\newcommand{\csFitTwLlSel}{13}
\newcommand{\csFitTwLlSigns}{agree}
\newcommand{\csFitTwNeg}{13}
\newcommand{\csFitTwPlBoot}{20}
\newcommand{\csFitTwPlC}{0.1}
\newcommand{\csFitTwPlCand}{449}
\newcommand{\csFitTwPlJm}{0.411}
\newcommand{\csFitTwPlJmin}{0.367}
\newcommand{\csFitTwPlSel}{101}
\newcommand{\csFitTwReasonBoot}{10}
\newcommand{\csFitTwReasonMean}{0.760}
\newcommand{\csFitTwReasonMin}{0.725}
\newcommand{\csFitTwReasonRows}{200}
\newcommand{\csFitTwTrain}{16{,}800}
\newcommand{\csFitTwXgKnots}{66}
\newcommand{\csFitTwXgRounds}{127}
\newcommand{\csFitTwXgViol}{0}
\newcommand{\csGtCnAuc}{0.5000}
\newcommand{\csGtCnAucHi}{0.5000}
\newcommand{\csGtCnAucLo}{0.5000}
\newcommand{\csGtCnBrier}{0.2100}
\newcommand{\csGtCnEce}{0.0467}
\newcommand{\csGtCnGini}{0.0000}
\newcommand{\csGtCnSkill}{0.000}
\newcommand{\csGtCnSpP}{0.500}
\newcommand{\csGtEbAuc}{0.7874}
\newcommand{\csGtEbAucHi}{0.8436}
\newcommand{\csGtEbAucLo}{0.7312}
\newcommand{\csGtEbBrier}{0.1668}
\newcommand{\csGtEbDlDiff}{$-$0.0009}
\newcommand{\csGtEbDlP}{0.955}
\newcommand{\csGtEbEce}{0.0578}
\newcommand{\csGtEbGini}{0.5748}
\newcommand{\csGtEbHlP}{\ensuremath{4.1\times10^{-14}}}
\newcommand{\csGtEbKs}{0.4651}
\newcommand{\csGtEbPsi}{0.0455}
\newcommand{\csGtEbPsiBench}{0.0720}
\newcommand{\csGtEbSkill}{0.206}
\newcommand{\csGtEbSpP}{0.007}
\newcommand{\csGtEvents}{90}
\newcommand{\csGtLgAuc}{0.7784}
\newcommand{\csGtLgAucHi}{0.8346}
\newcommand{\csGtLgAucLo}{0.7221}
\newcommand{\csGtLgBrier}{0.1735}
\newcommand{\csGtLgDlDiff}{$-$0.0099}
\newcommand{\csGtLgDlP}{0.615}
\newcommand{\csGtLgEce}{0.0762}
\newcommand{\csGtLgGini}{0.5568}
\newcommand{\csGtLgHlP}{\ensuremath{<10^{-300}}}
\newcommand{\csGtLgKs}{0.4254}
\newcommand{\csGtLgPsi}{0.0104}
\newcommand{\csGtLgPsiBench}{0.0645}
\newcommand{\csGtLgSkill}{0.174}
\newcommand{\csGtLgSpP}{\ensuremath{2.2\times10^{-7}}}
\newcommand{\csGtLlAuc}{0.7908}
\newcommand{\csGtLlAucHi}{0.8457}
\newcommand{\csGtLlAucLo}{0.7359}
\newcommand{\csGtLlBrier}{0.1614}
\newcommand{\csGtLlDlDiff}{+0.0025}
\newcommand{\csGtLlDlP}{0.124}
\newcommand{\csGtLlEce}{0.0517}
\newcommand{\csGtLlGini}{0.5817}
\newcommand{\csGtLlHlP}{0.654}
\newcommand{\csGtLlKs}{0.4587}
\newcommand{\csGtLlPsi}{0.0557}
\newcommand{\csGtLlPsiBench}{0.0866}
\newcommand{\csGtLlSkill}{0.231}
\newcommand{\csGtLlSpP}{0.617}
\newcommand{\csGtN}{300}
\newcommand{\csGtPlAuc}{0.7861}
\newcommand{\csGtPlAucHi}{0.8399}
\newcommand{\csGtPlAucLo}{0.7323}
\newcommand{\csGtPlBrier}{0.1658}
\newcommand{\csGtPlDlDiff}{$-$0.0022}
\newcommand{\csGtPlDlP}{0.870}
\newcommand{\csGtPlEce}{0.0475}
\newcommand{\csGtPlGini}{0.5723}
\newcommand{\csGtPlHlP}{0.597}
\newcommand{\csGtPlKs}{0.4460}
\newcommand{\csGtPlPsi}{0.0232}
\newcommand{\csGtPlPsiBench}{0.0866}
\newcommand{\csGtPlSkill}{0.210}
\newcommand{\csGtPlSpP}{0.371}
\newcommand{\csGtPzAuc}{0.7233}
\newcommand{\csGtPzAucHi}{0.7822}
\newcommand{\csGtPzAucLo}{0.6644}
\newcommand{\csGtPzBrier}{0.1821}
\newcommand{\csGtPzDlDiff}{$-$0.0650}
\newcommand{\csGtPzDlP}{0.019}
\newcommand{\csGtPzEce}{0.0735}
\newcommand{\csGtPzGini}{0.4466}
\newcommand{\csGtPzHlP}{0.328}
\newcommand{\csGtPzKs}{0.3984}
\newcommand{\csGtPzPsi}{0.0011}
\newcommand{\csGtPzPsiBench}{0.0400}
\newcommand{\csGtPzSkill}{0.133}
\newcommand{\csGtPzSpP}{0.713}
\newcommand{\csGtRate}{30.00}
\newcommand{\csGtScAuc}{0.7883}
\newcommand{\csGtScAucHi}{0.8433}
\newcommand{\csGtScAucLo}{0.7333}
\newcommand{\csGtScBrier}{0.1620}
\newcommand{\csGtScEce}{0.0455}
\newcommand{\csGtScGini}{0.5766}
\newcommand{\csGtScHalf}{0.0550}
\newcommand{\csGtScHlP}{0.876}
\newcommand{\csGtScKs}{0.4587}
\newcommand{\csGtScPsi}{0.0470}
\newcommand{\csGtScPsiBench}{0.0866}
\newcommand{\csGtScSkill}{0.229}
\newcommand{\csGtScSpP}{0.407}
\newcommand{\csGtTpAuc}{0.7769}
\newcommand{\csGtTpAucHi}{0.8331}
\newcommand{\csGtTpAucLo}{0.7206}
\newcommand{\csGtTpBrier}{0.1667}
\newcommand{\csGtTpDlDiff}{$-$0.0115}
\newcommand{\csGtTpDlP}{0.488}
\newcommand{\csGtTpEce}{0.0395}
\newcommand{\csGtTpGini}{0.5537}
\newcommand{\csGtTpHlP}{\ensuremath{<10^{-300}}}
\newcommand{\csGtTpKs}{0.4222}
\newcommand{\csGtTpPsi}{0.0367}
\newcommand{\csGtTpPsiBench}{0.0720}
\newcommand{\csGtTpSkill}{0.206}
\newcommand{\csGtTpSpP}{0.008}
\newcommand{\csHoBestAucGain}{0.0183}
\newcommand{\csHoBestName}{LightGBM}
\newcommand{\csHoCnAuc}{0.5000}
\newcommand{\csHoCnAucHi}{0.5000}
\newcommand{\csHoCnAucLo}{0.5000}
\newcommand{\csHoCnBrier}{0.1723}
\newcommand{\csHoCnEce}{0.0174}
\newcommand{\csHoCnGini}{0.0000}
\newcommand{\csHoCnSkill}{-0.000}
\newcommand{\csHoCnSpP}{0.497}
\newcommand{\csHoEbAuc}{0.7746}
\newcommand{\csHoEbAucHi}{0.7866}
\newcommand{\csHoEbAucLo}{0.7626}
\newcommand{\csHoEbBrier}{0.1367}
\newcommand{\csHoEbDlDiff}{+0.0159}
\newcommand{\csHoEbDlP}{\ensuremath{2.7\times10^{-10}}}
\newcommand{\csHoEbEce}{0.0118}
\newcommand{\csHoEbGini}{0.5492}
\newcommand{\csHoEbHlP}{0.079}
\newcommand{\csHoEbKs}{0.4258}
\newcommand{\csHoEbPsi}{0.0013}
\newcommand{\csHoEbPsiBench}{0.0029}
\newcommand{\csHoEbSkill}{0.207}
\newcommand{\csHoEbSpP}{0.231}
\newcommand{\csHoEceEbCal}{0.0118}
\newcommand{\csHoEceEbRaw}{0.0218}
\newcommand{\csHoEceLgCal}{0.0134}
\newcommand{\csHoEceLgRaw}{0.0121}
\newcommand{\csHoEceXgCal}{0.0124}
\newcommand{\csHoEceXgRaw}{0.0131}
\newcommand{\csHoEvents}{1{,}991}
\newcommand{\csHoLgAuc}{0.7770}
\newcommand{\csHoLgAucHi}{0.7889}
\newcommand{\csHoLgAucLo}{0.7652}
\newcommand{\csHoLgBrier}{0.1361}
\newcommand{\csHoLgDlDiff}{+0.0183}
\newcommand{\csHoLgDlP}{\ensuremath{6.5\times10^{-13}}}
\newcommand{\csHoLgEce}{0.0134}
\newcommand{\csHoLgGini}{0.5541}
\newcommand{\csHoLgHlP}{0.008}
\newcommand{\csHoLgKs}{0.4200}
\newcommand{\csHoLgPsi}{0.0004}
\newcommand{\csHoLgPsiBench}{0.0026}
\newcommand{\csHoLgSkill}{0.210}
\newcommand{\csHoLgSpP}{0.115}
\newcommand{\csHoLlAuc}{0.7587}
\newcommand{\csHoLlAucHi}{0.7713}
\newcommand{\csHoLlAucLo}{0.7460}
\newcommand{\csHoLlBrier}{0.1393}
\newcommand{\csHoLlDlDiff}{$-$0.0000}
\newcommand{\csHoLlDlP}{0.088}
\newcommand{\csHoLlEce}{0.0153}
\newcommand{\csHoLlGini}{0.5174}
\newcommand{\csHoLlHlP}{0.014}
\newcommand{\csHoLlKs}{0.3973}
\newcommand{\csHoLlPsi}{0.0009}
\newcommand{\csHoLlPsiBench}{0.0029}
\newcommand{\csHoLlSkill}{0.192}
\newcommand{\csHoLlSpP}{0.024}
\newcommand{\csHoN}{9{,}000}
\newcommand{\csHoPlAuc}{0.7749}
\newcommand{\csHoPlAucHi}{0.7870}
\newcommand{\csHoPlAucLo}{0.7629}
\newcommand{\csHoPlBrier}{0.1367}
\newcommand{\csHoPlDlDiff}{+0.0162}
\newcommand{\csHoPlDlP}{\ensuremath{2.6\times10^{-11}}}
\newcommand{\csHoPlEce}{0.0113}
\newcommand{\csHoPlGini}{0.5499}
\newcommand{\csHoPlHlP}{0.184}
\newcommand{\csHoPlKs}{0.4215}
\newcommand{\csHoPlPsi}{0.0008}
\newcommand{\csHoPlPsiBench}{0.0029}
\newcommand{\csHoPlSkill}{0.207}
\newcommand{\csHoPlSpP}{0.084}
\newcommand{\csHoPzAuc}{0.6917}
\newcommand{\csHoPzAucHi}{0.7047}
\newcommand{\csHoPzAucLo}{0.6787}
\newcommand{\csHoPzBrier}{0.1431}
\newcommand{\csHoPzDlDiff}{$-$0.0670}
\newcommand{\csHoPzDlP}{\ensuremath{1.3\times10^{-40}}}
\newcommand{\csHoPzEce}{0.0144}
\newcommand{\csHoPzGini}{0.3834}
\newcommand{\csHoPzHlP}{0.021}
\newcommand{\csHoPzKs}{0.3599}
\newcommand{\csHoPzPsi}{0.0006}
\newcommand{\csHoPzPsiBench}{0.0013}
\newcommand{\csHoPzSkill}{0.170}
\newcommand{\csHoPzSpP}{0.110}
\newcommand{\csHoRate}{22.12}
\newcommand{\csHoScAuc}{0.7587}
\newcommand{\csHoScAucHi}{0.7714}
\newcommand{\csHoScAucLo}{0.7461}
\newcommand{\csHoScBrier}{0.1393}
\newcommand{\csHoScEce}{0.0155}
\newcommand{\csHoScGini}{0.5174}
\newcommand{\csHoScHalf}{0.0127}
\newcommand{\csHoScHlP}{0.015}
\newcommand{\csHoScKs}{0.3973}
\newcommand{\csHoScOverPz}{0.0670}
\newcommand{\csHoScPsi}{0.0010}
\newcommand{\csHoScPsiBench}{0.0029}
\newcommand{\csHoScSkill}{0.192}
\newcommand{\csHoScSpP}{0.022}
\newcommand{\csHoXgAuc}{0.7755}
\newcommand{\csHoXgAucHi}{0.7875}
\newcommand{\csHoXgAucLo}{0.7635}
\newcommand{\csHoXgBrier}{0.1359}
\newcommand{\csHoXgDlDiff}{+0.0168}
\newcommand{\csHoXgDlP}{\ensuremath{2.2\times10^{-12}}}
\newcommand{\csHoXgEce}{0.0124}
\newcommand{\csHoXgGini}{0.5510}
\newcommand{\csHoXgHlP}{0.125}
\newcommand{\csHoXgKs}{0.4223}
\newcommand{\csHoXgPsi}{0.0010}
\newcommand{\csHoXgPsiBench}{0.0024}
\newcommand{\csHoXgSkill}{0.211}
\newcommand{\csHoXgSpP}{0.240}
\newcommand{\csLnCnAuc}{0.5000}
\newcommand{\csLnCnAucHi}{0.5000}
\newcommand{\csLnCnAucLo}{0.5000}
\newcommand{\csLnCnBrier}{0.1893}
\newcommand{\csLnCnEce}{0.0320}
\newcommand{\csLnCnGini}{0.0000}
\newcommand{\csLnCnSkill}{-0.005}
\newcommand{\csLnCnSpP}{\ensuremath{1.6\times10^{-12}}}
\newcommand{\csLnEbAuc}{0.7314}
\newcommand{\csLnEbAucHi}{0.7438}
\newcommand{\csLnEbAucLo}{0.7189}
\newcommand{\csLnEbBrier}{0.1600}
\newcommand{\csLnEbDlDiff}{+0.0139}
\newcommand{\csLnEbDlP}{\ensuremath{2.9\times10^{-8}}}
\newcommand{\csLnEbEce}{0.0341}
\newcommand{\csLnEbGini}{0.4628}
\newcommand{\csLnEbHlP}{\ensuremath{1.7\times10^{-25}}}
\newcommand{\csLnEbKs}{0.3569}
\newcommand{\csLnEbPsi}{0.0013}
\newcommand{\csLnEbPsiBench}{0.0029}
\newcommand{\csLnEbSkill}{0.150}
\newcommand{\csLnEbSpP}{\ensuremath{4.1\times10^{-28}}}
\newcommand{\csLnEvents}{2{,}265}
\newcommand{\csLnLgAuc}{0.7340}
\newcommand{\csLnLgAucHi}{0.7464}
\newcommand{\csLnLgAucLo}{0.7217}
\newcommand{\csLnLgBrier}{0.1594}
\newcommand{\csLnLgDlDiff}{+0.0166}
\newcommand{\csLnLgDlP}{\ensuremath{9.2\times10^{-11}}}
\newcommand{\csLnLgEce}{0.0374}
\newcommand{\csLnLgGini}{0.4680}
\newcommand{\csLnLgHlP}{\ensuremath{1.7\times10^{-40}}}
\newcommand{\csLnLgKs}{0.3558}
\newcommand{\csLnLgPsi}{0.0004}
\newcommand{\csLnLgPsiBench}{0.0026}
\newcommand{\csLnLgSkill}{0.154}
\newcommand{\csLnLgSpP}{\ensuremath{9.0\times10^{-32}}}
\newcommand{\csLnLlAuc}{0.7174}
\newcommand{\csLnLlAucHi}{0.7302}
\newcommand{\csLnLlAucLo}{0.7046}
\newcommand{\csLnLlBrier}{0.1626}
\newcommand{\csLnLlDlDiff}{$-$0.0000}
\newcommand{\csLnLlDlP}{0.069}
\newcommand{\csLnLlEce}{0.0420}
\newcommand{\csLnLlGini}{0.4348}
\newcommand{\csLnLlHlP}{\ensuremath{1.2\times10^{-28}}}
\newcommand{\csLnLlKs}{0.3315}
\newcommand{\csLnLlPsi}{0.0009}
\newcommand{\csLnLlPsiBench}{0.0029}
\newcommand{\csLnLlSkill}{0.137}
\newcommand{\csLnLlSpP}{\ensuremath{1.3\times10^{-32}}}
\newcommand{\csLnN}{9{,}000}
\newcommand{\csLnPlAuc}{0.7321}
\newcommand{\csLnPlAucHi}{0.7446}
\newcommand{\csLnPlAucLo}{0.7197}
\newcommand{\csLnPlBrier}{0.1601}
\newcommand{\csLnPlDlDiff}{+0.0147}
\newcommand{\csLnPlDlP}{\ensuremath{1.4\times10^{-9}}}
\newcommand{\csLnPlEce}{0.0410}
\newcommand{\csLnPlGini}{0.4643}
\newcommand{\csLnPlHlP}{\ensuremath{3.5\times10^{-27}}}
\newcommand{\csLnPlKs}{0.3529}
\newcommand{\csLnPlPsi}{0.0008}
\newcommand{\csLnPlPsiBench}{0.0029}
\newcommand{\csLnPlSkill}{0.150}
\newcommand{\csLnPlSpP}{\ensuremath{4.4\times10^{-31}}}
\newcommand{\csLnPzAuc}{0.6573}
\newcommand{\csLnPzAucHi}{0.6697}
\newcommand{\csLnPzAucLo}{0.6450}
\newcommand{\csLnPzBrier}{0.1661}
\newcommand{\csLnPzDlDiff}{$-$0.0601}
\newcommand{\csLnPzDlP}{\ensuremath{5.3\times10^{-33}}}
\newcommand{\csLnPzEce}{0.0379}
\newcommand{\csLnPzGini}{0.3147}
\newcommand{\csLnPzHlP}{\ensuremath{1.7\times10^{-24}}}
\newcommand{\csLnPzKs}{0.2943}
\newcommand{\csLnPzPsi}{0.0006}
\newcommand{\csLnPzPsiBench}{0.0013}
\newcommand{\csLnPzSkill}{0.118}
\newcommand{\csLnPzSpP}{\ensuremath{2.1\times10^{-26}}}
\newcommand{\csLnRate}{25.17}
\newcommand{\csLnScAuc}{0.7174}
\newcommand{\csLnScAucHi}{0.7302}
\newcommand{\csLnScAucLo}{0.7046}
\newcommand{\csLnScBrier}{0.1626}
\newcommand{\csLnScEce}{0.0421}
\newcommand{\csLnScGini}{0.4349}
\newcommand{\csLnScHalf}{0.0128}
\newcommand{\csLnScHlP}{\ensuremath{7.8\times10^{-29}}}
\newcommand{\csLnScKs}{0.3315}
\newcommand{\csLnScPsi}{0.0010}
\newcommand{\csLnScPsiBench}{0.0029}
\newcommand{\csLnScSkill}{0.137}
\newcommand{\csLnScSpP}{\ensuremath{7.6\times10^{-33}}}
\newcommand{\csLnXgAuc}{0.7327}
\newcommand{\csLnXgAucHi}{0.7451}
\newcommand{\csLnXgAucLo}{0.7203}
\newcommand{\csLnXgBrier}{0.1591}
\newcommand{\csLnXgDlDiff}{+0.0153}
\newcommand{\csLnXgDlP}{\ensuremath{1.4\times10^{-10}}}
\newcommand{\csLnXgEce}{0.0355}
\newcommand{\csLnXgGini}{0.4654}
\newcommand{\csLnXgHlP}{\ensuremath{2.8\times10^{-30}}}
\newcommand{\csLnXgKs}{0.3587}
\newcommand{\csLnXgPsi}{0.0010}
\newcommand{\csLnXgPsiBench}{0.0024}
\newcommand{\csLnXgSkill}{0.155}
\newcommand{\csLnXgSpP}{\ensuremath{3.2\times10^{-28}}}
\newcommand{\csPointsFeatures}{13}
\newcommand{\csPointsMatch}{50}
\newcommand{\csPointsMonotone}{13}
\newcommand{\csPointsN}{50}
\newcommand{\csPointsSpanLimit}{13}
\newcommand{\csPointsSpanPayZero}{59}
\newcommand{\csPsCnAuc}{0.5000}
\newcommand{\csPsCnAucHi}{0.5000}
\newcommand{\csPsCnAucLo}{0.5000}
\newcommand{\csPsCnBrier}{0.2339}
\newcommand{\csPsCnEce}{0.1106}
\newcommand{\csPsCnGini}{0.0000}
\newcommand{\csPsCnSkill}{-0.055}
\newcommand{\csPsCnSpP}{\ensuremath{2.9\times10^{-141}}}
\newcommand{\csPsEbAuc}{0.7838}
\newcommand{\csPsEbAucHi}{0.7940}
\newcommand{\csPsEbAucLo}{0.7736}
\newcommand{\csPsEbBrier}{0.1741}
\newcommand{\csPsEbDlDiff}{+0.0149}
\newcommand{\csPsEbDlP}{\ensuremath{1.1\times10^{-12}}}
\newcommand{\csPsEbEce}{0.0801}
\newcommand{\csPsEbGini}{0.5676}
\newcommand{\csPsEbHlP}{\ensuremath{1.6\times10^{-95}}}
\newcommand{\csPsEbKs}{0.4401}
\newcommand{\csPsEbPsi}{0.0189}
\newcommand{\csPsEbPsiBench}{0.0029}
\newcommand{\csPsEbSkill}{0.215}
\newcommand{\csPsEbSpP}{\ensuremath{8.0\times10^{-45}}}
\newcommand{\csPsEvents}{2{,}986}
\newcommand{\csPsLgAuc}{0.7852}
\newcommand{\csPsLgAucHi}{0.7953}
\newcommand{\csPsLgAucLo}{0.7750}
\newcommand{\csPsLgBrier}{0.1734}
\newcommand{\csPsLgDlDiff}{+0.0163}
\newcommand{\csPsLgDlP}{\ensuremath{9.5\times10^{-14}}}
\newcommand{\csPsLgEce}{0.0789}
\newcommand{\csPsLgGini}{0.5704}
\newcommand{\csPsLgHlP}{\ensuremath{1.1\times10^{-91}}}
\newcommand{\csPsLgKs}{0.4378}
\newcommand{\csPsLgPsi}{0.0149}
\newcommand{\csPsLgPsiBench}{0.0026}
\newcommand{\csPsLgSkill}{0.218}
\newcommand{\csPsLgSpP}{\ensuremath{6.2\times10^{-51}}}
\newcommand{\csPsLlAuc}{0.7689}
\newcommand{\csPsLlAucHi}{0.7796}
\newcommand{\csPsLlAucLo}{0.7581}
\newcommand{\csPsLlBrier}{0.1801}
\newcommand{\csPsLlDlDiff}{$-$0.0000}
\newcommand{\csPsLlDlP}{0.108}
\newcommand{\csPsLlEce}{0.0879}
\newcommand{\csPsLlGini}{0.5377}
\newcommand{\csPsLlHlP}{\ensuremath{5.1\times10^{-109}}}
\newcommand{\csPsLlKs}{0.4164}
\newcommand{\csPsLlPsi}{0.0144}
\newcommand{\csPsLlPsiBench}{0.0029}
\newcommand{\csPsLlSkill}{0.187}
\newcommand{\csPsLlSpP}{\ensuremath{1.3\times10^{-64}}}
\newcommand{\csPsN}{9{,}000}
\newcommand{\csPsPlAuc}{0.7811}
\newcommand{\csPsPlAucHi}{0.7915}
\newcommand{\csPsPlAucLo}{0.7707}
\newcommand{\csPsPlBrier}{0.1764}
\newcommand{\csPsPlDlDiff}{+0.0122}
\newcommand{\csPsPlDlP}{\ensuremath{9.0\times10^{-10}}}
\newcommand{\csPsPlEce}{0.0870}
\newcommand{\csPsPlGini}{0.5622}
\newcommand{\csPsPlHlP}{\ensuremath{1.5\times10^{-100}}}
\newcommand{\csPsPlKs}{0.4405}
\newcommand{\csPsPlPsi}{0.0161}
\newcommand{\csPsPlPsiBench}{0.0029}
\newcommand{\csPsPlSkill}{0.205}
\newcommand{\csPsPlSpP}{\ensuremath{4.1\times10^{-56}}}
\newcommand{\csPsPzAuc}{0.7061}
\newcommand{\csPsPzAucHi}{0.7169}
\newcommand{\csPsPzAucLo}{0.6953}
\newcommand{\csPsPzBrier}{0.1876}
\newcommand{\csPsPzDlDiff}{$-$0.0628}
\newcommand{\csPsPzDlP}{\ensuremath{9.3\times10^{-52}}}
\newcommand{\csPsPzEce}{0.0910}
\newcommand{\csPsPzGini}{0.4123}
\newcommand{\csPsPzHlP}{\ensuremath{1.6\times10^{-114}}}
\newcommand{\csPsPzKs}{0.3800}
\newcommand{\csPsPzPsi}{0.0113}
\newcommand{\csPsPzPsiBench}{0.0013}
\newcommand{\csPsPzSkill}{0.154}
\newcommand{\csPsPzSpP}{\ensuremath{2.7\times10^{-69}}}
\newcommand{\csPsRate}{33.18}
\newcommand{\csPsScAuc}{0.7689}
\newcommand{\csPsScAucHi}{0.7796}
\newcommand{\csPsScAucLo}{0.7581}
\newcommand{\csPsScBrier}{0.1801}
\newcommand{\csPsScEce}{0.0879}
\newcommand{\csPsScGini}{0.5378}
\newcommand{\csPsScHalf}{0.0107}
\newcommand{\csPsScHlP}{\ensuremath{2.2\times10^{-109}}}
\newcommand{\csPsScKs}{0.4166}
\newcommand{\csPsScPsi}{0.0144}
\newcommand{\csPsScPsiBench}{0.0029}
\newcommand{\csPsScSkill}{0.187}
\newcommand{\csPsScSpP}{\ensuremath{6.4\times10^{-65}}}
\newcommand{\csPsXgAuc}{0.7840}
\newcommand{\csPsXgAucHi}{0.7942}
\newcommand{\csPsXgAucLo}{0.7738}
\newcommand{\csPsXgBrier}{0.1735}
\newcommand{\csPsXgDlDiff}{+0.0151}
\newcommand{\csPsXgDlP}{\ensuremath{9.1\times10^{-14}}}
\newcommand{\csPsXgEce}{0.0789}
\newcommand{\csPsXgGini}{0.5680}
\newcommand{\csPsXgHlP}{\ensuremath{5.9\times10^{-87}}}
\newcommand{\csPsXgKs}{0.4381}
\newcommand{\csPsXgPsi}{0.0149}
\newcommand{\csPsXgPsiBench}{0.0024}
\newcommand{\csPsXgSkill}{0.218}
\newcommand{\csPsXgSpP}{\ensuremath{1.7\times10^{-47}}}
\newcommand{\csRjAccept}{60}
\newcommand{\csRjBands}{10}
\newcommand{\csRjBoot}{2{,}000}
\newcommand{\csRjCinv}{1}
\newcommand{\csRjDev}{700}
\newcommand{\csRjFloor}{0.05}
\newcommand{\csRjHoneHolds}{holds}
\newcommand{\csRjHoneRung}{Accepted only}
\newcommand{\csRjHthree}{0.01}
\newcommand{\csRjHthreeHolds}{does not hold}
\newcommand{\csRjHthreeRung}{Parcelling, prudence 2}
\newcommand{\csRjHtwo}{25}
\newcommand{\csRjHtwoHolds}{holds}
\newcommand{\csRjHtwoRung}{Parcelling, prudence 2}
\newcommand{\csRjMarAccRate}{18.2}
\newcommand{\csRjMarAoAuc}{0.7613}
\newcommand{\csRjMarAoBeats}{0}
\newcommand{\csRjMarAoBias}{$-$1.32}
\newcommand{\csRjMarAoBiasHi}{$-$0.81}
\newcommand{\csRjMarAoBiasLo}{$-$1.82}
\newcommand{\csRjMarAoBrier}{0.1777}
\newcommand{\csRjMarAoDAuc}{+0.0000}
\newcommand{\csRjMarAoDAucHi}{+0.0000}
\newcommand{\csRjMarAoDAucLo}{+0.0000}
\newcommand{\csRjMarAoRed}{0.0}
\newcommand{\csRjMarAoRedHi}{0.0}
\newcommand{\csRjMarAoRedLo}{0.0}
\newcommand{\csRjMarAuAuc}{0.7502}
\newcommand{\csRjMarAuBeats}{60}
\newcommand{\csRjMarAuBias}{$-$0.10}
\newcommand{\csRjMarAuBiasHi}{+0.45}
\newcommand{\csRjMarAuBiasLo}{$-$0.62}
\newcommand{\csRjMarAuBrier}{0.1849}
\newcommand{\csRjMarAuDAuc}{$-$0.0111}
\newcommand{\csRjMarAuDAucHi}{$-$0.0086}
\newcommand{\csRjMarAuDAucLo}{$-$0.0138}
\newcommand{\csRjMarAuRed}{92.7}
\newcommand{\csRjMarAuRedHi}{99.1}
\newcommand{\csRjMarAuRedLo}{47.1}
\newcommand{\csRjMarFuAuc}{0.7613}
\newcommand{\csRjMarFuBeats}{52}
\newcommand{\csRjMarFuBias}{$-$1.32}
\newcommand{\csRjMarFuBiasHi}{$-$0.83}
\newcommand{\csRjMarFuBiasLo}{$-$1.80}
\newcommand{\csRjMarFuBrier}{0.1777}
\newcommand{\csRjMarFuDAuc}{$-$0.0000}
\newcommand{\csRjMarFuDAucHi}{+0.0000}
\newcommand{\csRjMarFuDAucLo}{$-$0.0001}
\newcommand{\csRjMarFuRed}{0.1}
\newcommand{\csRjMarFuRedHi}{0.3}
\newcommand{\csRjMarFuRedLo}{-0.0}
\newcommand{\csRjMarHcAuc}{0.7525}
\newcommand{\csRjMarHcBeats}{42}
\newcommand{\csRjMarHcBias}{$-$1.06}
\newcommand{\csRjMarHcBiasHi}{$-$0.51}
\newcommand{\csRjMarHcBiasLo}{$-$1.62}
\newcommand{\csRjMarHcBrier}{0.1928}
\newcommand{\csRjMarHcDAuc}{$-$0.0088}
\newcommand{\csRjMarHcDAucHi}{$-$0.0074}
\newcommand{\csRjMarHcDAucLo}{$-$0.0103}
\newcommand{\csRjMarHcRed}{19.3}
\newcommand{\csRjMarHcRedHi}{39.3}
\newcommand{\csRjMarHcRedLo}{9.2}
\newcommand{\csRjMarKappa}{0}
\newcommand{\csRjMarOrAuc}{0.7827}
\newcommand{\csRjMarOrBeats}{73}
\newcommand{\csRjMarOrBias}{+0.09}
\newcommand{\csRjMarOrBiasHi}{+0.40}
\newcommand{\csRjMarOrBiasLo}{$-$0.21}
\newcommand{\csRjMarOrBrier}{0.1676}
\newcommand{\csRjMarOrDAuc}{+0.0214}
\newcommand{\csRjMarOrDAucHi}{+0.0248}
\newcommand{\csRjMarOrDAucLo}{+0.0179}
\newcommand{\csRjMarOrRed}{93.4}
\newcommand{\csRjMarOrRedHi}{99.6}
\newcommand{\csRjMarOrRedLo}{57.5}
\newcommand{\csRjMarPaAuc}{0.7533}
\newcommand{\csRjMarPaBeats}{43}
\newcommand{\csRjMarPaBias}{$-$0.77}
\newcommand{\csRjMarPaBiasHi}{$-$0.14}
\newcommand{\csRjMarPaBiasLo}{$-$1.37}
\newcommand{\csRjMarPaBrier}{0.1825}
\newcommand{\csRjMarPaDAuc}{$-$0.0080}
\newcommand{\csRjMarPaDAucHi}{$-$0.0058}
\newcommand{\csRjMarPaDAucLo}{$-$0.0103}
\newcommand{\csRjMarPaRed}{42.0}
\newcommand{\csRjMarPaRedHi}{82.8}
\newcommand{\csRjMarPaRedLo}{18.7}
\newcommand{\csRjMarPop}{30.0}
\newcommand{\csRjMarPpAuc}{0.7655}
\newcommand{\csRjMarPpBeats}{6}
\newcommand{\csRjMarPpBias}{+10.07}
\newcommand{\csRjMarPpBiasHi}{+10.69}
\newcommand{\csRjMarPpBiasLo}{+9.46}
\newcommand{\csRjMarPpBrier}{0.2037}
\newcommand{\csRjMarPpDAuc}{+0.0042}
\newcommand{\csRjMarPpDAucHi}{+0.0061}
\newcommand{\csRjMarPpDAucLo}{+0.0024}
\newcommand{\csRjMarPpRed}{-662.9}
\newcommand{\csRjMarPpRedHi}{-429.3}
\newcommand{\csRjMarPpRedLo}{-1190.2}
\newcommand{\csRjMarRejRate}{47.7}
\newcommand{\csRjMinAcc}{10}
\newcommand{\csRjModAccRate}{8.9}
\newcommand{\csRjModAoAuc}{0.6856}
\newcommand{\csRjModAoBeats}{0}
\newcommand{\csRjModAoBias}{$-$17.19}
\newcommand{\csRjModAoBiasHi}{$-$16.70}
\newcommand{\csRjModAoBiasLo}{$-$17.67}
\newcommand{\csRjModAoBrier}{0.2237}
\newcommand{\csRjModAoDAuc}{+0.0000}
\newcommand{\csRjModAoDAucHi}{+0.0000}
\newcommand{\csRjModAoDAucLo}{+0.0000}
\newcommand{\csRjModAoRed}{0.0}
\newcommand{\csRjModAoRedHi}{0.0}
\newcommand{\csRjModAoRedLo}{0.0}
\newcommand{\csRjModAuAuc}{0.6491}
\newcommand{\csRjModAuBeats}{57}
\newcommand{\csRjModAuBias}{$-$16.93}
\newcommand{\csRjModAuBiasHi}{$-$16.35}
\newcommand{\csRjModAuBiasLo}{$-$17.55}
\newcommand{\csRjModAuBrier}{0.2295}
\newcommand{\csRjModAuDAuc}{$-$0.0365}
\newcommand{\csRjModAuDAucHi}{$-$0.0313}
\newcommand{\csRjModAuDAucLo}{$-$0.0421}
\newcommand{\csRjModAuRed}{1.5}
\newcommand{\csRjModAuRedHi}{2.8}
\newcommand{\csRjModAuRedLo}{0.1}
\newcommand{\csRjModFuAuc}{0.6855}
\newcommand{\csRjModFuBeats}{54}
\newcommand{\csRjModFuBias}{$-$17.18}
\newcommand{\csRjModFuBiasHi}{$-$16.70}
\newcommand{\csRjModFuBiasLo}{$-$17.68}
\newcommand{\csRjModFuBrier}{0.2237}
\newcommand{\csRjModFuDAuc}{$-$0.0002}
\newcommand{\csRjModFuDAucHi}{$-$0.0001}
\newcommand{\csRjModFuDAucLo}{$-$0.0002}
\newcommand{\csRjModFuRed}{0.0}
\newcommand{\csRjModFuRedHi}{0.0}
\newcommand{\csRjModFuRedLo}{-0.0}
\newcommand{\csRjModHcAuc}{0.6544}
\newcommand{\csRjModHcBeats}{58}
\newcommand{\csRjModHcBias}{$-$17.08}
\newcommand{\csRjModHcBiasHi}{$-$16.54}
\newcommand{\csRjModHcBiasLo}{$-$17.60}
\newcommand{\csRjModHcBrier}{0.2366}
\newcommand{\csRjModHcDAuc}{$-$0.0313}
\newcommand{\csRjModHcDAucHi}{$-$0.0285}
\newcommand{\csRjModHcDAucLo}{$-$0.0342}
\newcommand{\csRjModHcRed}{0.6}
\newcommand{\csRjModHcRedHi}{1.3}
\newcommand{\csRjModHcRedLo}{-0.1}
\newcommand{\csRjModKappa}{0.5}
\newcommand{\csRjModOrAuc}{0.7827}
\newcommand{\csRjModOrBeats}{100}
\newcommand{\csRjModOrBias}{+0.09}
\newcommand{\csRjModOrBiasHi}{+0.39}
\newcommand{\csRjModOrBiasLo}{$-$0.21}
\newcommand{\csRjModOrBrier}{0.1676}
\newcommand{\csRjModOrDAuc}{+0.0971}
\newcommand{\csRjModOrDAucHi}{+0.1065}
\newcommand{\csRjModOrDAucLo}{+0.0870}
\newcommand{\csRjModOrRed}{99.5}
\newcommand{\csRjModOrRedHi}{100.0}
\newcommand{\csRjModOrRedLo}{97.7}
\newcommand{\csRjModPaAuc}{0.6695}
\newcommand{\csRjModPaBeats}{66}
\newcommand{\csRjModPaBias}{$-$16.79}
\newcommand{\csRjModPaBiasHi}{$-$16.21}
\newcommand{\csRjModPaBiasLo}{$-$17.35}
\newcommand{\csRjModPaBrier}{0.2256}
\newcommand{\csRjModPaDAuc}{$-$0.0161}
\newcommand{\csRjModPaDAucHi}{$-$0.0113}
\newcommand{\csRjModPaDAucLo}{$-$0.0214}
\newcommand{\csRjModPaRed}{2.3}
\newcommand{\csRjModPaRedHi}{3.5}
\newcommand{\csRjModPaRedLo}{1.0}
\newcommand{\csRjModPop}{30.0}
\newcommand{\csRjModPpAuc}{0.7064}
\newcommand{\csRjModPpBeats}{100}
\newcommand{\csRjModPpBias}{$-$9.57}
\newcommand{\csRjModPpBiasHi}{$-$8.77}
\newcommand{\csRjModPpBiasLo}{$-$10.41}
\newcommand{\csRjModPpBrier}{0.2079}
\newcommand{\csRjModPpDAuc}{+0.0208}
\newcommand{\csRjModPpDAucHi}{+0.0243}
\newcommand{\csRjModPpDAucLo}{+0.0171}
\newcommand{\csRjModPpRed}{44.3}
\newcommand{\csRjModPpRedHi}{47.6}
\newcommand{\csRjModPpRedLo}{41.2}
\newcommand{\csRjModRejRate}{61.7}
\newcommand{\csRjPrudence}{2}
\newcommand{\csRjReps}{100}
\newcommand{\csRjRows}{1000}
\newcommand{\csRjStrAccRate}{2.9}
\newcommand{\csRjStrAoAuc}{0.5388}
\newcommand{\csRjStrAoBeats}{0}
\newcommand{\csRjStrAoBias}{$-$26.69}
\newcommand{\csRjStrAoBiasHi}{$-$26.49}
\newcommand{\csRjStrAoBiasLo}{$-$26.89}
\newcommand{\csRjStrAoBrier}{0.2801}
\newcommand{\csRjStrAoDAuc}{+0.0000}
\newcommand{\csRjStrAoDAucHi}{+0.0000}
\newcommand{\csRjStrAoDAucLo}{+0.0000}
\newcommand{\csRjStrAoRed}{0.0}
\newcommand{\csRjStrAoRedHi}{0.0}
\newcommand{\csRjStrAoRedLo}{0.0}
\newcommand{\csRjStrAuAuc}{0.5017}
\newcommand{\csRjStrAuBeats}{25}
\newcommand{\csRjStrAuBias}{$-$26.91}
\newcommand{\csRjStrAuBiasHi}{$-$26.72}
\newcommand{\csRjStrAuBiasLo}{$-$27.11}
\newcommand{\csRjStrAuBrier}{0.2831}
\newcommand{\csRjStrAuDAuc}{$-$0.0371}
\newcommand{\csRjStrAuDAucHi}{$-$0.0327}
\newcommand{\csRjStrAuDAucLo}{$-$0.0413}
\newcommand{\csRjStrAuRed}{-0.8}
\newcommand{\csRjStrAuRedHi}{-0.6}
\newcommand{\csRjStrAuRedLo}{-1.1}
\newcommand{\csRjStrFuAuc}{0.5389}
\newcommand{\csRjStrFuBeats}{53}
\newcommand{\csRjStrFuBias}{$-$26.69}
\newcommand{\csRjStrFuBiasHi}{$-$26.49}
\newcommand{\csRjStrFuBiasLo}{$-$26.88}
\newcommand{\csRjStrFuBrier}{0.2801}
\newcommand{\csRjStrFuDAuc}{+0.0001}
\newcommand{\csRjStrFuDAucHi}{+0.0002}
\newcommand{\csRjStrFuDAucLo}{$-$0.0001}
\newcommand{\csRjStrFuRed}{0.0}
\newcommand{\csRjStrFuRedHi}{0.0}
\newcommand{\csRjStrFuRedLo}{-0.0}
\newcommand{\csRjStrHcAuc}{0.5132}
\newcommand{\csRjStrHcBeats}{42}
\newcommand{\csRjStrHcBias}{$-$26.75}
\newcommand{\csRjStrHcBiasHi}{$-$26.54}
\newcommand{\csRjStrHcBiasLo}{$-$26.96}
\newcommand{\csRjStrHcBrier}{0.2831}
\newcommand{\csRjStrHcDAuc}{$-$0.0256}
\newcommand{\csRjStrHcDAucHi}{$-$0.0220}
\newcommand{\csRjStrHcDAucLo}{$-$0.0293}
\newcommand{\csRjStrHcRed}{-0.2}
\newcommand{\csRjStrHcRedHi}{-0.0}
\newcommand{\csRjStrHcRedLo}{-0.4}
\newcommand{\csRjStrKappa}{1}
\newcommand{\csRjStrOrAuc}{0.7827}
\newcommand{\csRjStrOrBeats}{100}
\newcommand{\csRjStrOrBias}{+0.09}
\newcommand{\csRjStrOrBiasHi}{+0.40}
\newcommand{\csRjStrOrBiasLo}{$-$0.21}
\newcommand{\csRjStrOrBrier}{0.1676}
\newcommand{\csRjStrOrDAuc}{+0.2439}
\newcommand{\csRjStrOrDAucHi}{+0.2594}
\newcommand{\csRjStrOrDAucLo}{+0.2288}
\newcommand{\csRjStrOrRed}{99.7}
\newcommand{\csRjStrOrRedHi}{100.0}
\newcommand{\csRjStrOrRedLo}{98.5}
\newcommand{\csRjStrPaAuc}{0.5307}
\newcommand{\csRjStrPaBeats}{54}
\newcommand{\csRjStrPaBias}{$-$26.63}
\newcommand{\csRjStrPaBiasHi}{$-$26.39}
\newcommand{\csRjStrPaBiasLo}{$-$26.84}
\newcommand{\csRjStrPaBrier}{0.2808}
\newcommand{\csRjStrPaDAuc}{$-$0.0081}
\newcommand{\csRjStrPaDAucHi}{$-$0.0016}
\newcommand{\csRjStrPaDAucLo}{$-$0.0145}
\newcommand{\csRjStrPaRed}{0.2}
\newcommand{\csRjStrPaRedHi}{0.6}
\newcommand{\csRjStrPaRedLo}{-0.1}
\newcommand{\csRjStrPop}{30.0}
\newcommand{\csRjStrPpAuc}{0.5621}
\newcommand{\csRjStrPpBeats}{98}
\newcommand{\csRjStrPpBias}{$-$24.99}
\newcommand{\csRjStrPpBiasHi}{$-$24.62}
\newcommand{\csRjStrPpBiasLo}{$-$25.34}
\newcommand{\csRjStrPpBrier}{0.2723}
\newcommand{\csRjStrPpDAuc}{+0.0233}
\newcommand{\csRjStrPpDAucHi}{+0.0295}
\newcommand{\csRjStrPpDAucLo}{+0.0174}
\newcommand{\csRjStrPpRed}{6.4}
\newcommand{\csRjStrPpRedHi}{7.2}
\newcommand{\csRjStrPpRedLo}{5.6}
\newcommand{\csRjStrRejRate}{70.7}
\newcommand{\csRjTest}{300}
\newcommand{\csRjWeakAccRate}{13.3}
\newcommand{\csRjWeakAoAuc}{0.7398}
\newcommand{\csRjWeakAoBeats}{0}
\newcommand{\csRjWeakAoBias}{$-$9.31}
\newcommand{\csRjWeakAoBiasHi}{$-$8.76}
\newcommand{\csRjWeakAoBiasLo}{$-$9.89}
\newcommand{\csRjWeakAoBrier}{0.1913}
\newcommand{\csRjWeakAoDAuc}{+0.0000}
\newcommand{\csRjWeakAoDAucHi}{+0.0000}
\newcommand{\csRjWeakAoDAucLo}{+0.0000}
\newcommand{\csRjWeakAoRed}{0.0}
\newcommand{\csRjWeakAoRedHi}{0.0}
\newcommand{\csRjWeakAoRedLo}{0.0}
\newcommand{\csRjWeakAuAuc}{0.7122}
\newcommand{\csRjWeakAuBeats}{72}
\newcommand{\csRjWeakAuBias}{$-$8.52}
\newcommand{\csRjWeakAuBiasHi}{$-$7.89}
\newcommand{\csRjWeakAuBiasLo}{$-$9.11}
\newcommand{\csRjWeakAuBrier}{0.1989}
\newcommand{\csRjWeakAuDAuc}{$-$0.0276}
\newcommand{\csRjWeakAuDAucHi}{$-$0.0241}
\newcommand{\csRjWeakAuDAucLo}{$-$0.0311}
\newcommand{\csRjWeakAuRed}{8.4}
\newcommand{\csRjWeakAuRedHi}{11.5}
\newcommand{\csRjWeakAuRedLo}{5.6}
\newcommand{\csRjWeakFuAuc}{0.7398}
\newcommand{\csRjWeakFuBeats}{55}
\newcommand{\csRjWeakFuBias}{$-$9.31}
\newcommand{\csRjWeakFuBiasHi}{$-$8.74}
\newcommand{\csRjWeakFuBiasLo}{$-$9.84}
\newcommand{\csRjWeakFuBrier}{0.1913}
\newcommand{\csRjWeakFuDAuc}{$-$0.0000}
\newcommand{\csRjWeakFuDAucHi}{+0.0000}
\newcommand{\csRjWeakFuDAucLo}{$-$0.0001}
\newcommand{\csRjWeakFuRed}{0.0}
\newcommand{\csRjWeakFuRedHi}{0.0}
\newcommand{\csRjWeakFuRedLo}{-0.0}
\newcommand{\csRjWeakHcAuc}{0.7186}
\newcommand{\csRjWeakHcBeats}{59}
\newcommand{\csRjWeakHcBias}{$-$9.14}
\newcommand{\csRjWeakHcBiasHi}{$-$8.55}
\newcommand{\csRjWeakHcBiasLo}{$-$9.75}
\newcommand{\csRjWeakHcBrier}{0.2071}
\newcommand{\csRjWeakHcDAuc}{$-$0.0213}
\newcommand{\csRjWeakHcDAucHi}{$-$0.0192}
\newcommand{\csRjWeakHcDAucLo}{$-$0.0235}
\newcommand{\csRjWeakHcRed}{1.8}
\newcommand{\csRjWeakHcRedHi}{3.1}
\newcommand{\csRjWeakHcRedLo}{0.5}
\newcommand{\csRjWeakKappa}{0.25}
\newcommand{\csRjWeakOrAuc}{0.7827}
\newcommand{\csRjWeakOrBeats}{97}
\newcommand{\csRjWeakOrBias}{+0.09}
\newcommand{\csRjWeakOrBiasHi}{+0.41}
\newcommand{\csRjWeakOrBiasLo}{$-$0.22}
\newcommand{\csRjWeakOrBrier}{0.1676}
\newcommand{\csRjWeakOrDAuc}{+0.0429}
\newcommand{\csRjWeakOrDAucHi}{+0.0483}
\newcommand{\csRjWeakOrDAucLo}{+0.0376}
\newcommand{\csRjWeakOrRed}{99.1}
\newcommand{\csRjWeakOrRedHi}{99.9}
\newcommand{\csRjWeakOrRedLo}{95.5}
\newcommand{\csRjWeakPaAuc}{0.7286}
\newcommand{\csRjWeakPaBeats}{74}
\newcommand{\csRjWeakPaBias}{$-$8.51}
\newcommand{\csRjWeakPaBiasHi}{$-$7.89}
\newcommand{\csRjWeakPaBiasLo}{$-$9.13}
\newcommand{\csRjWeakPaBrier}{0.1946}
\newcommand{\csRjWeakPaDAuc}{$-$0.0112}
\newcommand{\csRjWeakPaDAucHi}{$-$0.0082}
\newcommand{\csRjWeakPaDAucLo}{$-$0.0142}
\newcommand{\csRjWeakPaRed}{8.6}
\newcommand{\csRjWeakPaRedHi}{11.7}
\newcommand{\csRjWeakPaRedLo}{5.8}
\newcommand{\csRjWeakPop}{30.0}
\newcommand{\csRjWeakPpAuc}{0.7492}
\newcommand{\csRjWeakPpBeats}{88}
\newcommand{\csRjWeakPpBias}{+1.24}
\newcommand{\csRjWeakPpBiasHi}{+1.98}
\newcommand{\csRjWeakPpBiasLo}{+0.48}
\newcommand{\csRjWeakPpBrier}{0.1927}
\newcommand{\csRjWeakPpDAuc}{+0.0094}
\newcommand{\csRjWeakPpDAucHi}{+0.0116}
\newcommand{\csRjWeakPpDAucLo}{+0.0072}
\newcommand{\csRjWeakPpRed}{86.7}
\newcommand{\csRjWeakPpRedHi}{95.1}
\newcommand{\csRjWeakPpRedLo}{77.5}
\newcommand{\csRjWeakRejRate}{55.1}
\newcommand{\csScFactor}{28.85}
\newcommand{\csScIntercept}{-1.252}
\newcommand{\csScOffset}{487.12}
\newcommand{\csSsCnAuc}{0.5000}
\newcommand{\csSsCnAucHi}{0.5000}
\newcommand{\csSsCnAucLo}{0.5000}
\newcommand{\csSsCnBrier}{0.1723}
\newcommand{\csSsCnEce}{0.0174}
\newcommand{\csSsCnGini}{0.0000}
\newcommand{\csSsCnSkill}{-0.000}
\newcommand{\csSsCnSpP}{0.480}
\newcommand{\csSsEbAuc}{0.7649}
\newcommand{\csSsEbAucHi}{0.7773}
\newcommand{\csSsEbAucLo}{0.7525}
\newcommand{\csSsEbBrier}{0.1389}
\newcommand{\csSsEbDlDiff}{+0.0124}
\newcommand{\csSsEbDlP}{\ensuremath{4.9\times10^{-4}}}
\newcommand{\csSsEbEce}{0.0166}
\newcommand{\csSsEbGini}{0.5298}
\newcommand{\csSsEbHlP}{\ensuremath{8.0\times10^{-4}}}
\newcommand{\csSsEbKs}{0.4143}
\newcommand{\csSsEbPsi}{0.0063}
\newcommand{\csSsEbPsiBench}{0.0095}
\newcommand{\csSsEbSkill}{0.194}
\newcommand{\csSsEbSpP}{0.228}
\newcommand{\csSsEvents}{1{,}991}
\newcommand{\csSsLgAuc}{0.7601}
\newcommand{\csSsLgAucHi}{0.7725}
\newcommand{\csSsLgAucLo}{0.7478}
\newcommand{\csSsLgBrier}{0.1396}
\newcommand{\csSsLgDlDiff}{+0.0077}
\newcommand{\csSsLgDlP}{0.019}
\newcommand{\csSsLgEce}{0.0172}
\newcommand{\csSsLgGini}{0.5203}
\newcommand{\csSsLgHlP}{\ensuremath{2.1\times10^{-5}}}
\newcommand{\csSsLgKs}{0.3943}
\newcommand{\csSsLgPsi}{0.0010}
\newcommand{\csSsLgPsiBench}{0.0095}
\newcommand{\csSsLgSkill}{0.190}
\newcommand{\csSsLgSpP}{0.086}
\newcommand{\csSsLlAuc}{0.7511}
\newcommand{\csSsLlAucHi}{0.7639}
\newcommand{\csSsLlAucLo}{0.7382}
\newcommand{\csSsLlBrier}{0.1404}
\newcommand{\csSsLlDlDiff}{$-$0.0014}
\newcommand{\csSsLlDlP}{0.074}
\newcommand{\csSsLlEce}{0.0135}
\newcommand{\csSsLlGini}{0.5021}
\newcommand{\csSsLlHlP}{0.010}
\newcommand{\csSsLlKs}{0.3894}
\newcommand{\csSsLlPsi}{0.0035}
\newcommand{\csSsLlPsiBench}{0.0103}
\newcommand{\csSsLlSkill}{0.185}
\newcommand{\csSsLlSpP}{\ensuremath{8.0\times10^{-4}}}
\newcommand{\csSsN}{9{,}000}
\newcommand{\csSsPlAuc}{0.7528}
\newcommand{\csSsPlAucHi}{0.7653}
\newcommand{\csSsPlAucLo}{0.7402}
\newcommand{\csSsPlBrier}{0.1424}
\newcommand{\csSsPlDlDiff}{+0.0003}
\newcommand{\csSsPlDlP}{0.936}
\newcommand{\csSsPlEce}{0.0261}
\newcommand{\csSsPlGini}{0.5055}
\newcommand{\csSsPlHlP}{\ensuremath{1.3\times10^{-21}}}
\newcommand{\csSsPlKs}{0.3812}
\newcommand{\csSsPlPsi}{0.0051}
\newcommand{\csSsPlPsiBench}{0.0103}
\newcommand{\csSsPlSkill}{0.174}
\newcommand{\csSsPlSpP}{\ensuremath{2.3\times10^{-14}}}
\newcommand{\csSsPzAuc}{0.6928}
\newcommand{\csSsPzAucHi}{0.7049}
\newcommand{\csSsPzAucLo}{0.6806}
\newcommand{\csSsPzBrier}{0.1431}
\newcommand{\csSsPzDlDiff}{$-$0.0597}
\newcommand{\csSsPzDlP}{\ensuremath{5.3\times10^{-38}}}
\newcommand{\csSsPzEce}{0.0140}
\newcommand{\csSsPzGini}{0.3855}
\newcommand{\csSsPzHlP}{0.101}
\newcommand{\csSsPzKs}{0.3599}
\newcommand{\csSsPzPsi}{0.0009}
\newcommand{\csSsPzPsiBench}{0.0048}
\newcommand{\csSsPzSkill}{0.169}
\newcommand{\csSsPzSpP}{0.041}
\newcommand{\csSsRate}{22.12}
\newcommand{\csSsScAuc}{0.7525}
\newcommand{\csSsScAucHi}{0.7653}
\newcommand{\csSsScAucLo}{0.7396}
\newcommand{\csSsScBrier}{0.1403}
\newcommand{\csSsScEce}{0.0153}
\newcommand{\csSsScGini}{0.5049}
\newcommand{\csSsScHalf}{0.0128}
\newcommand{\csSsScHlP}{\ensuremath{3.5\times10^{-5}}}
\newcommand{\csSsScKs}{0.3925}
\newcommand{\csSsScPsi}{0.0033}
\newcommand{\csSsScPsiBench}{0.0103}
\newcommand{\csSsScSkill}{0.186}
\newcommand{\csSsScSpP}{\ensuremath{2.9\times10^{-5}}}
\newcommand{\csTpEstimators}{8}


A credit score maps the characteristics $x$ of an applicant or a borrower to the probability of default over a stated horizon,
$\PD(x)=\Prob(Y=1\mid X=x)$, where $Y=1$ marks a default and $X$ the vector of inputs. The same quantity drives the approval
decision, the price of the loan, the capital of a bank that uses the internal ratings-based approach and the expected credit loss
of its provisions, so a PD model is the most common model a validation function reviews and the one it reviews most often. This
section states the methods Contraste fits and tests on three real samples: a weight-of-evidence scorecard as the reference model
(a monotone binning of each input, a logistic regression on the binned inputs and an integer points table), a penalised logistic
regression, a penalised logistic regression with tree-derived rules, an explainable boosting machine, a monotone gradient boosting
machine with an explicit calibration map, a tabular foundation model, and a family of treatments of the rejected applicants whose
outcomes are never observed.

The regulation that governs the area is the Basel framework for banks that estimate PD internally. Its minimum requirements for
statistical rating models place on the bank the burden of showing good predictive power, a reasonable set of predictors, and
estimates that are accurate on average with no known material bias; they require documentation of the theory, the assumptions and
the mathematical and empirical basis of the model, and a rigorous statistical process, including out-of-time and out-of-sample
performance tests, for validating it \citep{credit-scoring:bcbscre36}. The validation requirement compares the realised default
rate with the estimated PD of each grade and asks that the two agree within an expected range, at least once a year
\citep{credit-scoring:bcbscre36}. The PD floors of the same framework (0.05\,\% for retail exposures, 0.10\,\% for revolving
retail) bound the estimate from below \citep{credit-scoring:bcbscre32}. The Basel Committee's study of validation methods warns
that no really powerful test of calibration is available, because default events are correlated
\citep{credit-scoring:bcbswp14}, and the supervisory reporting instructions of the European Central Bank fix the statistics and
their $p$-values without prescribing pass-or-fail thresholds \citep{credit-scoring:ecb2019}. For model risk in general, the US
guidance SR~11-7 \citep{credit-scoring:sr117} shaped the practice that other supervisors echo, and was superseded in the United
States by SR~26-2 on 17 April 2026 \citep{credit-scoring:sr262}. For consumer credit in the United States, the official
interpretation of Regulation~B states that disclosing more than four reasons for an adverse action is not likely to help the
applicant, which fixes the number of principal reasons that the reason codes of the machine-learning methods below return
\citep{credit-scoring:regb}.

The published evidence on the choice of method is consistent across benchmarks, and the statistical classification methods of
consumer credit scoring have been reviewed for decades \citep{credit-scoring:handhenley1997}. On eight real credit data sets, logistic
regression and linear discriminant analysis performed very well beside support vector machines and neural networks
\citep{credit-scoring:baesens2003}; on 41 classifiers and eight retail data sets, heterogeneous ensembles beat logistic regression
significantly, with gains of a few AUC points, and the authors recommend reporting at least the AUC, the partial Gini coefficient
and the Brier score \citep{credit-scoring:lessmann2015}. Deep networks have not shown an advantage over boosted trees on tabular
credit data \citep{credit-scoring:gunnarsson2021,credit-scoring:grinsztajn2022}, and random forests and gradient boosting stay
robust when the minority class of a credit sample is strongly undersampled \citep{credit-scoring:brown2012}. The size of each gain has to be measured on the
data at hand, and calibration has to be tested beside discrimination, since a method that ranks better can still state the wrong
level. The methods below are therefore run on one protocol and compared with the scorecard by a paired test.

The section covers three cases. The first case (C01) asks whether a monotone scorecard matches the challengers the literature says
can beat it once calibration is required, and what the validation tests report when the population shifts; its data are the
Taiwan credit-card default sample, \csFitTwTrain{} training cards and \csHoN{} locked holdout cards \citep{credit-scoring:yeh2009},
and the German credit sample of 1,000 applicants as a small-sample twin \citep{credit-scoring:hofmann1994}. The second case (C02)
asks how fast discrimination falls from the one-year to the five-year horizon on the Polish company bankruptcy files
\citep{credit-scoring:tomczak2016}, whether the learned methods beat the scorecard by a margin a paired test supports, and whether
they stay calibrated. The third case (C03) asks whether a PD model built on approval cohorts old enough to have been observed for
five years stays calibrated for the cohorts that lived through the 2008--2010 recession, using the loans of the SBA 7(a) program
\citep{credit-scoring:sbafoia}.

All methods of the section share one protocol. For C01 and C02 a locked stratified holdout of 30\,\% of each file is opened once;
20\,\% of the remainder is a calibration slice that fits every calibration map and the rest is the training slice; binning,
weights of evidence, imputation, scaling, rule extraction and every hyperparameter are fitted on the training slice only, the
hyperparameters by repeated stratified five-fold cross-validation (two repeats) inside it, with seed 42. No resampling is applied
before calibration is measured, because undersampling, oversampling and synthetic minority methods produced strongly
miscalibrated risk models without a higher AUC \citep{credit-scoring:goorbergh2022}. For C03 the split is out of time by approval
fiscal year. The tests that judge every method are those of the validation engine \texttt{riskvalidation} 1.1.0, called from the
data pipeline of Contraste (release 1.00.000).

% ---------------------------------------------------------------------------------------------------------------------
\subsection{Weight-of-evidence monotone scorecard}
\label{sec:m1}

The scorecard is the reference model of the section. Its theory is the weight of evidence of a discretised input. For an input
cut into bins $i=1,\dots,n$, let $r_i^{E}$ and $r_i^{NE}$ be the numbers of defaults (events) and non-defaults in bin $i$, and
$r_T^{E}$ and $r_T^{NE}$ their totals. Following \citet{credit-scoring:navas2020}, the weight of evidence and the information
value of the input are
\begin{equation}
\mathrm{WoE}_i=\log\frac{r_i^{NE}/r_T^{NE}}{r_i^{E}/r_T^{E}},\qquad
\mathrm{IV}=\sum_{i=1}^{n}(p_i-q_i)\log\frac{p_i}{q_i},\quad p_i=\frac{r_i^{NE}}{r_T^{NE}},\; q_i=\frac{r_i^{E}}{r_T^{E}},
\label{eq:cs-woe}
\end{equation}
so that a higher weight of evidence means a lower default rate, and the information value is the symmetrised Kullback--Leibler
divergence between the distributions of non-defaults and defaults over the bins. The bins are not chosen by hand. The pre-binning
produces $n\le 20$ cells from a decision tree, and an integer linear program merges consecutive cells to maximise the information
value of the result under a monotone trend of the default rate, a minimum bin share and bounds on the number of bins
\citep{credit-scoring:navas2020}:
\begin{equation}
\max_{X}\ \sum_{i=1}^{n}\Big(V_{ii}X_{ii}+\sum_{j<i}(V_{ij}-V_{i,j+1})X_{ij}\Big)
\quad\text{s.t.}\quad \sum_{i\ge j}X_{ij}=1,\ \ X_{ij}\le X_{i,j+1},\ \ b_{\min}\le\sum_i X_{ii}\le b_{\max},
\label{eq:cs-binning}
\end{equation}
where $X_{ij}\in\{0,1\}$ is lower triangular, $X_{ij}=1$ meaning that cells $j,\dots,i$ form one bin, and $V_{ij}$ is the
information value of that merged bin; the monotonicity and bin-size constraints (omitted above) act on the merged default rates
and counts. The scorecard is then a logistic regression on the weights of evidence of the $k$ retained inputs,
\begin{equation}
\eta(x)=\beta_0+\sum_{j=1}^{k}\beta_j\,\mathrm{WoE}_j(x_j),\qquad \PD(x)=\frac{1}{1+e^{-\eta(x)}},
\label{eq:cs-logit}
\end{equation}
and, with this sign convention, every coefficient of a sound scorecard is negative. The score is a linear function of the log-odds
scaled so that a fixed number of points doubles the odds \citep{credit-scoring:siddiqi2017,credit-scoring:skorecard}:
\begin{equation}
\begin{aligned}
\text{Factor}&=\frac{\text{PDO}}{\ln 2},\qquad
\text{Offset}=S_{\mathrm{ref}}-\text{Factor}\,\ln(\text{odds}_{\mathrm{ref}}),\\
\text{Points}_{j,b}&=\frac{\text{Offset}}{k}-\Big(\beta_j\,\mathrm{WoE}_{j,b}+\frac{\beta_0}{k}\Big)\text{Factor},
\end{aligned}
\label{eq:cs-points}
\end{equation}
rounded to integers; the score of an applicant is the sum of the points of the bins its inputs fall in, which equals
$\text{Offset}-\text{Factor}\,\eta(x)$ before rounding. PDO is the number of points that doubles the odds, $S_{\mathrm{ref}}$ the
reference score and $\text{odds}_{\mathrm{ref}}$ the odds at that score.

The assumptions are four. The log-odds of default is additive in the inputs once each input is binned; the default rate is
monotone in each binned input in a direction fixed from its meaning before any fit (more months of delay raise the risk, a larger
credit limit and larger payments lower it); the training sample represents the population the model scores; and a category that
the training slice never held falls in the zero-weight ``missing'' bin. The parameters, all fixed before fitting, are a PDO of 20
points, a reference score of 600 points at odds of 50 to 1, an information-value threshold of 0.02 below which an input is
screened out (the threshold is the convention attributed to \citet{credit-scoring:siddiqi2017}; the page of the book was not
checked), a minimum bin share of 5\,\% of the records, at most 20 pre-bins and the constraint-programming solver of the binning
library \citep{credit-scoring:optbinning}. After the screen, a sign check drops, one at a time, the wrong-signed input with the
largest $p$-value and refits, until every coefficient is negative. Table~\ref{tab:cs-fit} lists the fitted configuration on the
Taiwan training slice: of \csFitTwInputs{} inputs, \csFitTwDropIv{} fall below the threshold and \csFitTwKept{} remain, with
\csFitTwCells{} cells in the points table. On the 2,000-row refit the sign check also acts, dropping \csFitSmDropSign{} inputs in
addition to \csFitSmDropIv{} screened out by information value; this is the estimation risk that a small sample adds.

The computational cost has three parts. The pre-binning sorts each input once, $O(N\log N)$ per input for $N$ training rows; the
program of Equation~\eqref{eq:cs-binning} has at most $n(n+1)/2\le 210$ binary variables per input; and the logistic fit by Newton
iterations costs $O(Nk^2)$ per iteration, with at most 100 iterations. Scoring an applicant costs $k$ bisection searches in the
split lists. The implementation is the module \texttt{pipeline/model/scorecard.py} of the data pipeline (\texttt{fit\_scorecard},
\texttt{Scorecard.points}, \texttt{table\_points}), on \texttt{optbinning} 1.0.0 and \texttt{statsmodels} 0.15.0, and the tests that
judge it are the functions \texttt{disc\_auc}, \texttt{disc\_delong}, \texttt{disc\_ks}, \texttt{pd\_hosmer\_lemeshow},
\texttt{pd\_spiegelhalter}, \texttt{pd\_brier}, \texttt{pd\_ece} and \texttt{stability\_psi} of \texttt{riskvalidation} 1.1.0.

\begin{table}[htbp]
\centering
\small
\caption{Fitted configuration of each method on the training slice of the Taiwan card-default sample (C01): the parameters chosen on training data only, the number of retained terms and the size of the slices.}
\label{tab:cs-fit}
\begin{tabular}{lp{11cm}}
\toprule
Method & Fitted configuration on the card-default training slice \\
\midrule
\relax WoE scorecard & \relax 13 of 20 inputs kept (information value at least 0.02, sign check); 95 points-table cells; PDO 20, reference 600 points at odds 50 \\
\relax L1 logistic & \relax penalty $1/C$ with $C=1$ (grid 0.003 to 1, five-fold cross-validated log loss); 13 non-zero coefficients \\
\relax PLTR & \relax $C=0.1$; 449 candidate rules; 101 selected; mean Jaccard 0.411 over 20 resamples \\
\relax EBM & \relax 25 terms (5 pairwise); at most 64 bins per main term and 16 per pair; 8 outer bags; isotonic map with 58 knots \\
\relax LightGBM & \relax 94 trees by early stopping; 15 leaves; learning rate 0.05; isotonic map with 56 knots \\
\relax XGBoost & \relax 127 trees by early stopping; depth 4; learning rate 0.05 \\
\relax Training slice & \relax 16,800 rows (3,716 defaults); calibration slice 4,200; locked holdout 9,000 \\
\bottomrule
\end{tabular}
\end{table}


The scorecard is compared with every other method by the paired DeLong test (Equation~\eqref{eq:cs-delong}), and its acceptance
criterion has three parts: the binned inputs, fitted on the training slice only, reproduce monotone default rates; the points
table recomputes every score exactly (checked on the \csPointsN{} scored applicants stored with the holdout); and the AUC with its DeLong interval, the Kolmogorov--Smirnov separation, the Brier score,
the calibration by grade and the population stability index out of sample are reported. On the Taiwan fit all
\csPointsMonotone{} of \csPointsFeatures{} retained inputs have monotone default rates over their bins, and the points table
recomputes the points of \csPointsMatch{} of the \csPointsN{} scored applicants stored with the holdout. Table~\ref{tab:cs-points}
shows the points of the two inputs with the largest information value: the September repayment status
(information value \csFitTwIvMax{}) spans \csPointsSpanPayZero{} points from its best bin to its worst, and the credit limit
spans only \csPointsSpanLimit{} points, so the ranking of the scorecard is carried by the repayment history. The score of a borrower is read directly from the table, which is what makes the
model auditable line by line.

\begin{table}[htbp]
\centering
\small
\caption{Points table of the two strongest groups of the Taiwan scorecard: the September repayment status and the credit limit, with the bin counts, default rates, weight of evidence and integer points on the training slice.}
\label{tab:cs-points}
\begin{tabular}{llrrrrr}
\toprule
Input & Bin & Rows & Defaults & Default rate & WoE & Points \\
\midrule
\relax PAY\_0 & \relax \texttt{(-inf, -1.50)} & 1,562 & 204 & \relax 0.131 & \relax +0.637 & \relax 54 \\
\relax PAY\_0 & \relax \texttt{[-1.50, 0.50)} & 11,339 & 1,554 & \relax 0.137 & \relax +0.581 & \relax 53 \\
\relax PAY\_0 & \relax \texttt{[0.50, 1.50)} & 2,112 & 713 & \relax 0.338 & \relax $-$0.585 & \relax 28 \\
\relax PAY\_0 & \relax \texttt{[1.50, inf)} & 1,787 & 1,245 & \relax 0.697 & \relax $-$2.090 & \relax $-$5 \\
\relax LIMIT\_BAL & \relax \texttt{(-inf, 45000.00)} & 2,396 & 881 & \relax 0.368 & \relax $-$0.717 & \relax 34 \\
\relax LIMIT\_BAL & \relax \texttt{[45000.00, 75000.00)} & 2,764 & 734 & \relax 0.266 & \relax $-$0.241 & \relax 38 \\
\relax LIMIT\_BAL & \relax \texttt{[75000.00, 105000.00)} & 1,820 & 441 & \relax 0.242 & \relax $-$0.119 & \relax 39 \\
\relax LIMIT\_BAL & \relax \texttt{[105000.00, 145000.00)} & 1,575 & 370 & \relax 0.235 & \relax $-$0.078 & \relax 40 \\
\relax LIMIT\_BAL & \relax \texttt{[145000.00, 165000.00)} & 1,023 & 189 & \relax 0.185 & \relax +0.226 & \relax 42 \\
\relax LIMIT\_BAL & \relax \texttt{[165000.00, 245000.00)} & 3,310 & 575 & \relax 0.174 & \relax +0.301 & \relax 43 \\
\relax LIMIT\_BAL & \relax \texttt{[245000.00, 375000.00)} & 2,577 & 371 & \relax 0.144 & \relax +0.524 & \relax 45 \\
\relax LIMIT\_BAL & \relax \texttt{[375000.00, inf)} & 1,335 & 155 & \relax 0.116 & \relax +0.771 & \relax 47 \\
\bottomrule
\end{tabular}
\end{table}


\begin{algorithm}[htbp]
\caption{Fit of the monotone weight-of-evidence scorecard on the training slice}
\label{alg:cs-scorecard}
\begin{algorithmic}[1]
\Require training inputs $X$ and defaults $y$; monotone direction $d_j\in\{-1,0,+1\}$ per input; $\mathrm{IV}_{\min}=0.02$; PDO, $S_{\mathrm{ref}}$, $\text{odds}_{\mathrm{ref}}$
\For{each input $j$}
  \State bin $x_j$ with the program of Equation~\eqref{eq:cs-binning} under trend $d_j$ (automatic ascending or descending trend if $d_j=0$)
  \State compute $\mathrm{WoE}_{j,b}$ and $\mathrm{IV}_j$ by Equation~\eqref{eq:cs-woe}
\EndFor
\State keep $\mathcal{K}=\{j:\mathrm{IV}_j\ge \mathrm{IV}_{\min}\}$
\Repeat
  \State fit the logistic regression of $y$ on $\{\mathrm{WoE}_j(x_j)\}_{j\in\mathcal{K}}$ by Newton iterations
  \State $\mathcal{W}\gets\{j\in\mathcal{K}:\hat\beta_j>0\}$
  \If{$\mathcal{W}\neq\emptyset$} \State remove from $\mathcal{K}$ the $j\in\mathcal{W}$ with the largest $p$-value \EndIf
\Until{$\mathcal{W}=\emptyset$}
\State compute $\text{Factor}$, $\text{Offset}$ and $\text{Points}_{j,b}$ by Equation~\eqref{eq:cs-points}; round to integers
\State \Return bins, $\hat\beta$, points table
\end{algorithmic}
\end{algorithm}

% ---------------------------------------------------------------------------------------------------------------------
\subsection{Penalised logistic regression}
\label{sec:m2}

The penalised logistic regression keeps the structure of the scorecard (a logistic model on the weights of evidence) and replaces
the sign-check procedure by a continuous shrinkage. Its theory is the lasso of \citet{credit-scoring:tibshirani1996} applied to the
logistic log-likelihood: with $z_i$ the vector of weights of evidence of applicant $i$ and $\ell$ the logistic loss, the estimate is
\begin{equation}
(\hat\beta_0,\hat\beta)=\arg\min_{\beta_0,\beta}\ \|\beta\|_1+C\sum_{i=1}^{N}\ell\big(y_i,\ \beta_0+z_i^{\top}\beta\big),\qquad
\ell(y,\eta)=\log\!\big(1+e^{\eta}\big)-y\,\eta,
\label{eq:cs-l1}
\end{equation}
where $C>0$ is the inverse of the penalty strength, as the library defines it. A small $C$ sets coefficients exactly to zero, so
the penalty both shrinks and selects; the weights of evidence of Equation~\eqref{eq:cs-woe} keep every coefficient interpretable
as a log-odds contribution. The inputs are the weights of evidence of every variable that passed the information-value screen,
including the variables that the sign check of the scorecard would have dropped, so the two models differ only in how they deal
with collinearity: the scorecard removes inputs by a rule, the penalty shrinks them by a continuous amount. The strength is chosen
on the training slice only, by the mean log loss over the $|\mathcal{F}|=10$ validation folds (five folds, two repeats):
\begin{equation}
\hat C=\arg\min_{C\in\mathcal{G}}\ \frac{1}{|\mathcal{F}|}\sum_{(\mathrm{tr},\mathrm{va})\in\mathcal{F}}
\mathrm{LL}\big(y_{\mathrm{va}},\ \hat p_{C}^{\,\mathrm{tr}}(z_{\mathrm{va}})\big),\qquad
\mathcal{G}=\{0.003,\,0.01,\,0.03,\,0.1,\,0.3,\,1\},
\label{eq:cs-cv}
\end{equation}
where $\mathrm{LL}$ is the mean logistic log loss and $\hat p_C^{\mathrm{tr}}$ the model fitted on the training part of a fold with
strength $C$. The final model is refitted on the whole training slice with $\hat C$ (Algorithm~\ref{alg:cs-penalised}).

The assumptions are those of the scorecard, plus sparsity: that a few inputs carry the signal, which holds for the repayment
history and the credit limit. The one parameter of the method is $C$, selected from a six-point logarithmic grid by
Equation~\eqref{eq:cs-cv}; the solver is \texttt{liblinear} with at most 2,000 iterations. On the Taiwan fit the grid maximum
$\hat C=\csFitTwLlC{}$ was selected, the weakest penalty of the grid, and all \csFitTwLlSel{} coefficients are non-zero with signs that \csFitTwLlSigns{} with the weights of evidence. On the German fit
$\hat C=\csFitGeLlC{}$ and the signs also \csFitGeLlSigns{}; on the 2,000-row Taiwan refit $\hat C=\csFitSmLlC{}$ and the signs
\csFitSmLlSigns{}, which shows that a small sample gives coefficients the scorecard's sign check would have refused. The cost of
selection is $6\times10=60$ fits of the penalised model plus the final fit, each on the $N\times k$ matrix of weights of evidence;
scoring costs $k$ table lookups and a dot product. The implementation is \texttt{pipeline/model/penalised.py}
(\texttt{fit\_penalised\_woe}) on \texttt{scikit-learn} 1.9.1 \texttt{LogisticRegression}, and the fitted record stores the
coefficients, the selected inputs and whether the signs agree with the weights of evidence.

Two properties of the estimator bear on its use as a validated model. First, the penalty biases the coefficients toward zero, so a
penalised PD is shrunk toward the portfolio mean unless the penalty is weak; neither this model nor the scorecard carries a
calibration map, because the level of a logistic model is tied to the training default rate by the likelihood, and the calibration
tests of the engine are the check on that level, not a map fitted to repair it. Second, because the model is logistic on the same
weights of evidence, it converts to a points table by Equation~\eqref{eq:cs-points} with the penalised coefficients in place of the
unpenalised ones, so it is as auditable as the scorecard; what it gives up is the sign guarantee, which the fitted record keeps as
a flag (\texttt{signs\_agree\_with\_woe}) rather than as a constraint. A coefficient with the wrong sign on a weight of evidence
signals that the input stands in for another input the model also holds, the collinearity that the sign check of the scorecard
removes by dropping the input; the penalty can leave it in with a small coefficient, which is why the flag is reported on every
fit. The expected failure modes are a penalty so weak that it adds nothing to the scorecard, as on the Taiwan fit, and a penalty so
strong that an informative input disappears; the grid of Equation~\eqref{eq:cs-cv} spans both.

The method is compared with the scorecard by the paired DeLong test for two AUCs computed on the same applicants
\citep{credit-scoring:delong1988}. With structural components $V^{(r)}_{10}$ (the share of non-defaulters ranked safer than each
defaulter) and $V^{(r)}_{01}$ (the share of defaulters ranked riskier than each non-defaulter) for model $r\in\{1,2\}$, and
$m$ defaulters and $n$ non-defaulters,
\begin{equation}
\mathbf S=\frac{S_{10}}{m}+\frac{S_{01}}{n},\qquad
z=\frac{\widehat{\mathrm{AUC}}_1-\widehat{\mathrm{AUC}}_2}{\sqrt{\mathbf S_{11}+\mathbf S_{22}-2\mathbf S_{12}}},
\label{eq:cs-delong}
\end{equation}
where $S_{10}$ and $S_{01}$ are the sample covariance matrices of the components across defaulters and across non-defaulters, and
$z$ is asymptotically standard normal under the null hypothesis that the two AUCs are equal. The AUC itself is the probability that
a defaulter is ranked riskier than a non-defaulter, with ties counting one half \citep{credit-scoring:hanley1982}; the engine
function \texttt{disc\_delong} returns the statistic, the two-sided $p$-value and the difference. The acceptance criterion of the
method is that this paired comparison against the scorecard is reported. On the Taiwan holdout the penalised model and the
scorecard have the same AUC of \csHoLlAuc{} (difference \csHoLlDlDiff{}, $p=\csHoLlDlP{}$), so the shrinkage changes nothing
measurable when cross-validation selects the weakest penalty of the grid; the value of the method is as the control that separates the effect of
a different estimator from the effect of a different model class in the comparisons of the sections that follow.

\begin{algorithm}[htbp]
\caption{Fit of the penalised logistic regression on weights of evidence, with the strength chosen by cross-validation}
\label{alg:cs-penalised}
\begin{algorithmic}[1]
\Require weights of evidence $Z$ of the inputs that passed the information-value screen, defaults $y$ of the training slice; the folds $\mathcal{F}$ (five folds, two repeats); the grid $\mathcal{G}$
\For{each $C\in\mathcal{G}$, in increasing order}
  \For{each fold $(\mathrm{tr},\mathrm{va})\in\mathcal{F}$}
    \State fit the L1 logistic regression of Equation~\eqref{eq:cs-l1} on $(Z_{\mathrm{tr}},y_{\mathrm{tr}})$ with strength $C$ (\texttt{liblinear}, at most 2,000 iterations)
    \State $\mathrm{LL}_{C,\mathrm{va}}\gets$ mean log loss of its PD on $(Z_{\mathrm{va}},y_{\mathrm{va}})$
  \EndFor
  \State $\bar{\mathrm{LL}}_C\gets$ mean of $\mathrm{LL}_{C,\mathrm{va}}$ over the folds
\EndFor
\State $\hat C\gets$ the $C$ with the smallest $\bar{\mathrm{LL}}_C$ (the first of equal values in increasing order)
\State refit on all of $(Z,y)$ with strength $\hat C$; $\;\mathcal{S}\gets\{j:\hat\beta_j\neq0\}$; $\;$ flag whether every $\hat\beta_j\le0$
\State \Return $\hat\beta_0$, $\hat\beta$, $\hat C$, $\mathcal{S}$, the sign flag
\end{algorithmic}
\end{algorithm}

% ---------------------------------------------------------------------------------------------------------------------
\subsection{Penalised logistic tree regression (PLTR)}
\label{sec:m3}

Penalised logistic tree regression adds non-linear and interaction effects to the logistic model without giving up its
interpretation. \citet{credit-scoring:dumitrescu2022} extract threshold rules from short decision trees and enter them, beside the
original variables, in a penalised logistic regression; the authors report that the method beats logistic regression significantly
and is competitive with a random forest on credit data. A rule is a conjunction of at most two threshold conditions,
\begin{equation}
r_m(x)=\prod_{(j,\,\circ,\,t)\in\mathcal{C}_m}\mathbf 1\big[x_j\ \circ\ t\big],\qquad \circ\in\{\le,>\},
\label{eq:cs-rule}
\end{equation}
where $\mathcal{C}_m$ is the set of conditions along the path to a leaf. In Contraste a tree of depth one is grown on each variable
(one split, giving one rule after the redundant leaf is dropped) and a tree of depth two on each pair of variables (up to four
leaves, three rules), so that the number of candidate rules is at most $k+3\binom{k}{2}$ for $k$ inputs; leaves hold at least 5\,\%
of the training rows, duplicate rules are merged by their text, and categorical inputs enter the trees through their weight of
evidence so that their categories are ordered by risk. The design vector stacks the standardised inputs and the rule indicators,
$u(x)=\big(\tilde x_1,\dots,\tilde x_k,\,r_1(x),\dots,r_M(x)\big)$, with $\tilde x_j=(x_j-\mu_j)/\sigma_j$ fitted on the training
slice, and the model is the L1 logistic regression of Equation~\eqref{eq:cs-l1} on $u(x)$ with strength chosen by
Equation~\eqref{eq:cs-cv}. A rule is selected when its coefficient is non-zero, $\mathcal{S}=\{m:\hat\beta_m\neq0\}$.

The stability of the selection is measured by refitting on bootstrap resamples. For resample $b$ of the training rows, with the
strength $\hat C$ kept and the rules kept, the selected set is $\mathcal{S}_b$, and
\begin{equation}
J_b=\frac{|\mathcal{S}\cap\mathcal{S}_b|}{|\mathcal{S}\cup\mathcal{S}_b|},\qquad
\bar J=\frac1B\sum_{b=1}^{B}J_b,\quad B=20,
\label{eq:cs-jaccard}
\end{equation}
the Jaccard index between the rules selected on the full training slice and on the resample, with $J_b=1$ when both sets are
empty. The acceptance criterion of the method is that the rule set is stable across 20 bootstraps and
that the Jaccard index is reported; the index is a measurement, and Section~\ref{sec:results-credit-scoring} reads its value.

The assumptions are those of the penalised logistic regression, and one more: that interactions of at most two variables and
thresholds learned by a tree capture the non-linearity that matters. The parameters are the tree depths (one and two), the
minimum leaf share (5\,\%), the strength grid of Equation~\eqref{eq:cs-cv} and $B=20$; the depths follow the construction of
\citet{credit-scoring:dumitrescu2022} as the pipeline implements it, and the leaf share is a convention that keeps every rule
supported by at least 5\,\% of the rows its tree was grown on. On the Taiwan fit the strength is $\hat C=\csFitTwPlC{}$, there
are \csFitTwPlCand{} candidate rules out of at most $20+3\cdot190=590$, and \csFitTwPlSel{} are selected. The number of trees is
$k+\binom{k}{2}=210$ for $k=20$, each fitted in $O(N\log N)$ at these depths; the penalised fit then runs on $N\times(k+M)$
columns, the strength search repeats it 60 times and the bootstrap adds 20 refits at the fixed strength. Rule extraction
dominates the run time: a pilot run of this method took 13 minutes at 100,000 rows (a figure from the design of the case, not from a committed artifact), which is why the
vintage windows of about a million loans are fitted on a stratified sample of \csCthreeFitRows{} rows. The implementation is
\texttt{pipeline/model/penalised.py} (\texttt{extract\_rules}, \texttt{fit\_pltr}) on \texttt{scikit-learn} 1.9.1
\texttt{DecisionTreeClassifier} and \texttt{LogisticRegression}, with the rule extraction written in the pipeline.

Three features of the construction shape how the result is read. The bootstrap of Equation~\eqref{eq:cs-jaccard} keeps the
candidate rules and the strength fixed, so it measures the variability of the selection only and not that of the rule
extraction; a protocol that regrew the trees on each resample would show more variability, and the index therefore understates the
instability of the whole procedure. The rule indicators are not standardised, only the original variables are, so the L1 penalty
puts a rule that covers a twentieth of the rows and a rule that covers half of them on the same scale of coefficients, which
favours the broad rules. And the prediction is additive in the log-odds, so the contribution of every variable and of every
fired rule can be listed as a reason, a property that the boosted tree ensembles of Section~\ref{sec:m5} obtain only through an
attribution method. The expected failure modes are a rule explosion, where candidates outnumber the rows that can support them,
and an unstable selection across resamples; the Jaccard index is the measurement of the second.

\begin{algorithm}[htbp]
\caption{Fit of the penalised logistic tree regression and the stability of its rule set}
\label{alg:cs-pltr}
\begin{algorithmic}[1]
\Require training inputs $X$ (categories replaced by their weight of evidence) and defaults $y$; strength grid $\mathcal{G}$; $B=20$
\State $\mathcal{R}\gets\emptyset$
\For{each variable $j$} \State grow a tree of depth one on $x_j$ with leaf share at least $5\,\%$; add its rules, without the last leaf, to $\mathcal{R}$ \EndFor
\For{each pair $(j,j')$} \State grow a tree of depth two on $(x_j,x_{j'})$ with leaf share at least $5\,\%$; add its rules, without the last leaf, to $\mathcal{R}$ \EndFor
\State merge the rules of $\mathcal{R}$ with the same text
\State build $u(x)$ from the standardised variables and the indicators of $\mathcal{R}$
\State choose $\hat C$ by Equation~\eqref{eq:cs-cv} and fit the L1 logistic regression on $u(x)$; $\mathcal{S}\gets\{m:\hat\beta_m\neq0\}$
\For{$b=1,\dots,B$}
  \State draw $N$ training rows with replacement; refit with strength $\hat C$ on the resample; $\mathcal{S}_b\gets$ its selected rules
  \State $J_b\gets|\mathcal{S}\cap\mathcal{S}_b|/|\mathcal{S}\cup\mathcal{S}_b|$ (1 if both are empty)
\EndFor
\State \Return $\hat\beta$, $\mathcal{S}$, $\bar J=\frac1B\sum_b J_b$
\end{algorithmic}
\end{algorithm}

% ---------------------------------------------------------------------------------------------------------------------
\subsection{Explainable boosting machine (GA2M)}
\label{sec:m4}

The explainable boosting machine is a generalised additive model with pairwise interactions, learned by boosting.
\citet{credit-scoring:lou2013} introduce the model and the detection of the interaction pairs, and \citet{credit-scoring:nori2019}
describe its implementation in the InterpretML package. The model form is
\begin{equation}
g\big(\E[y\mid x]\big)=\beta_0+\sum_{j}f_j(x_j)+\sum_{(i,j)\in\mathcal{P}}f_{ij}(x_i,x_j),
\label{eq:cs-ebm}
\end{equation}
where $g$ is the logit link, each $f_j$ is a shape function learned by cyclic gradient boosting (a very small step applied to one
term at a time, in turn, with bagging), and $\mathcal{P}$ is the set of detected pairs \citep{credit-scoring:interpretdocs}. For
credit the model is a boosted scorecard: each shape function is a table of log-odds contributions by bin, which is a points table
in another unit, and the contribution of each main effect to a reason code is exact, with no Shapley approximation. After the
fit, the additive logit is exported as data. With $T$ terms, $b_t(x)$ the bin of the record on the inputs of term $t$ and $s_t$ the
score array of that term,
\begin{equation}
\eta(x)=\beta_0+\sum_{t=1}^{T}s_t\big[b_t(x)\big],\qquad p(x)=\hat g\!\Big(\frac{1}{1+e^{-\eta(x)}}\Big),
\label{eq:cs-ebm-export}
\end{equation}
where $\hat g$ is the isotonic calibration map of Section~\ref{sec:m5}, fitted on the calibration slice as the nondecreasing
function that minimises $\sum_i(y_i-g(\hat p_i))^2$ over the calibration rows \citep{credit-scoring:zadrozny2002}, and clipped to
$[10^{-6},1-10^{-6}]$ so that a PD stays inside the open interval that the data contracts require. A scorer outside Python needs
only $\beta_0$, the bins and the arrays $s_t$. Algorithm~\ref{alg:cs-ebm} states the fit, the export and its checks.

The assumptions are that the response is additive up to the detected pairs, that the direction of each constrained input is known
(the same directions as for the scorecard are passed as monotone constraints), and that the pairs found on the training slice
represent the interactions of the population. The parameters are set below the library defaults on purpose: at most
\csEbMaxBins{} bins per main term (the library default is 1,024), \csEbMaxIntBins{} bins per pairwise term, \csEbInteractions{}
pairwise terms, \csEbOuterBags{} outer bags, one thread and a fixed seed. The default number of bins would produce shape
functions too large to ship as a document within the 256\,KiB size limit of the artifacts and too fine to read as a points table; the
bin counts are therefore a design choice, not a tuned value. On the Taiwan fit the model has \csFitTwEbTerms{} terms (20 main
effects and \csEbmPairs{} pairs), and the isotonic map has \csFitTwEbKnots{} knots. The cost of scoring is $T$ array lookups, and
the cost of fitting is $O(B\,R\,T\,N)$ for $B$ outer bags, $R$ boosting rounds (chosen by the early stopping of the library), $T$
terms and $N$ training rows, because each round updates one term from a pass over the rows binned on that term, plus the screening of
the candidate pairs, which this order leaves out; the export has the size of the sum of the term arrays. In every exported array, index 0 holds the missing values and the last index the categories the training slice never
held, so that a scorer outside Python handles both without a rule of its own. Pairwise terms are the part of the model that is
hardest to turn into a reason code, since the contribution depends on two inputs at once, and Contraste limits them to
\csEbInteractions{}; the monotone constraints are set on the main effects, so the check of the exported arrays below covers the main
effects and says nothing about the shape of a pairwise term.

The acceptance criterion of the method has four parts: monotone shape functions where a direction is declared; a points table per
term; the reliability of the calibrated score; and live additive scoring that equals the offline score within $10^{-9}$. The
exported arrays carry the evidence for the first and the last. Of the \csEbmConstrained{} main effects with a declared direction,
\csEbmMonotone{} are monotone over their interior bins in the exported arrays. Evaluating Equation~\eqref{eq:cs-ebm-export} from the
export alone on the \csPointsN{} scored applicants stored with the holdout reproduces their stored logits with a largest absolute
difference of \csEbmGap{}, so the additive document is the model; the live scorer of the application is held to the same document
by a separate test of the application. The reliability of the calibrated score is read in
Section~\ref{sec:results-credit-scoring}: on the holdout the expected calibration error falls from \csHoEceEbRaw{} before the map
to \csHoEceEbCal{} after it. Contraste implements the method with \texttt{InterpretML} \texttt{ExplainableBoostingClassifier}
0.7.8 (module \texttt{pipeline/model/ebm.py}, functions \texttt{fit\_ebm}, \texttt{export} and \texttt{score\_logit}).

On the paired DeLong comparison the method gains \csHoEbDlDiff{} AUC over the scorecard on the Taiwan holdout ($p=\csHoEbDlP{}$),
a gain of the same size as the penalised logistic tree regression (\csHoPlDlDiff{}), which suggests that the non-linearity of the
repayment history and the credit limit, not the choice of method, produces the gain. That reading is a hypothesis: the artifacts
do not isolate the components of the gain, and the pairwise terms are not ablated.

\begin{algorithm}[htbp]
\caption{Fit of the explainable boosting machine, its export, the isotonic map and the checks of the export}
\label{alg:cs-ebm}
\begin{algorithmic}[1]
\Require training inputs $X$ and defaults $y$; directions $c_j\in\{-1,0,+1\}$ (0 for categorical inputs); calibration slice $(X_{\mathrm{cal}},y_{\mathrm{cal}})$; at most 64 bins per main term, 16 per pair, 5 pairs, 8 outer bags, seed
\State fit the model of Equation~\eqref{eq:cs-ebm} by cyclic boosting on the training slice, with the monotone constraints on the main effects
\State export the intercept $\beta_0$, for each input its bins and for each term $t$ its array $s_t$ (index 0 for missing values, the last index for unseen categories)
\State $\hat p_i\gets\sigma(\eta(x_i))$ for the calibration rows, with $\eta$ of Equation~\eqref{eq:cs-ebm-export} evaluated from the export
\State fit the nondecreasing isotonic map $\hat g$ on $(\hat p_i,y_i)$ of the calibration slice; clip its values to $[10^{-6},1-10^{-6}]$
\For{each main effect with $c_j\neq0$} \State check that $c_j\,s_t$ is nondecreasing over the interior bins of the export \EndFor
\State evaluate $\eta$ from the export alone on the applicants stored with the holdout; report the largest absolute difference from the stored logits
\State \Return the export, $\hat g$, the monotonicity record and the difference
\end{algorithmic}
\end{algorithm}

% ---------------------------------------------------------------------------------------------------------------------
\subsection{Monotone gradient boosting with explicit calibration}
\label{sec:m5}

Gradient boosting fits an additive ensemble of trees, $F(x)=\sum_{t=1}^{T}\nu\,f_t(x)$, each tree correcting the residual of the
previous ones, and the predicted PD is $\sigma(F(x))$ with $\sigma$ the logistic function. A monotone constraint restricts the
ensemble to be non-decreasing, or non-increasing, in a chosen input. For an input $j$ with declared direction $c_j\in\{-1,+1\}$,
\begin{equation}
c_j\,\big[F(x+\delta e_j)-F(x)\big]\ \ge\ 0\qquad\text{for all }x\text{ and all }\delta>0,
\label{eq:cs-mono}
\end{equation}
where $e_j$ is the $j$th unit vector. LightGBM enforces the constraint on the splits and the leaf values; its \texttt{advanced}
method is the less over-constraining of the three it offers \citep{credit-scoring:lightgbmdocs}. XGBoost enforces it too, with the
warning that with histogram binning the constraints can wipe out split candidates, so the number of bins is raised
\citep{credit-scoring:xgboostdocs}. A constraint is verified rather than assumed: on the partial-dependence grid, for the
individual conditional expectation curves of $n_r=200$ training records and $G$ quantile levels $g_1<\dots<g_G$ of each
constrained input, the largest violation is
\begin{equation}
v_j=\max\Big(0,\ -\min_{i,\,k}\ c_j\big[F(x_i^{(j\leftarrow g_{k+1})})-F(x_i^{(j\leftarrow g_{k})})\big]\Big),
\label{eq:cs-pdp}
\end{equation}
and the monotonicity is accepted when every $v_j$ is at most $10^{-12}$; here $x^{(j\leftarrow g)}$ is the record with input $j$
set to $g$.

A boosted score is not a calibrated probability, so an explicit map is fitted on a slice that the trees never saw. The isotonic map
$\hat g$ of Section~\ref{sec:m4} turns the raw score into a PD, fitted on the calibration slice; the choice of an isotonic map
rather than a sigmoid is a convention of the pipeline, and the quantity that justifies it, the reliability of the result, is
measured on the holdout. Reason codes follow the additive attribution of the Shapley value, $F(x)=\phi_0+\sum_j\phi_j(x)$
\citep{credit-scoring:lundberg2017}, computed exactly for trees by TreeSHAP \citep{credit-scoring:lundberg2020}: the reasons of
record $i$ are the four inputs with the largest risk-raising contributions, $T_i=\operatorname{top}_4\{\phi_{ij}\}$, the practice that
the four-reason limit of Regulation~B fixes \citep{credit-scoring:regb}. A boosted model with such an attribution can be made
as transparent and auditable as a scorecard while keeping its predictive gain, in a real credit case study
\citep{credit-scoring:bucker2022}. Because correlated inputs split credit arbitrarily, the
stability of the reasons is measured by refitting on bootstrap resamples $b=1,\dots,B_r$ of the training slice and computing the
mean overlap of the top-four sets for the same holdout records,
\begin{equation}
O_b=\frac1{n_r}\sum_{i=1}^{n_r}\frac{|T_i\cap T_i^{(b)}|}{4},\qquad B_r=10,\ n_r=200.
\label{eq:cs-overlap}
\end{equation}

The assumptions are that the declared directions are true of the population, that the calibration slice is exchangeable with the
data to be scored, and that tree boosting captures the structure that matters. The parameters of LightGBM are a learning rate of
0.05, 15 leaves, at least 100 records per leaf, a feature and a bagging fraction of 0.9, an L2 penalty of 1, the \texttt{advanced}
constraint method and deterministic training; the number of trees is chosen by early stopping with a patience of 50 rounds on
the cross-validation folds of the training slice, with a ceiling of 2,000 rounds. XGBoost, which serves as the cross-check of the
engine, uses a learning rate of 0.05, depth 4, a minimum child weight of 20, subsampling of 0.9 and 512 histogram bins. On the
Taiwan fit early stopping gave \csFitTwLgRounds{} trees for LightGBM and \csFitTwXgRounds{} for XGBoost, and the isotonic maps have
\csFitTwLgKnots{} and \csFitTwXgKnots{} knots. The cost of a boosting round is linear in the number of rows and features, the
number of rounds is selected by the cross-validation, the bootstrap of the reasons adds $B_r$ refits at the selected number of
rounds, and scoring costs $T$ tree traversals. The implementation is \texttt{pipeline/model/gbm.py} (\texttt{fit\_lightgbm},
\texttt{fit\_xgboost}, \texttt{monotonicity}, \texttt{reason\_codes}, \texttt{reason\_stability}) on \texttt{LightGBM} 4.7.0 and
\texttt{XGBoost} 3.4.1, with the map in \texttt{pipeline/core/calibration.py}; Algorithm~\ref{alg:cs-gbm} states the procedure.

The acceptance criterion has four parts: monotonicity verified on the partial-dependence grids; an isotonic map fitted on a
separate slice; the Brier score and the reliability before and after the map; and the stability of the reason codes, measured as
the top-four overlap over bootstrap refits. On the Taiwan fit the largest monotone violation over the
\csFitTwLgConstrained{} constrained inputs is \csFitTwLgViol{} for LightGBM and \csFitTwXgViol{} for XGBoost, on a grid of
\csFitTwLgGrid{} levels and \csFitTwLgRows{} records. The mean top-four overlap over \csFitTwReasonBoot{} refits is \csFitTwReasonMean{} (smallest \csFitTwReasonMin{}),
that is, about three of the four principal reasons survive a refit; on the German fit the mean falls to \csFitGeReasonMean{}
(smallest \csFitGeReasonMin{}). The reliability before and after the map is read in
Section~\ref{sec:results-credit-scoring}: for LightGBM the expected calibration error is \csHoEceLgRaw{} before the map and
\csHoEceLgCal{} after it, so on this holdout the isotonic map does not improve it, whereas for the explainable boosting machine it
nearly halves it. The Brier score is \csHoLgBrier{} for LightGBM and \csHoXgBrier{} for XGBoost, against \csHoScBrier{} for the
scorecard. LightGBM is the method with the highest AUC on the holdout (\csHoLgAuc{}), a gain of \csHoLgDlDiff{} over the scorecard
($p=\csHoLgDlP{}$); XGBoost reaches \csHoXgAuc{} with the same constraints, so the gain does not depend on the engine.

\begin{algorithm}[htbp]
\caption{Fit of the monotone gradient boosting machine, its isotonic map, the partial-dependence check and the stability of the reasons}
\label{alg:cs-gbm}
\begin{algorithmic}[1]
\Require training slice $(X,y)$; directions $c_j$; folds $\mathcal{F}$; calibration slice; patience 50, ceiling 2,000 rounds; $B_r=10$, $n_r=200$
\State run the cross-validation over $\mathcal{F}$ with the monotone constraints and early stopping; $\hat T\gets$ the number of rounds at which it stops
\State train the ensemble on the whole training slice for $\hat T$ rounds
\State fit the isotonic map $\hat g$ on the PD of the calibration slice (Section~\ref{sec:m4})
\State draw $n_r$ training records; for each constrained input $j$ and each of the $G$ quantile levels, set the input to the level and evaluate the raw score $F$
\State $v_j\gets$ the largest fall of $c_j F$ along the grid, floored at 0; accept the monotonicity when every $v_j\le10^{-12}$ \Comment{Equation~\eqref{eq:cs-pdp}}
\State compute the TreeSHAP contributions of $n_r$ holdout records; $T_i\gets$ the four inputs with the largest contributions
\For{$b=1,\dots,B_r$}
  \State draw $N$ training rows with replacement; train for $\hat T$ rounds; $T_i^{(b)}\gets$ its four largest contributions for the same records
  \State $O_b\gets\frac1{n_r}\sum_i|T_i\cap T_i^{(b)}|/4$ \Comment{Equation~\eqref{eq:cs-overlap}}
\EndFor
\State \Return the ensemble, $\hat g$, $(v_j)_j$, $\bar O$ and $\min_b O_b$
\end{algorithmic}
\end{algorithm}

% ---------------------------------------------------------------------------------------------------------------------
\subsection{TabPFN tabular foundation model}
\label{sec:m6}

TabPFN is a transformer pretrained on a large family of synthetic tabular problems that classifies a new data set by
in-context learning, without training on that data set: the training rows and labels are given to the network as context, and the
prediction for a test row is one forward pass \citep{credit-scoring:hollmann2025}. The reported strength of the model is small
data, up to about ten thousand rows, which is where tree ensembles are also reported to be strongest among general tabular
models \citep{credit-scoring:grinsztajn2022}. With $E$ ensemble members whose predictions the package averages, the raw PD of a
record $x$ given the training sample $\mathcal{D}$ is $\bar q(x\mid\mathcal{D})=\frac1E\sum_{e=1}^{E}q_e(x\mid\mathcal{D})$ and,
after the isotonic map $\hat g$ fitted on the calibration slice,
\begin{equation}
\hat p(x)=\hat g\big(\bar q(x\mid\mathcal{D}_{\mathrm{train}})\big),\qquad E=8.
\label{eq:cs-tabpfn}
\end{equation}
Because no weights are learned from the data, the estimate has no coefficient to inspect and no monotone constraint: the method
is a challenger judged on the same tests as the boosted models, and not a model to deploy where the direction of an effect is a
requirement.

The Brier score of the result is read through the Murphy partition into reliability, resolution and uncertainty, and through the
skill against the trivial forecast \citep{credit-scoring:brier1950,credit-scoring:murphy1973}:
\begin{equation}
\mathrm{BS}=\frac1N\sum_{i=1}^{N}(\hat p_i-y_i)^2=\mathrm{REL}-\mathrm{RES}+\mathrm{UNC},\qquad
\mathrm{UNC}=\bar y(1-\bar y),\qquad \mathrm{BSS}=1-\frac{\mathrm{BS}}{\mathrm{UNC}},
\label{eq:cs-brier}
\end{equation}
with $\bar y$ the default rate of the sample; the partition is exact when the forecast is constant within each group, and the
engine function \texttt{pd\_brier} reports whether it is. The skill is the more informative figure on a sample with a low default
rate, because every Brier score looks small there.

The assumptions of the method are those of the pretrained prior (the new problem resembles the synthetic problems it was trained
on) and the exchangeability of the context with the rows to be scored. The licence is a condition of use. The code is Apache-2.0,
the weights of the default model of the package are licensed for non-commercial use only, and the weights of version~2 carry the
Prior Labs License 1.1, Apache 2.0 with an attribution clause (\emph{Built with PriorLabs-TabPFN}) that allows commercial use
\citep{credit-scoring:tabpfnlicence}. Contraste therefore selects the weights of version~2, shows the licence and the attribution
in every view of the method, caches the weights in the models root of the device, and refuses to fit when the weights did not
come from there. The parameters are $E=8$ ensemble members on the CPU, the training codes of categorical inputs learned on the
training slice (an unseen category is $-1$), and the isotonic map; none is tuned (Algorithm~\ref{alg:cs-tabpfn}). The method runs only on the small case. The
German twin has \csFitGeTrain{} training rows, \csFitGeCal{} calibration rows and a holdout of \csGtN{} applicants with
\csGtEvents{} defaults; the context of a transformer holds all training rows, so cost grows with the training size, as the next paragraph states, and a
context limit is the expected failure on the samples of \csFitTwTrain{} rows and more, which is why the method is restricted to
small cases.
On the vintage samples the method is absent from the environment that produced the committed results (the package is an optional
extra), and no number stands in for it. The implementation is \texttt{pipeline/model/tabpfn\_rung.py} (\texttt{fit\_tabpfn}) on
the \texttt{tabpfn} package, version \csTpVersion{}, with the version~2 weights file \texttt{\csTpWeights{}} on the CPU.

The cost of one forward pass follows from the attention layers of the architecture (as implemented in the package: attention
along a row between the features of one record, and attention along a column between the cells of one feature, in which the
training cells attend to one another and each test cell attends to the training cells). For $N$ training rows, $M$ test rows,
$d$ features and $L$ layers, one member costs $O\bigl(L\,[(N+M)\,d^{2}+(N^{2}+MN)\,d]\bigr)$, and the $E=8$ members multiply it by $E$.
The order is quadratic in the training size and in the number of features, which is the cost behind the restriction to small
cases. The statement is derived from the layers of the package; the original paper was not opened to compare it, and no wall time
of the fit on the German twin is recorded in a committed artifact, so none is quoted.

The acceptance criterion is that the method runs on small cases only, that the weights licence is shown on every view, and that
it meets the checks that apply to it: discrimination, the Brier score and the reliability of the calibrated score (monotonicity
and the stability of reason codes do not apply, because the method has no constraint and no tree attribution). On the German twin the method reaches an AUC of \csGtTpAuc{} against \csGtScAuc{}
for the scorecard, a paired difference of \csGtTpDlDiff{} ($p=\csGtTpDlP{}$): it does not beat the scorecard. Its Brier score is
\csGtTpBrier{} against \csGtScBrier{}, and its expected calibration error \csGtTpEce{} after the isotonic map (\csGtEceTpRaw{} before it) against \csGtScEce{} for the scorecard, the lowest of the
methods on this sample. The sample does not reproduce a gain of the foundation model on small data, and it should not be read as a
ranking of the methods either: with \csGtEvents{} defaults in the holdout the half-width of the interval of each AUC is about
\csGtScHalf{}, and none of the differences in AUC between a learned method and the scorecard is significant.

\begin{algorithm}[htbp]
\caption{Fit of the TabPFN challenger and its isotonic map on a small sample}
\label{alg:cs-tabpfn}
\begin{algorithmic}[1]
\Require training slice $(X,y)$ and calibration slice with at most about ten thousand rows in all; $E=8$ members; weights of version~2 in the models root
\State learn the integer codes of the categorical inputs on the training slice (an unseen category is $-1$)
\State load the weights from the models root; refuse to fit when they did not come from there
\State store the training rows and labels as the context of the network (no weight is updated)
\State $\bar q_i\gets\frac1E\sum_{e=1}^{E}q_e(x_i\mid\mathcal{D}_{\mathrm{train}})$ for the calibration rows, one forward pass per member
\State fit the isotonic map $\hat g$ on $(\bar q_i,y_i)$ of the calibration slice
\State \Return the context, $\hat g$, and the package and weights record
\end{algorithmic}
\end{algorithm}

% ---------------------------------------------------------------------------------------------------------------------
\subsection{Reject inference on a sample with a known rejected population}
\label{sec:m35}

An application scorecard is fitted on the applicants that a legacy policy accepted, so the sample it learns from is selected.
Reject inference tries to repair the selection with the applicants the policy rejected, whose outcomes are never observed. Write
$A_i=1$ when applicant $i$ is accepted. If acceptance depends on the inputs and on noise only, $A\perp Y\mid X$, then
$\Prob(Y=1\mid X,A=1)=\Prob(Y=1\mid X)$ and the accepted-only model is unbiased for the quantity of interest; if the policy also
used information that the model does not see (an underwriter's judgement correlated with the outcome), the missingness is not at
random and the accepted-only model is biased. Whether any method can repair that bias without further information is contested:
\citet{credit-scoring:handhenley1993} argue that reject inference cannot reliably work without untestable assumptions, and
\citet{credit-scoring:crookbanasik2004} ask whether it improves application scoring at all, and both are the references for the
position that a selection driven by a judgement the model does not see cannot be reliably repaired. The augmentation by the
inverse of the acceptance probability is due to \citet{credit-scoring:banasikcrook2007}; the hard cut-off, the parcelling with a
prudence factor and the fuzzy augmentation are standard in the practitioner literature \citep{credit-scoring:siddiqi2017}; and a
comparison of reject-inference methods finds none uniformly best \citep{credit-scoring:ehrhardt2021}. The only design that can test a method is
an accept-all population with known outcomes, as in the framework of \citet{credit-scoring:kozodoi2025}.

Contraste builds that design on the German twin. A legacy policy is simulated on the fully labelled development applicants, the
outcomes of the rejected are hidden, each treatment is fitted on what a lender would have, and each is scored on a held-out
accept-all population. The policy ranks applicants by $u_i=z_i+\kappa\,s_i+\varepsilon_i$ and accepts the lowest share
$\alpha=\csRjAccept{}\,\%$:
\begin{equation}
u_i=z_i+\kappa\,\frac{y_i-\bar y}{\mathrm{sd}(y)}+\varepsilon_i,\qquad \varepsilon_i\sim\mathcal N(0,1),\qquad
A_i=\mathbf 1\big[u_i\ \text{among the lowest }\alpha N_{\mathrm{dev}}\big],
\label{eq:cs-policy}
\end{equation}
where $z_i$ is the standardised logit of a model fitted on every development outcome (the policy's view of the inputs) and
$\kappa\ge0$ is the weight of the private signal; $\kappa=0$ is a policy blind to the outcome given the inputs. The treatments,
all on one L2-penalised logistic learner with $C=\csRjCinv{}$ fixed and never tuned, so that only the treatment of the rejected
differs, are
\begin{align}
\text{augmentation:}\quad & w_i=\frac{1}{\max(\hat\pi_i,\ 0.05)},\quad i\in\text{accepted},\label{eq:cs-aug}\\
\text{parcelling:}\quad & \tilde y_i\sim\mathrm{Bernoulli}\big(\min\{1,\ \rho\,\bar y_{b(i)}^{\,\mathrm{acc}}\}\big),\quad i\in\text{rejected},\label{eq:cs-parcel}\\
\text{fuzzy:}\quad & (y,w)\in\{(1,\hat p_i),\ (0,1-\hat p_i)\},\quad i\in\text{rejected},\label{eq:cs-fuzzy}
\end{align}
where $\hat\pi_i$ is the probability of acceptance estimated on all development applicants, $b(i)$ the score band of applicant $i$
among ten deciles of the accepted-only PD $\hat p_i$, $\bar y_b^{\mathrm{acc}}$ the default rate of the accepted applicants of the band
(the mean $\hat p$ of the rejected of the band when fewer than \csRjMinAcc{} applicants of it were accepted), and $\rho$ the prudence
factor, $\rho=1$ for parcelling and $\rho=\csRjPrudence{}$ for parcelling with prudence. The hard cut-off labels bad the
$\operatorname{round}(\bar{\hat p}\,n_{\mathrm{rej}})$ rejected applicants with the largest accepted-only PD, and the oracle, fitted
on every development outcome, is the bound that no lender can reach. Each treatment is judged by the level bias of its mean PD on
the accept-all test set, by the reduction of the bias against the accepted-only model, and by the change of the AUC:
\begin{equation}
\mathrm{b}_r=\overline{\hat p_r}-\bar y_{\mathrm{test}},\qquad
\mathrm{red}_r=1-\frac{|\bar{\mathrm b}_r|}{|\bar{\mathrm b}_0|},\qquad
\Delta\mathrm{AUC}_r=\mathrm{AUC}_r-\mathrm{AUC}_0,
\label{eq:cs-bias}
\end{equation}
with the bars denoting means over repetitions and the index $0$ the accepted-only model. Each quantity is reported with a 95\,\%
percentile interval over \csRjBoot{} bootstrap resamples of the \csRjReps{} repetitions, paired for the reduction because both
models see the same split, policy and noise.

\begin{algorithm}[htbp]
\caption{One repetition of the reject-inference experiment on a labelled population}
\label{alg:cs-reject}
\begin{algorithmic}[1]
\Require labelled applicants; weight $\kappa$ of the private signal; accepted share $\alpha=0.6$; prudence $\rho=2$
\State split the applicants into development ($70\,\%$) and an accept-all test set ($30\,\%$), stratified by outcome
\State fit the design (scaling and one-hot coding) on development rows only
\State $z\gets$ standardised logit of a model fitted on all development outcomes; draw $\varepsilon$; compute $u$ by Equation~\eqref{eq:cs-policy}; accept the lowest $\alpha$ of $u$
\State hide the outcomes of the rejected
\State fit the oracle on all development outcomes; fit the accepted-only model $M_0$ on the accepted
\State $\hat\pi\gets$ acceptance model on all development rows; augmentation: refit the accepted with weights of Equation~\eqref{eq:cs-aug}
\State $\hat p\gets M_0$ on all development rows; hard cut-off: label bad the $\operatorname{round}(\bar{\hat p}\,n_{\mathrm{rej}})$ riskiest rejected; refit on accepted and rejected
\State parcelling and parcelling with prudence: assign random outcomes to the rejected by Equation~\eqref{eq:cs-parcel}; refit on accepted and rejected
\State fuzzy augmentation: add each rejected twice by Equation~\eqref{eq:cs-fuzzy}; refit with the weights
\State score the test set with every model; \Return $\mathrm b_r$, $\mathrm{AUC}_r$, Brier score of each treatment
\end{algorithmic}
\end{algorithm}

The experiment has \csRjRows{} applicants, \csRjDev{} in development and \csRjTest{} in the test set, \csRjReps{} repetitions per
policy, and four policies with $\kappa\in\{0,\,0.25,\,0.5,\,1\}$. The accepted share, the number of bands, the weight floor of
\csRjFloor{} (no applicant weighs more than 20 times), the prudence factor of \csRjPrudence{} (the rejected assumed twice as bad
as the accepted of their band, after \citealt{credit-scoring:siddiqi2017}) and the minimum of \csRjMinAcc{} accepted applicants per
band are fixed values. The private signal takes four strengths because at $\kappa=1$ alone the accepted book holds
\csRjStrAccRate{}\,\% defaulters against \csRjMarPop{}\,\% in the population, a book no lender holds; the moderate policy, $\kappa=0.5$,
was fixed as the judged regime before it was run, and no treatment, parameter or hypothesis was changed after any result was seen.
The three hypotheses are: under the judged regime the 95\,\% interval of the accepted-only bias lies below zero (H1); at least one
treatment cuts the absolute bias by \csRjHtwo{}\,\% or more with an interval above zero (H2); and no treatment improves the AUC by
more than \csRjHthree{} over the accepted-only model (H3). Each repetition fits nine logistic models: one for the policy, one for
the oracle, one for the accepted-only model, one for the acceptance probability, and one each for augmentation, the hard cut-off,
the two parcellings and the fuzzy augmentation. The experiment therefore fits $4\times100\times9=3{,}600$ models, and each of the
28 combinations of a policy and a treatment adds \csRjBoot{} bootstrap resamples of 100 values per interval. The implementation is
\texttt{pipeline/cases/c01\_reject.py} (\texttt{legacy\_acceptance}, \texttt{repetition}, \texttt{parcel\_rates},
\texttt{augmentation\_weights}, \texttt{interval}, \texttt{reduction}), in numpy, pandas and \texttt{scikit-learn} 1.9.1.

The acceptance criterion is that the bias of the accepted-only model on the full population is measured and that each treatment's
reduction of it is reported with its interval. Table~\ref{tab:cs-reject} and Table~\ref{tab:cs-reject-regimes} report both, and
Section~\ref{sec:results-credit-scoring} reads them: H1 \csRjHoneHolds{}, H2 \csRjHtwoHolds{} (for \csRjHtwoRung{} only) and
H3 \csRjHthreeHolds{}.

% ---------------------------------------------------------------------------------------------------------------------
\subsection{Results}
\label{sec:results-credit-scoring}

\paragraph{Retail cards, the observed holdout.}
Table~\ref{tab:cs-c01-holdout} compares the methods on the locked holdout of the Taiwan sample, \csHoN{} cards with \csHoEvents{}
defaults (a default rate of \csHoRate{}\,\%). The scorecard reaches an AUC of \csHoScAuc{} (95\,\% interval \csHoScAucLo{} to
\csHoScAucHi{}); the single-variable benchmark on the September repayment status already reaches \csHoPzAuc{}, so the scorecard
lifts the AUC by \csHoScOverPz{} over it. The four learned methods sit between \csHoEbAuc{} and \csHoLgAuc{}, a gain of
\csHoEbDlDiff{} to \csHoLgDlDiff{} over the scorecard, and every paired DeLong $p$-value is below $10^{-9}$, so the gain is
statistically unambiguous at \csHoN{} cards. The table does not show whether a gain of this size changes a decision; the decision
impact at a cut-off is a separate output of the case. The penalised logistic regression repeats the scorecard exactly, which
shows that the estimator does not produce the gain; whether the gain comes from non-linear structure or from the
\csFitTwDropIv{} inputs that the scorecard screens out (the learned methods use all \csFitTwInputs{} raw inputs, the scorecard the
\csFitTwKept{} binned ones) is not isolated by the artifacts.
Figure~\ref{fig:cs-roc} shows where the gain sits, in the right panel as the difference in true positive rate; it does not show
the uncertainty of the curves, which the intervals of the table give.

\begin{figure}[htbp]
\centering
\includegraphics[width=\textwidth]{credit-scoring-roc}
\caption{ROC curves of the single-variable benchmark, the scorecard, the penalised logistic tree regression, the explainable
boosting machine and LightGBM on the locked holdout of the Taiwan card-default sample (left), and the difference in true positive
rate between each method and the scorecard along the false positive rate (right).}
\label{fig:cs-roc}
\end{figure}

\begin{table}[htbp]
\centering
\footnotesize
\setlength{\tabcolsep}{3.5pt}
\caption{Discrimination and calibration of each method on the locked holdout of the Taiwan card-default sample (9,000 cards, 1,991 defaults): AUC with the DeLong interval, Kolmogorov-Smirnov separation, Brier score, expected calibration error, the paired DeLong difference against the scorecard and the Hosmer-Lemeshow $p$-value.}
\label{tab:cs-c01-holdout}
\begin{tabular}{lrrrrrr}
\toprule
Method & AUC (95\% CI) & KS & Brier & ECE & \shortstack{$\Delta$AUC vs scorecard\\(DeLong $p$)} & \shortstack{H-L\\$p$} \\
\midrule
\relax Single variable & \relax $0.692\pm0.013$ & \relax 0.360 & \relax 0.1431 & \relax 0.0144 & \relax $-$0.0670 (\ensuremath{1.3\times10^{-40}}) & \relax 0.021 \\
\relax WoE scorecard & \relax $0.759\pm0.013$ & \relax 0.397 & \relax 0.1393 & \relax 0.0155 & \relax -- & \relax 0.015 \\
\relax L1 logistic & \relax $0.759\pm0.013$ & \relax 0.397 & \relax 0.1393 & \relax 0.0153 & \relax $-3.0\times10^{-5}$ (0.088) & \relax 0.014 \\
\relax PLTR & \relax $0.775\pm0.012$ & \relax 0.422 & \relax 0.1367 & \relax 0.0113 & \relax +0.0162 (\ensuremath{2.6\times10^{-11}}) & \relax 0.184 \\
\relax EBM & \relax $0.775\pm0.012$ & \relax 0.426 & \relax 0.1367 & \relax 0.0118 & \relax +0.0159 (\ensuremath{2.7\times10^{-10}}) & \relax 0.079 \\
\relax LightGBM & \relax $0.777\pm0.012$ & \relax 0.420 & \relax 0.1361 & \relax 0.0134 & \relax +0.0183 (\ensuremath{6.5\times10^{-13}}) & \relax 0.008 \\
\relax XGBoost & \relax $0.775\pm0.012$ & \relax 0.422 & \relax 0.1359 & \relax 0.0124 & \relax +0.0168 (\ensuremath{2.2\times10^{-12}}) & \relax 0.125 \\
\bottomrule
\end{tabular}
\end{table}


\paragraph{Calibration.}
The scorecard has a Brier score of \csHoScBrier{}, a Brier skill of \csHoScSkill{} against the trivial forecast, and a
Hosmer--Lemeshow $p$-value of \csHoScHlP{}; the learned methods lower the Brier score to between \csHoXgBrier{} and
\csHoPlBrier{}. The Hosmer--Lemeshow statistic \citep{credit-scoring:hosmer1980} flags the scorecard and LightGBM (\csHoLgHlP{}),
and does not flag the penalised logistic tree regression (\csHoPlHlP{}), so better ranking does not imply better calibration
across the methods, and the calibration map of Section~\ref{sec:m5} does not guarantee a passing test. Figure~\ref{fig:cs-reliability}
shows the reliability of the scorecard, of the explainable boosting machine and of LightGBM in ten quantile bins: the isotonic map
brings the explainable boosting machine to the diagonal (expected calibration error \csHoEceEbRaw{} to \csHoEceEbCal{}) but leaves the
LightGBM curve where it was (\csHoEceLgRaw{} to \csHoEceLgCal{}). The figure does not show the sampling error of each bin, a tenth
of the \csHoN{} holdout cards, so differences between the curves of the order of the bin error carry little information, and the
test statistics of the table are the measure to rely on. The Spiegelhalter test, which works on the obligors without binning
\citep{credit-scoring:spiegelhalter1986}, agrees with the Hosmer--Lemeshow test on the scorecard ($p=\csHoScSpP{}$). The expected
calibration error used above is the mean absolute gap between the predicted and the observed default rate over the ten quantile
bins, weighted by bin size \citep{credit-scoring:naeini2015}.

\begin{figure}[htbp]
\centering
\includegraphics[width=\textwidth]{credit-scoring-reliability}
\caption{Reliability of the scorecard (left), the explainable boosting machine (centre) and LightGBM (right) on the Taiwan
holdout: the observed default rate against the mean predicted PD in ten quantile bins, before and after the isotonic calibration
map where one is fitted, with the diagonal of perfect calibration.}
\label{fig:cs-reliability}
\end{figure}

\paragraph{Shifted populations and the small sample.}
Table~\ref{tab:cs-c01-variants} and Figure~\ref{fig:cs-variants} give the AUC on the seven populations of the case. The order of
the methods is the same on the observed holdout, on the two covariate shifts, on the prior shift and on the label noise: the
scorecard and the penalised logistic regression are lowest, tied, and LightGBM is highest. The AUC of every method rises under
the covariate shifts (the severe shift gives the scorecard \csDsScAuc{} against \csHoScAuc{}), which is consistent with a
resampling toward delinquent records that are easier to separate, an explanation the artifacts do not test, and label noise lowers
every method (the scorecard falls to \csLnScAuc{}). The table shows ranking only, and the prior shift moves calibration far more than ranking: it raises the expected
calibration error of the scorecard to \csPsScEce{}. That shift resamples rows so that the distribution of the inputs given the
outcome is unchanged and only the default rate moves, so the log-odds shift by a constant that the correction of the intercept for
a known change of the event fraction removes \citep{credit-scoring:king2001}; ranking is not the property at risk. The population
stability index of the scorecard score, the symmetrised divergence between the binned score distributions of the training slice
and the sample, is \csHoScPsi{} on the holdout, below the 5\,\% critical value of \csHoScPsiBench{} of the chi-square benchmark for samples of
these sizes \citep{credit-scoring:yurdakul2020}, and \csDsScPsi{} under the severe covariate shift, far above it. On the 2,000-row refit the ordering
changes: the explainable boosting machine leads (\csSsEbAuc{}), LightGBM follows (\csSsLgAuc{}) and the penalised logistic tree
regression (\csSsPlAuc{}) is no better than the scorecard (\csSsScAuc{}), which is consistent with the low stability of its rule
set at that size. On the German twin no learned method differs significantly from the scorecard.

\begin{table}[htbp]
\centering
\small
\caption{AUC of the methods on each evaluation population of the Taiwan sample and on the German twin: the observed holdout, two covariate shifts, a prior shift, label noise, the methods refitted on 2,000 rows, and the German twin, where TabPFN runs.}
\label{tab:cs-c01-variants}
\begin{tabular}{lrrrrrr}
\toprule
Variant & Scorecard & L1 & PLTR & EBM & LightGBM & TabPFN \\
\midrule
\relax Holdout & \relax 0.7587 & \relax 0.7587 & \relax 0.7749 & \relax 0.7746 & \relax 0.7770 & \relax -- \\
\relax Drift 0.25 & \relax 0.7842 & \relax 0.7842 & \relax 0.7966 & \relax 0.7953 & \relax 0.7994 & \relax -- \\
\relax Drift 0.6 & \relax 0.8178 & \relax 0.8178 & \relax 0.8244 & \relax 0.8230 & \relax 0.8277 & \relax -- \\
\relax Prior shift x1.5 & \relax 0.7689 & \relax 0.7689 & \relax 0.7811 & \relax 0.7838 & \relax 0.7852 & \relax -- \\
\relax Label noise 5\% & \relax 0.7174 & \relax 0.7174 & \relax 0.7321 & \relax 0.7314 & \relax 0.7340 & \relax -- \\
\relax 2,000 rows & \relax 0.7525 & \relax 0.7511 & \relax 0.7528 & \relax 0.7649 & \relax 0.7601 & \relax -- \\
\relax German twin & \relax 0.7883 & \relax 0.7908 & \relax 0.7861 & \relax 0.7874 & \relax 0.7784 & \relax 0.7769 \\
\bottomrule
\end{tabular}
\end{table}


\begin{figure}[htbp]
\centering
\includegraphics[width=\textwidth]{credit-scoring-variants}
\caption{AUC with its 95\,\% DeLong interval of the scorecard, the penalised logistic tree regression, the explainable boosting
machine and LightGBM on the observed Taiwan holdout, the two covariate shifts, the prior shift, the label noise and the
2,000-row refit.}
\label{fig:cs-variants}
\end{figure}

\paragraph{Structural checks.}
Table~\ref{tab:cs-checks} and Figure~\ref{fig:cs-stability} report the checks that the acceptance criteria name. Every scorecard
coefficient is negative on the three fits, no monotone violation is found for either boosted model on the partial-dependence
grid, the top-four reasons overlap by \csFitTwReasonMean{} on the Taiwan fit and \csFitGeReasonMean{} on the German fit, and the
rule set of the penalised logistic tree regression has a mean Jaccard index of \csFitTwPlJm{} (Taiwan, \csFitTwTrain{} rows),
\csFitSmPlJm{} (2,000 rows) and \csFitGePlJm{} (German, \csFitGeTrain{} rows). Rule membership is therefore unstable even where the
predictions are not: the rules selected on a resample and on the full slice share only about \csFitTwPlJm{} of their union on the
large sample, and less on the small ones. The checks do not say why; correlated rules that substitute for each other would produce the pattern,
and the artifacts do not test that explanation. What follows for use is that the rules of the method should be read as examples
of the structure the data carry, not as a stable list of risk drivers.

\begin{table}[htbp]
\centering
\scriptsize
\setlength{\tabcolsep}{3.5pt}
\caption{Structural checks of the fitted methods on the three fits of C01: the sign of every scorecard coefficient, the stability of the PLTR rule set over 20 bootstrap resamples, the largest violation of the declared monotone direction on the partial-dependence grid, and the overlap of the four principal reasons.}
\label{tab:cs-checks}
\begin{tabular}{lrrrrr}
\toprule
Fit & \shortstack{Negative WoE\\coefficients} & \shortstack{PLTR Jaccard\\mean / min} & \shortstack{Largest monotone\\violation} & \shortstack{Constrained\\inputs} & \shortstack{Top-4 reason overlap\\mean / min} \\
\midrule
\relax Taiwan, 16,800 rows & \relax 13 of 13 & \relax 0.411 / 0.367 & \relax 0 & \relax 13 & \relax 0.760 / 0.725 \\
\relax Taiwan, 2,000 rows & \relax 10 of 10 & \relax 0.349 / 0.281 & \relax 0 & \relax 13 & \relax 0.719 / 0.662 \\
\relax German, 560 rows & \relax 10 of 10 & \relax 0.280 / 0.196 & \relax 0 & \relax 2 & \relax 0.731 / 0.624 \\
\bottomrule
\end{tabular}
\end{table}


\begin{figure}[htbp]
\centering
\includegraphics[width=\textwidth]{credit-scoring-stability}
\caption{Stability of the learned methods over bootstrap resamples on the three fits of the retail case: the Jaccard index of the
rule set of the penalised logistic tree regression over 20 resamples (left) and the mean overlap of the four principal reasons of
LightGBM over 10 refits (right).}
\label{fig:cs-stability}
\end{figure}

\paragraph{Corporate bankruptcy across horizons.}
Table~\ref{tab:cs-c02-auc}, Table~\ref{tab:cs-c02-diff} and Figure~\ref{fig:cs-c02} give the C02 results. The scorecard AUC
falls from \csCtwoAScAuc{} at one year to \csCtwoEScAuc{} at five years, not monotonically (\csCtwoBScAuc{} at two years and
\csCtwoCScAuc{} at three), whereas LightGBM stays between \csCtwoDLgAuc{} and \csCtwoALgAuc{}. Every learned method except the
penalised logistic regression beats the scorecard at every horizon (the paired differences of LightGBM run from \csCtwoALgDiff{} at
one year to \csCtwoELgDiff{} at five years, each with a DeLong $p$ below 0.01), and the penalised logistic regression does so at
four of five horizons and on the pooled file. The holdouts are small: \csCtwoHoldEvE{} bankruptcies at five years and
\csCtwoHoldEvA{} at one year, so each AUC interval is wide, with a half-width of \csCtwoEScHalf{} for the scorecard at five years;
in the right panel of Figure~\ref{fig:cs-c02} the interval of the paired difference excludes zero for every learned method at every
horizon except the penalised logistic regression at one year. The tables do not show out-of-time behaviour: the five files are
independent samples of different forecast years, the file-to-horizon mapping follows the data page and the source paper was not
checked, and the event rates fall with the horizon (from \csCtwoRateA{}\,\% to \csCtwoRateE{}\,\%) so no cumulative PD is read from
them. The overfit flag of the engine, which is raised when the train-minus-holdout AUC gap exceeds 0.05, is raised for
\csCtwoFlagTotal{} of the \csCtwoRungTotal{} fitted methods across the five horizons, and at four and five years it reaches even
the scorecard; on the pooled file, with \csCtwoHoldEvP{} holdout bankruptcies, none is flagged. What follows is that on small
corporate samples a learned method can gain much in AUC and still be the less trustworthy estimate, so the gain must be read
with the gap.

\begin{table}[htbp]
\centering
\small
\caption{AUC of each method on the locked holdout of the Polish company bankruptcy files, by forecast horizon and for the five files pooled with a horizon indicator, with the number of bankruptcies in each holdout.}
\label{tab:cs-c02-auc}
\begin{tabular}{lrrrrrrr}
\toprule
Horizon & Holdout defaults & Scorecard & L1 & PLTR & EBM & LightGBM & XGBoost \\
\midrule
\relax 1 year & \relax 123 & \relax 0.9106 & \relax 0.9115 & \relax 0.9343 & \relax 0.9538 & \relax 0.9520 & \relax 0.9412 \\
\relax 2 years & \relax 155 & \relax 0.8100 & \relax 0.8348 & \relax 0.9036 & \relax 0.9217 & \relax 0.9228 & \relax 0.9124 \\
\relax 3 years & \relax 149 & \relax 0.8388 & \relax 0.8670 & \relax 0.8895 & \relax 0.8983 & \relax 0.9253 & \relax 0.9009 \\
\relax 4 years & \relax 120 & \relax 0.7920 & \relax 0.8123 & \relax 0.8973 & \relax 0.9011 & \relax 0.9117 & \relax 0.8887 \\
\relax 5 years & \relax 81 & \relax 0.7279 & \relax 0.7583 & \relax 0.9029 & \relax 0.9320 & \relax 0.9428 & \relax 0.9269 \\
\relax Pooled & \relax 628 & \relax 0.8648 & \relax 0.8824 & \relax 0.9168 & \relax 0.9450 & \relax 0.9703 & \relax 0.9642 \\
\bottomrule
\end{tabular}
\end{table}


\begin{table}[htbp]
\centering
\small
\caption{Paired AUC difference of each learned method against the scorecard on the Polish bankruptcy holdouts (positive favours the learned method; a star marks a DeLong $p$ below 0.05) and the number of the seven fitted methods whose train-minus-holdout AUC gap exceeds the 0.05 tolerance.}
\label{tab:cs-c02-diff}
\begin{tabular}{lrrrrrr}
\toprule
Horizon & L1 & PLTR & EBM & LightGBM & XGBoost & Overfit flags \\
\midrule
\relax 1 year & \relax +0.001 & \relax +0.024$^{*}$ & \relax +0.043$^{*}$ & \relax +0.041$^{*}$ & \relax +0.031$^{*}$ & \relax 1 of 7 \\
\relax 2 years & \relax +0.025$^{*}$ & \relax +0.094$^{*}$ & \relax +0.112$^{*}$ & \relax +0.113$^{*}$ & \relax +0.102$^{*}$ & \relax 4 of 7 \\
\relax 3 years & \relax +0.028$^{*}$ & \relax +0.051$^{*}$ & \relax +0.059$^{*}$ & \relax +0.086$^{*}$ & \relax +0.062$^{*}$ & \relax 4 of 7 \\
\relax 4 years & \relax +0.020$^{*}$ & \relax +0.105$^{*}$ & \relax +0.109$^{*}$ & \relax +0.120$^{*}$ & \relax +0.097$^{*}$ & \relax 6 of 7 \\
\relax 5 years & \relax +0.030$^{*}$ & \relax +0.175$^{*}$ & \relax +0.204$^{*}$ & \relax +0.215$^{*}$ & \relax +0.199$^{*}$ & \relax 6 of 7 \\
\relax Pooled & \relax +0.018$^{*}$ & \relax +0.052$^{*}$ & \relax +0.080$^{*}$ & \relax +0.106$^{*}$ & \relax +0.099$^{*}$ & \relax 0 of 7 \\
\bottomrule
\end{tabular}
\end{table}


\begin{figure}[htbp]
\centering
\includegraphics[width=\textwidth]{credit-scoring-c02}
\caption{Holdout AUC of the scorecard and the learned methods on the Polish bankruptcy files, by forecast horizon, with the 95\,\%
DeLong interval (left), and the paired AUC difference of each learned method against the scorecard with its DeLong interval
(right).}
\label{fig:cs-c02}
\end{figure}

\paragraph{Small-business loans by vintage.}
Table~\ref{tab:cs-c03-auc}, Table~\ref{tab:cs-c03-vintage}, Figure~\ref{fig:cs-c03} and Figure~\ref{fig:cs-vintage} give the C03
results. Discrimination holds through the cycle: the AUC of the scorecard is between \csCthreeDScAuc{} and \csCthreeEScAuc{} on
the five cohorts, and no learned method dominates, since LightGBM leads on two cohorts, the penalised logistic regression on two
and the penalised logistic tree regression on one, while the latter and the explainable boosting machine fall to
\csCthreeEPlAuc{} and \csCthreeEEbAuc{} on the last cohort, which may reflect fitting on five origination attributes and
a stratified sample of \csCthreeFitRows{} loans; the artifacts do not test that explanation. Level is another matter. The scorecard for the FY2006--07 cohort, fitted on
approvals through FY\csCthreeALast{}, predicts a 5-year charge-off probability of \csCthreeVAPred{}\,\% for the FY2006 vintage that went on to
\csCthreeVAObs{}\,\%, an under-prediction of \csCthreeVAGap{} points, and \csCthreeVBPred{}\,\% against \csCthreeVBObs{}\,\% for FY2007
(\csCthreeVBGap{} points); the gap is \csCthreeVCGap{} points for FY2008, \csCthreeVDGap{} for FY2009, and the sign reverses from FY2010
(\csCthreeVEGap{} points), with the models trained through the crisis over-predicting by \csCthreeLaterMin{} to
\csCthreeLaterMax{} points in every later vintage (\csCthreeVFGap{} for FY2011 and \csCthreeVJGap{} for FY2015).
Figure~\ref{fig:cs-vintage} shows that the shortfall of the crisis vintages is already present at 12 months of life. The tables
do not show whether a model
fitted on a bank's own book would have behaved so: the loans are guaranteed small-business loans made under a federal program,
the file carries no delinquency history, balance or rating, and the sample is not the SME book of a bank.

\begin{table}[htbp]
\centering
\footnotesize
\setlength{\tabcolsep}{3.5pt}
\caption{AUC with its 95\% interval of each method on the five approval cohorts of the SBA 7(a) loans, each scored against the 5-year charge-off outcome with a model fitted only on approval years whose 5-year mark closed before the cohort began.}
\label{tab:cs-c03-auc}
\begin{tabular}{lrrrrrr}
\toprule
Test cohort & \shortstack{Training\\charge-offs} & Scorecard & L1 & PLTR & EBM & LightGBM \\
\midrule
\relax FY2006--07 & \relax 23,320 & \relax $0.825\pm0.002$ & \relax $0.826\pm0.002$ & \relax $0.770\pm0.003$ & \relax $0.734\pm0.003$ & \relax $0.797\pm0.002$ \\
\relax FY2008--10 & \relax 28,393 & \relax $0.815\pm0.003$ & \relax $0.814\pm0.003$ & \relax $0.775\pm0.004$ & \relax $0.777\pm0.003$ & \relax $0.832\pm0.003$ \\
\relax FY2011--14 & \relax 50,614 & \relax $0.817\pm0.006$ & \relax $0.835\pm0.005$ & \relax $0.814\pm0.006$ & \relax $0.828\pm0.005$ & \relax $0.873\pm0.004$ \\
\relax FY2015--17 & \relax 110,349 & \relax $0.793\pm0.005$ & \relax $0.806\pm0.005$ & \relax $0.849\pm0.004$ & \relax $0.795\pm0.005$ & \relax $0.837\pm0.004$ \\
\relax FY2018--20 & \relax 115,237 & \relax $0.826\pm0.004$ & \relax $0.834\pm0.004$ & \relax $0.739\pm0.007$ & \relax $0.746\pm0.006$ & \relax $0.830\pm0.004$ \\
\bottomrule
\end{tabular}
\end{table}


\begin{table}[htbp]
\centering
\small
\caption{Calibration of the scorecard by approval fiscal year on the SBA 7(a) loans: the observed 5-year charge-off rate, the mean predicted 5-year PD of the scorecard fitted for that year's cohort, the gap in percentage points and the AUC within the year.}
\label{tab:cs-c03-vintage}
\begin{tabular}{lrrrrr}
\toprule
Approval year & Loans & Observed 5-year PD (\%) & Scorecard PD (\%) & Gap (points) & AUC \\
\midrule
\relax 2006 & \relax 86,393 & \relax 22.31 & \relax 8.80 & \relax +13.51 & \relax 0.8332 \\
\relax 2007 & \relax 88,134 & \relax 27.58 & \relax 8.76 & \relax +18.82 & \relax 0.8204 \\
\relax 2008 & \relax 61,539 & \relax 21.27 & \relax 7.35 & \relax +13.91 & \relax 0.8132 \\
\relax 2009 & \relax 36,594 & \relax 8.38 & \relax 6.15 & \relax +2.23 & \relax 0.7845 \\
\relax 2010 & \relax 39,890 & \relax 4.90 & \relax 5.41 & \relax $-$0.51 & \relax 0.8102 \\
\relax 2011 & \relax 45,592 & \relax 3.65 & \relax 6.49 & \relax $-$2.84 & \relax 0.7983 \\
\relax 2012 & \relax 38,861 & \relax 3.27 & \relax 6.55 & \relax $-$3.27 & \relax 0.8261 \\
\relax 2013 & \relax 40,400 & \relax 3.07 & \relax 6.20 & \relax $-$3.13 & \relax 0.8236 \\
\relax 2014 & \relax 45,942 & \relax 3.52 & \relax 6.36 & \relax $-$2.84 & \relax 0.8235 \\
\relax 2015 & \relax 55,387 & \relax 3.85 & \relax 8.82 & \relax $-$4.97 & \relax 0.8085 \\
\relax 2016 & \relax 56,759 & \relax 3.87 & \relax 8.36 & \relax $-$4.49 & \relax 0.7932 \\
\relax 2017 & \relax 56,057 & \relax 4.31 & \relax 8.04 & \relax $-$3.73 & \relax 0.7841 \\
\relax 2018 & \relax 54,171 & \relax 4.66 & \relax 6.68 & \relax $-$2.01 & \relax 0.8104 \\
\relax 2019 & \relax 45,664 & \relax 3.91 & \relax 5.76 & \relax $-$1.85 & \relax 0.8321 \\
\relax 2020 & \relax 36,456 & \relax 2.59 & \relax 5.55 & \relax $-$2.96 & \relax 0.8418 \\
\bottomrule
\end{tabular}
\end{table}


\begin{figure}[htbp]
\centering
\includegraphics[width=\textwidth]{credit-scoring-c03}
\caption{AUC with its 95\,\% interval of the scorecard and the learned methods on the five approval cohorts of the SBA 7(a) loans
(left), and the observed 5-year charge-off rate against the mean 5-year PD predicted by the scorecard of the cohort, by approval
fiscal year (right).}
\label{fig:cs-c03}
\end{figure}

\begin{figure}[htbp]
\centering
\includegraphics[width=0.78\textwidth]{credit-scoring-vintage}
\caption{Cumulative charge-off probability by months since approval for the SBA 7(a) vintages FY2007, FY2008, FY2010 and FY2013:
observed (solid lines with filled markers) and predicted by the scorecard of the vintage's cohort, with the 5-year PD allocated in
time by the share of charge-offs of the training window (dotted lines with open markers).}
\label{fig:cs-vintage}
\end{figure}

\paragraph{Reject inference.}
Table~\ref{tab:cs-reject} and Table~\ref{tab:cs-reject-regimes} and Figure~\ref{fig:cs-reject} give the result of the experiment.
Under the policy with no private signal ($\kappa=0$) the accepted-only model is nearly unbiased (\csRjMarAoBias{} points), as the
derivation predicts, and augmentation recovers \csRjMarAuRed{}\,\% of the small bias. As the private signal grows, the accepted
book thins from \csRjMarAccRate{}\,\% defaulters to \csRjStrAccRate{}\,\% against \csRjMarPop{}\,\% in the population, and the bias of
the accepted-only model grows to \csRjWeakAoBias{}, \csRjModAoBias{} and \csRjStrAoBias{} points. Under the judged regime
($\kappa=0.5$) four of the five treatments cut the bias by at most \csRjModPaRed{}\,\% and rank the full population worse than the
accepted-only model (changes in AUC of \csRjModAuDAuc{} for augmentation, \csRjModHcDAuc{} for the hard cut-off and \csRjModPaDAuc{} for
parcelling); only parcelling with the prudence factor reduces it, by \csRjModPpRed{}\,\% (interval \csRjModPpRedLo{} to
\csRjModPpRedHi{}), with an AUC gain of \csRjModPpDAuc{}. H1 holds, with a bias of \csRjModAoBias{} points (interval \csRjModAoBiasLo{}
to \csRjModAoBiasHi{}); H2 holds for that one treatment; H3 does not hold, because the gain of \csRjModPpDAuc{} exceeds
\csRjHthree{}. The same treatment over-corrects when the policy is blind to the outcome, to \csRjMarPpBias{} points under
$\kappa=0$: a prudence factor is a bet on the unseen outcomes of the rejected, and it pays only when the rejected are worse than
their score band says. The tables do not show a result on a real rejected population, whose outcomes are unknown; the
experiment measures what a treatment does when the truth is known by construction on \csRjRows{} applicants, and the bias and the AUC
differences are means over repetitions of one policy design. What follows for validation is that the accepted-only bias must be
measured, not assumed, and that no standard treatment should be adopted on its AUC alone.

\begin{table}[htbp]
\centering
\scriptsize
\setlength{\tabcolsep}{3.5pt}
\caption{Reject inference on the German twin under the moderate private-signal policy: level bias of the mean PD against the default rate of the accept-all population, its reduction against the accepted-only model, the AUC change and the Brier score, each a mean over 100 repetitions with a 95\% bootstrap interval.}
\label{tab:cs-reject}
\begin{tabular}{lrrrrr}
\toprule
Treatment & \shortstack{Level bias\\(points)} & \shortstack{Bias reduction\\(\%)} & \shortstack{$\Delta$AUC vs\\accepted only} & AUC & Brier \\
\midrule
\relax Oracle & \relax $0.09$ [$-0.21$, $0.39$] & \relax $99.5$ [$97.7$, $100.0$] & \relax $0.0971$ [$0.0870$, $0.1065$] & \relax 0.7827 & \relax 0.1676 \\
\relax Accepted only & \relax $-17.19$ [$-17.67$, $-16.70$] & \relax -- & \relax -- & \relax 0.6856 & \relax 0.2237 \\
\relax Augmentation & \relax $-16.93$ [$-17.55$, $-16.35$] & \relax $1.5$ [$0.1$, $2.8$] & \relax $-0.0365$ [$-0.0421$, $-0.0313$] & \relax 0.6491 & \relax 0.2295 \\
\relax Hard cut-off & \relax $-17.08$ [$-17.60$, $-16.54$] & \relax $0.6$ [$-0.1$, $1.3$] & \relax $-0.0313$ [$-0.0342$, $-0.0285$] & \relax 0.6544 & \relax 0.2366 \\
\relax Parcelling & \relax $-16.79$ [$-17.35$, $-16.21$] & \relax $2.3$ [$1.0$, $3.5$] & \relax $-0.0161$ [$-0.0214$, $-0.0113$] & \relax 0.6695 & \relax 0.2256 \\
\relax Parcelling, prudence 2 & \relax $-9.57$ [$-10.41$, $-8.77$] & \relax $44.3$ [$41.2$, $47.6$] & \relax $0.0208$ [$0.0171$, $0.0243$] & \relax 0.7064 & \relax 0.2079 \\
\relax Fuzzy augmentation & \relax $-17.18$ [$-17.68$, $-16.70$] & \relax $0.0$ [$-0.0$, $0.0$] & \relax $-0.0002$ [$-0.0002$, $-0.0001$] & \relax 0.6855 & \relax 0.2237 \\
\bottomrule
\end{tabular}
\end{table}


\begin{table}[htbp]
\centering
\scriptsize
\setlength{\tabcolsep}{3.5pt}
\caption{Reject inference across the four legacy policies of the German twin, ordered by the weight $\kappa$ of the underwriter's private signal: the accepted book, the accepted-only bias and AUC, the treatment with the largest bias reduction and the bias of the prudence-2 parcelling.}
\label{tab:cs-reject-regimes}
\begin{tabular}{rrrrlrr}
\toprule
$\kappa$ & \shortstack{Defaulters among\\accepted (\%)} & \shortstack{Accepted-only\\bias (points)} & \shortstack{Accepted-only\\AUC} & \shortstack{Best\\treatment} & \shortstack{Reduction (\%)\\(95\% CI)} & \shortstack{Prudence-2\\bias (points)} \\
\midrule
\relax 0 & \relax 18.2 & \relax $-$1.32 & \relax 0.7613 & \relax Augmentation & \relax 92.7 [47.1, 99.1] & \relax +10.07 \\
\relax 0.25 & \relax 13.3 & \relax $-$9.31 & \relax 0.7398 & \relax Parcelling, prudence 2 & \relax 86.7 [77.5, 95.1] & \relax +1.24 \\
\relax 0.5 & \relax 8.9 & \relax $-$17.19 & \relax 0.6856 & \relax Parcelling, prudence 2 & \relax 44.3 [41.2, 47.6] & \relax $-$9.57 \\
\relax 1 & \relax 2.9 & \relax $-$26.69 & \relax 0.5388 & \relax Parcelling, prudence 2 & \relax 6.4 [5.6, 7.2] & \relax $-$24.99 \\
\bottomrule
\end{tabular}
\end{table}


\begin{figure}[htbp]
\centering
\includegraphics[width=\textwidth]{credit-scoring-reject}
\caption{Level bias of the mean PD on the accept-all German test set with its 95\,\% bootstrap interval (left) and AUC (right) of
the oracle, the accepted-only model, augmentation, parcelling and parcelling with prudence factor 2, against the weight of the
underwriter's private signal in the legacy policy.}
\label{fig:cs-reject}
\end{figure}

\paragraph{What the results support.}
Across the three cases the scorecard is a strong reference. It matches the penalised regression exactly on the Taiwan cards,
where the learned methods lead it by about two AUC points; the penalised regression matches or exceeds it on the Polish files
and on the SBA cohorts; on the SBA cohorts no learned method dominates, each leading the scorecard on some cohorts and trailing
it on others (Table~\ref{tab:cs-c03-auc}); and on the small Polish corporate samples the learned methods lead it by more
(Table~\ref{tab:cs-c02-diff}). It is the only method whose every parameter can be read from a table. The learned methods earn their place where the paired test
and the structural checks both agree, as for the explainable boosting machine and LightGBM on the Taiwan holdout, and they do
not where the sample is small, as on the German twin. Calibration, stability of the selected structure and the level of the
estimate through a cycle are separate properties from ranking, and each needed its own measurement.
